\chapter{Detailed computations}
\label{chap:detailed-computations}
% Regimes I--III : tables span all families (Convention~\ref{conv:regime-tags}).

Explicit basis elements, differential matrices with all signs, and
residue computations at each collision divisor. The chapter separates
chain groups, bar cohomology, the dual coalgebra $A^i=H^\star B(A)$,
the Verdier--Koszul dual algebra $A^!$, and the derived chiral centre
$Z^{\mathrm{der}}_{\mathrm{ch}}(A)$ at every table where those objects
can otherwise be confused.

\noindent\emph{Degree-by-degree bar complex computations for $\mathfrak{sl}_3$, $\mathcal{W}_3$, $E_8$, and the Virasoro algebra, complementing the Heisenberg calculations of Chapter~\ref{ch:heisenberg-frame}. The tables record structure constants, differential matrices, residue terms, and the exact object computed.}

\section{Heisenberg computations}

\subsection{Bar complex: degrees 3 and 4}

\begin{computation}[Heisenberg bar complex: full degree 3]\label{comp:heisenberg-deg3-full}
The degree 3 component of $\B(\cH)$ is:
\[
\B_3(\cH) = \bigoplus_{m,n,p > 0} \bC \cdot [a_{-m}|a_{-n}|a_{-p}] \otimes \Omega^2(\overline{\Conf}_3)
\]

\emph{Basis for $\Omega^2(\overline{\Conf}_3)$.}
The three 2-forms $\eta_{12} \wedge \eta_{13}$, $\eta_{12} \wedge \eta_{23}$, and $\eta_{13} \wedge \eta_{23}$ (where $\eta_{ij} = d\log(z_i - z_j)$) satisfy the Arnold relation:
\[
\eta_{12} \wedge \eta_{23} + \eta_{23} \wedge \eta_{31} + \eta_{31} \wedge \eta_{12} = 0
\]
(with $\eta_{31} = \eta_{13}$, since $d\log(z_3 - z_1) = \frac{dz_3 - dz_1}{z_3 - z_1} = \frac{-(dz_1 - dz_3)}{-(z_1 - z_3)} = d\log(z_1 - z_3)$). Therefore $\dim \Omega^2(\overline{\Conf}_3) = 2$, consistent with the Poincar\'{e} polynomial $P_t(C_3(\mathbb{C})) = 1 + 3t + 2t^2$. A basis is $\{\eta_{12} \wedge \eta_{13},\; \eta_{13} \wedge \eta_{23}\}$, and the Arnold relation gives $\eta_{12} \wedge \eta_{23} = \eta_{12} \wedge \eta_{13} + \eta_{13} \wedge \eta_{23}$.

\emph{Differential computation.}
Let
\[
x=[a|a|a]\otimes\eta_{12}\wedge\eta_{23}
\]
be the generator-level chain; mode labels insert the usual central
pairing $[a_m,a_n]=k\,m\,\delta_{m+n,0}$ in the contracted slots.
The Heisenberg OPE has no simple pole, so the bracket channel
$a_{(0)}a$ vanishes. The double-pole channel
$a_{(1)}a=k\mathbf 1$ does not disappear from the geometric bar
complex. In the Fulton--MacPherson residue convention of
Theorem~\ref{thm:frame-heisenberg-bar}, the adjacent collision
divisors contribute
\[
d_{\mathrm{res}}x
 =
k\,[a]\otimes\eta_{23}
-k\,[a]\otimes\eta_{13},
\]
up to the orientation convention fixed in
Lemma~\ref{lem:orientation-freefields}. The non-adjacent divisor
$D_{13}$ vanishes on this chain basis by form-squaring. The
de~Rham component is zero because the logarithmic forms are closed.

Thus the degree-$3$ top-form chain is not a cycle in general. The
central Heisenberg class is the class left by the alternating
adjacent-residue complex: applying $d_{\mathrm{res}}$ once more gives
$k^2-k^2=0$, and this cancellation is the Arnold relation in this
chain basis. The bar cohomology is the dual coalgebra
$\cH_k^i=H^\star B(\cH_k)$; the Verdier--Koszul dual algebra is
obtained only after dualizing that coalgebra, and carries curvature
$m_0=-k\,\omega$.
\end{computation}

\begin{computation}[Heisenberg: degree 4 differential]\label{comp:heisenberg-deg4}
For degree 4 with 4 tensor factors:
\[
\B_4(\cH) \ni [a_{-m}|a_{-n}|a_{-p}|a_{-q}] \otimes \omega_4
\]

The form space $\Omega^3(\overline{\Conf}_4)$ has dimension 6, generated by
products of three $\eta_{ij}$.

The differential $d_{\text{res}}$ involves the six collision divisors
$D_{ij}$ ($1 \leq i < j \leq 4$). On the chain form
$\eta_{12}\wedge\eta_{23}\wedge\eta_{34}$ the adjacent divisors
$D_{12},D_{23},D_{34}$ contribute the alternating central contractions
coming from $a_{(1)}a=k\mathbf 1$, while the non-adjacent divisors
vanish by form-squaring. In field notation,
\[
d_{\mathrm{res}}\bigl([a|a|a|a]\otimes
\eta_{12}\wedge\eta_{23}\wedge\eta_{34}\bigr)
 =
k\,[a|a]\otimes\eta_{23}\wedge\eta_{34}
-k\,[a|a]\otimes\eta_{13}\wedge\eta_{34}
+k\,[a|a]\otimes\eta_{12}\wedge\eta_{24},
\]
with the signs determined by the FM boundary orientation.

The degree-$4$ top-form channel is therefore a piece of the
Heisenberg residue complex, not a cycle subcoalgebra. The Arnold
relation supplies $d_{\mathrm{res}}^2=0$. This statement is about
the raw bar complex. It does not identify $A^!$ or the derived
chiral centre.
\end{computation}


\subsection{Twisting morphism and curvature}

\begin{computation}[Maurer--Cartan through degree 4]The twisting morphism satisfies the bar--cobar Maurer--Cartan equation
in the curved convention determined by the Heisenberg pairing.

\emph{Setup.} The twisting morphism $\tau: \B(\cH) \to \cH$ is:
\[
\tau([a_{-n}]) = a_{-n}, \quad \tau(\text{higher}) = 0
\]

The convolution product $\tau \star \tau: \B(\cH) \to \cH$ is computed via:
\[
(\tau \star \tau)(x) = \mu(\tau \otimes \tau)\Delta(x)
\]
where $\Delta$ is the coproduct on $\B(\cH)$ and $\mu$ is the product on $\cH$.

\emph{On degree~2.}
\[
\Delta([a_{-m}|a_{-n}]) = 1 \otimes [a_{-m}|a_{-n}] + [a_{-m}] \otimes [a_{-n}] + [a_{-m}|a_{-n}] \otimes 1
\]
(the standard deconcatenation coproduct; the terms $1 \otimes (\cdots)$ and $(\cdots) \otimes 1$ give zero under $\tau \otimes \tau$ since $\tau$ vanishes on degree $\neq 1$).

Thus the quadratic convolution records the normally ordered product,
not the central mode commutator:
\[
(\tau \star \tau)([a_{-m}|a_{-n}])
=a_{-m}a_{-n}.
\]
The central extension enters through the residue part of the bar
differential:
\[
a_{(0)}a=0,\qquad a_{(1)}a=k\mathbf 1.
\]
For negative creation modes alone there is no condition $m+n=0$.
The identity $[a_{-m},a_m]=k\,m$ belongs to the full Heisenberg
mode algebra and cannot be used to manufacture curvature inside the
raw reduced bar complex.

The precise separation is:
\[
\Omega B(\cH_k)\simeq \cH_k
\qquad\text{(bar--cobar inversion),}
\]
whereas
\[
\cH_k^! \simeq
\bigl(\mathrm{Sym}^{\mathrm{ch}}(V^*),\,m_0=-k\,\omega\bigr)
\qquad\text{(Verdier--Koszul dual algebra).}
\]
Consequently
\[
\kappa(\cH_k)=k,\qquad
\kappa(\cH_k^!)=-k,
\]
and $\kappa(\cH_k)+\kappa(\cH_k^!)=0$. The sign belongs to the
Verdier-dual lane, not to the bar--cobar counit.
\end{computation}


\section{Kac--Moody computations}

\begin{computation}[\texorpdfstring{$\widehat{\mathfrak{sl}}_3$}{sl-hat_3} OPE]\label{comp:sl3-ope}
The Lie algebra $\mathfrak{sl}_3$ has basis:
\[
\{H_1, H_2, E_1, E_2, E_3, F_1, F_2, F_3\}
\]
where $E_3 = [E_1, E_2]$ and $F_3 = [F_2, F_1]$.

\emph{Cartan matrix.}
\[
A = \begin{pmatrix} 2 & -1 \\ -1 & 2 \end{pmatrix}
\]

\emph{Structure constants.}
\begin{align*}
[H_i, E_j] &= A_{ij} E_j, \quad [H_i, F_j] = -A_{ij} F_j \\
[E_i, F_j] &= \delta_{ij} H_i \quad (i,j \in \{1,2\}) \\
[E_1, E_2] &= E_3, \quad [F_2, F_1] = F_3 \\
[E_1, E_3] &= 0 = [E_2, E_3], \quad \text{similarly for } F
\end{align*}

The normalized trace form is $(H_i, H_j) = \mathrm{tr}_{\mathrm{fund}}(H_i H_j) = A_{ij}$, $(E_i, F_j) = \delta_{ij}$.
(The Killing form is $\kappa = 2h^\vee \cdot \mathrm{tr}_{\mathrm{fund}} = 6\,\mathrm{tr}_{\mathrm{fund}}$ for $\mathfrak{sl}_3$;
the OPEs below use the normalized form with level $k$ measured in units of the long root.)

The affine OPEs at level $k$ are:
\begin{align*}
H_i(z) H_j(w) &\sim \frac{kA_{ij}}{(z-w)^2} \\
H_i(z) E_j(w) &\sim \frac{A_{ij} E_j(w)}{z-w} \\
E_i(z) F_j(w) &\sim \frac{k\delta_{ij}}{(z-w)^2} + \frac{\delta_{ij} H_i(w)}{z-w} \\
E_1(z) E_2(w) &\sim \frac{E_3(w)}{z-w} \\
E_1(z) E_3(w) &\sim 0 \quad \text{(since $[E_1, E_3] = 0$ in $\mathfrak{sl}_3$)}
\end{align*}
\end{computation}

\begin{computation}[\texorpdfstring{$\widehat{\mathfrak{sl}}_3$}{sl-hat_3} bar differential]\label{comp:sl3-bar}
The degree-2 differential is as follows.

With constant form:
\begin{align*}
d[H_{1,-m}|E_{1,-n}] &= 2[E_{1,-(m+n)}] \\
d[H_{1,-m}|E_{2,-n}] &= -[E_{2,-(m+n)}] \\
d[E_{1,-m}|E_{2,-n}] &= [E_{3,-(m+n)}] \\
d[E_{1,-m}|F_{1,-n}] &= [H_{1,-(m+n)}] + km\delta_{m+n,0}[1]
\end{align*}

With $\eta_{12}$: All differentials vanish (no simple poles when paired with $d\log$).

The degree-3 differential encodes the Serre relations.

The simple-root Serre relation appears as the following bar cycle:
\[
[E_{1,-\ell}|E_{1,-m}|E_{2,-n}] - 2[E_{1,-\ell}|E_{2,-m}|E_{1,-n}] + [E_{2,-\ell}|E_{1,-m}|E_{1,-n}]
\]

The differential of this combination gives the Serre identity:
\[
d(\text{above}) = E_1 E_1 E_2 - 2 E_1 E_2 E_1 + E_2 E_1 E_1 = (\text{ad}\, E_1)^2(E_2) = 0 \quad \text{in } U(\hat{\mathfrak{sl}}_3)
\]
\end{computation}


\subsection{Twisted acyclicity}

\begin{computation}[Kac--Moody acyclicity for the universal affine module]\label{comp:km-acyclic}
The twisted Koszul complex $\B(\hat{\mathfrak{sl}}_2) \otimes_\tau V_k$
is acyclic for the universal affine vacuum module at every level $k$.

\emph{Filtration.} Define $F_p$ by the PBW degree (number of generators applied to vacuum).

\emph{Associated graded.}
\[
\mathrm{gr}_F(\B(\hat{\mathfrak{sl}}_2) \otimes_\tau V_k) \cong
B\bigl(U(\mathfrak{sl}_2[t^{-1}]t^{-1})\bigr)
\otimes_{\tau_0} U(\mathfrak{sl}_2[t^{-1}]t^{-1})
\]
with $\tau_0$ the standard bar twisting cochain. This is the ordinary
augmented bar resolution of the algebra
$U(\mathfrak{sl}_2[t^{-1}]t^{-1})$.

\emph{$E_1$ page.} The homology of the associated graded complex is
\[
E_1^{p,q} =
\begin{cases}
\mathbb{C} & p=q=0,\\
0 & \text{otherwise}.
\end{cases}
\]
The vanishing is the exactness of the bar resolution, not a Whitehead
lemma computation. Whitehead's lemmas do not annihilate
$H_\bullet(\mathfrak{sl}_2;M)$ in all positive degrees for arbitrary
coefficient modules.

\emph{Convergence.} The filtration is bounded below on each conformal
weight and exhaustive, hence the spectral sequence gives:
\[
H_n(\B(\hat{\mathfrak{sl}}_2) \otimes_\tau V_k) = \begin{cases}
\bC & n = 0 \\
0 & n > 0
\end{cases}
\]

\emph{Level scope.} This is the bar-cobar resolution of the universal
affine module, not the raw critical centre and not a statement about a
simple admissible quotient. At $k=-h^\vee=-2$ the Sugawara denominator
vanishes and the completed critical vacuum algebra carries the
Feigin--Frenkel centre. Bar-cobar inversion itself has no pole: the
classes that jump live in the critical centre/Hochschild lane. At
non-critical admissible levels the universal module still has the PBW
bar resolution; the simple quotient $L_k$ requires a quotient-specific
associated-variety and spectral-sequence check.
\end{computation}


\section{$\mathcal{W}$-algebra computations}

\subsection{Virasoro bar complex through degree 3}

\begin{computation}[Virasoro OPE products]
\index{Virasoro algebra!bar complex}
The Virasoro algebra has a single strong generator $T$ (conformal weight~$2$) with OPE:
\begin{align*}
T(z)\,T(w) &= \frac{c/2}{(z-w)^4} + \frac{2T(w)}{(z-w)^2} + \frac{\partial T(w)}{z-w} + \normord{TT}(w) + O(z-w)
\end{align*}
The $n$-th products (singular terms):
\begin{center}
\begin{tabular}{clll}
$n$ & $T_{(n)}T$ & Pole order & Role \\
\hline
$3$ & $c/2$ & $(z\!-\!w)^{-4}$ & curvature \\
$2$ & $0$ & $(z\!-\!w)^{-3}$ & (absent) \\
$1$ & $2T$ & $(z\!-\!w)^{-2}$ & conformal weight \\
$0$ & $\partial T$ & $(z\!-\!w)^{-1}$ & translation
\end{tabular}
\end{center}

\emph{Key difference from Kac--Moody:} the affine Lie algebra $\widehat{\mathfrak{g}}_k$ has OPE poles up to order~$2$, while the Virasoro has a quartic pole. This extra pole produces the curvature element $m_0 = (c/2)\cdot\mathbf{1}$ in the $A_\infty$ structure (cf.\ Computation~\ref{comp:virasoro-bar-diff} below).
On the uniform-weight scalar lane, this curvature $m_0 = c/2$ is
the scalar component of the
universal Maurer--Cartan class
$\Theta_{\mathrm{Vir}_c}^{\min} = (c/2)\cdot\eta\otimes\Lambda$
(Theorem~\ref{thm:explicit-theta}). At $c = 0$ (the Witt algebra),
$\Theta_{\mathrm{Vir}_0}^{\min} = 0$ and the bar complex is
uncurved; this does not determine the higher-degree shadow obstruction tower. At
general~$c$, the scalar curvature propagates to genus~$1$ universally
and, for this single-generator algebra, to all genera via the Hodge classes
$\lambda_g$.
\end{computation}

\begin{computation}[Virasoro vacuum module: low-weight states]\label{comp:virasoro-vacuum}
The augmentation ideal $\bar{V} = V/\mathbb{C}\cdot|0\rangle$ of the Virasoro vacuum module has a basis of states created by $L_{-n}$ (for $n \geq 2$) applied to the vacuum:
\begin{center}
{\small
\begin{tabular}{c|ccccccccccccccccccc}
weight $h$ & 2 & 3 & 4 & 5 & 6 & 7 & 8 & 9 & 10 & 11 & 12 & 13 & 14 & 15 & 16 & 17 & 18 & 19 & 20 \\
\hline
$\dim \bar{V}_h$ & 1 & 1 & 2 & 2 & 4 & 4 & 7 & 8 & 12 & 14 & 21 & 24 & 34 & 41 & 55 & 66 & 88 & 105 & 137
\end{tabular}
}
\end{center}
where $p_{\geq 2}(h) = p(h) - p(h-1)$ denotes the number of partitions of~$h$ into parts $\geq 2$, with generating function $\prod_{n=2}^{\infty}(1-q^n)^{-1}$. The constraint arises from $L_{-1}|0\rangle = 0$ (translation invariance).
\end{computation}

\begin{computation}[Bar differential: degree~2 \texorpdfstring{$\to$}{math} degree~1]\label{comp:virasoro-bar-diff}
\index{bar differential!Virasoro}
The bar complex $\bar{B}(\mathrm{Vir}_c)$ has differential $D$ whose action on degree-2 elements extracts the full OPE data (all singular terms), as in the Kac--Moody formula~\eqref{eq:bar-diff-general}:
\begin{align}
D(T \otimes T \otimes \eta_{12}) &= T_{(3)}T \cdot |0\rangle + T_{(2)}T + T_{(1)}T + T_{(0)}T \notag \\
&= \frac{c}{2}\cdot|0\rangle + 0 + 2T + \partial T \end{align}
The three contributions correspond to:
\begin{itemize}
\item $c/2\cdot|0\rangle \in \bar{B}^0$: the \emph{vacuum term} from the quartic pole (curvature)
\item $2T \in \bar{B}^1$: the \emph{conformal weight term} from the double pole
\item $\partial T \in \bar{B}^1$: the \emph{translation term} from the simple pole
\end{itemize}

\emph{Comparison with Kac--Moody:}
\begin{center}
\begin{tabular}{lll}
\textbf{Algebra} & \textbf{$D(v \otimes w \otimes \eta)$} & \textbf{Highest OPE pole} \\
\hline
Heisenberg $\mathcal{H}_\kappa$ & $\kappa\cdot|0\rangle$ & double ($z^{-2}$) \\
$\widehat{\mathfrak{sl}}_{2,k}$ & $k(J^a,J^b)|0\rangle + f^{ab}{}_c J^c$ & double ($z^{-2}$) \\
Virasoro $\mathrm{Vir}_c$ & $(c/2)|0\rangle + 2T + \partial T$ & quartic ($z^{-4}$)
\end{tabular}
\end{center}

At weight $h = 5$, the bar differential on mixed elements $T \otimes \partial T$ and $\partial T \otimes T$:
\begin{align}
D(T \otimes \partial T \otimes \eta_{12}) &= T_{(3)}(\partial T) \cdot |0\rangle + T_{(1)}(\partial T) + T_{(0)}(\partial T) \notag \\
&= 0 + 3\partial T + \partial^2 T \end{align}
using $T_{(3)}(\partial T) = \partial(T_{(3)}T) + 3 T_{(2)}T = 0 + 0 = 0$, $T_{(1)}(\partial T) = \partial(T_{(1)}T) + T_{(0)}T = 2\partial T + \partial T = 3\partial T$, and $T_{(0)}(\partial T) = \partial(T_{(0)}T) = \partial^2 T$.
\end{computation}

\begin{computation}[Curvature element and \texorpdfstring{$c=0$}{c=0} special value]\label{comp:virasoro-curvature}
The curvature element of the Virasoro bar complex is:
\begin{equation}m_0 = \frac{c}{2} \in \bar{B}^0 = \mathbb{C}
\end{equation}
arising from the quartic pole $T_{(3)}T = c/2$.

The $A_\infty$ relation $m_1^2 = [m_0, -]$ gives:
\begin{equation}
m_1^2(v) = m_2(m_0, v) - m_2(v, m_0)
\end{equation}
for any $v \in \bar{V}$. The curvature controls the failure of $m_1$ (the linearized bar differential) to be a differential: $m_1^2 \neq 0$ unless $c = 0$.

\emph{Three regimes.}
\begin{itemize}
\item $c = 0$: Uncurved. The bar complex is a genuine cochain complex ($D^2 = 0$ on $\bar{V}$). The \emph{quadratic} OPE is self-dual in the classical sense ($\mathrm{Vir}_0^{!_{\mathrm{quad}}} \simeq \mathrm{Vir}_0$, Theorem~\ref{thm:virasoro-self-duality}); however, the same-family scalar shadow of the chiral Koszul dual lies at central charge~$26$, not at~$0$. The $TT$ OPE is purely quadratic ($2T/(z-w)^2 + \partial T/(z-w)$). The scalar self-dual point of this same-family shadow is $c = 13$ (Example~\ref{ex:virasoro-koszul-dual}).
\item $c \neq 0$, generic: Curved $A_\infty$. The completed bar construction is needed (Appendix~\ref{app:nilpotent-completion}). The Verdier--Koszul dual algebra is obtained from $A^i=H^\star B(A)$ by duality; its Virasoro same-family scalar shadow has central charge $26-c$ (Example~\ref{ex:virasoro-koszul-dual}), where the shift $c \mapsto 26-c$ arises from the Verdier duality pairing on $\overline{C}_{n+1}(X)$ interacting with the quartic pole data.
\item $c = 26$: The same-family shadow of the dual lane lies at $c=0$, the uncurved Virasoro value. Equivalently, after adjoining the ghost sector, the critical bosonic string satisfies the genus-$0$ BRST/bar comparison of Theorem~\ref{thm:brst-bar-genus0}.
\end{itemize}

\emph{The curvature--central charge dictionary:} the scalar curvature
of the Virasoro bar complex is
$m_0=\kappa(\mathrm{Vir}_c)=c/2$. The remaining degree-$2$
residue terms $2T$ and $\partial T$ are fixed by the conformal
structure, while the same-family Verdier partner has scalar
curvature $(26-c)/2$. Thus $c$ determines the universal Virasoro
OPE, but the table records only the scalar curvature lane when it
writes $\kappa(\mathrm{Vir}_c)=c/2$.
\end{computation}

\begin{computation}[\texorpdfstring{$\mathcal{W}_3$}{W_3} OPE coefficients]The $\mathcal{W}_3$ algebra has generators $T$ (spin 2) and $W$ (spin 3), with OPEs given in Definition~\ref{def:w3-algebra} and the complete $W$--$W$ expansion in Theorem~\ref{thm:w-w-ope-complete}.

\emph{Verification of Jacobi.} The $T$-$W$-$W$ Jacobi identity gives:
\[
[L_m, [W_n, W_p]] + \text{cyclic} = 0
\]
Using $[L_m, W_n] = (2m - n)W_{m+n}$:
\begin{align*}
[L_m, [W_n, W_p]] &= [L_m, \text{(sum of } L_{n+p}, \Lambda_{n+p}, \text{etc.)}] \\
&= (m - n - p)[\cdots] + \cdots
\end{align*}
The coefficient $\frac{16}{22+5c}$ is uniquely determined by requiring Jacobi to hold.

\emph{Central charge values.}
\begin{itemize}
\item $c = -2$: A non-unitary $W_3$ algebra. The DS parametrization gives $k = -3/2$, corresponding to $(p,q) = (3,2)$ with $\gcd(3,2) = 1$, but $W_3$ minimal models require $p, q \geq 3$, so $c = -2$ lies outside the minimal model series. (Not to be confused with symplectic fermions, which also have $c=-2$ but are a Virasoro, not $W_3$, vertex algebra.)
\item $c = 0$: Trivial/degenerate $W_3(4,3)$ model (not to be confused with symplectic fermions, which have $c=-2$).
\item $c = 2$: Free field (Heisenberg + bc ghosts).
\item $c \to \infty$: Classical limit ($\mathcal{W}_3^{\mathrm{cl}}$).
\end{itemize}
\end{computation}


\begin{computation}[DS reduction for \texorpdfstring{$\mathcal{W}_3$}{W_3}]\label{comp:ds-w3}
Starting from $\hat{\mathfrak{sl}}_3$ at level $k$:

\emph{BRST complex.}
\[
C^\bullet = V_k(\mathfrak{sl}_3) \otimes \mathrm{Cl}(\mathfrak{n}((t)) \oplus \mathfrak{n}^*((t)))
\]
where $\mathfrak{n} = \bC E_1 \oplus \bC E_2 \oplus \bC E_3$ is the nilpotent radical ($E_3 = [E_1, E_2]$) and $\mathrm{Cl}$ denotes the Clifford algebra (ghost--anti-ghost system).

\emph{Ghost fields.} $c_1, c_2, c_3$ (fermionic) with $|c_i| = 1$.

\emph{BRST differential.}
\begin{align*}
Q &= \sum_n (E_{1,n} - \chi_1\delta_{n,0})c_{1,-n} + \sum_n (E_{2,n} - \chi_2\delta_{n,0})c_{2,-n} \\
&\quad + \sum_n E_{3,n} c_{3,-n} + \sum_{m,n} c_{1,m}c_{2,n}b_{3,-(m+n)} + \cdots
\end{align*}
where $\chi_1, \chi_2$ determine the nilpotent element $f = \chi_1 F_1 + \chi_2 F_2$.

\emph{Principal nilpotent.} $\chi_1 = \chi_2 = 1$.

\emph{$H^0(Q)$.} Generated by:
\begin{itemize}
\item $T = $ Sugawara tensor (survives reduction)
\item $W = $ new spin-3 generator from reduction
\end{itemize}

\emph{Central charge formula.}
\[
c_{\mathcal{W}_3}(k) = 2 - \frac{24(k+2)^2}{k+3}
\]
\end{computation}


\subsection{\texorpdfstring{$\mathcal{W}_3$ bar complex: degree-2 computation}{W-3 bar complex: degree-2 computation}}

\begin{computation}[Complete \texorpdfstring{$n$}{n}-th products for \texorpdfstring{$\mathcal{W}_3$}{W_3}]\label{comp:w3-nthproducts}
\index{W3 algebra@$\mathcal{W}_3$!$n$-th products}
The $\mathcal{W}_3$ algebra has generators $T$ (weight~$2$) and $W$ (weight~$3$).
The singular $n$-th products ($a_{(n)}b$ for $n \geq 0$ with nonzero coefficient)
for each ordered generator pair are:

\emph{$T \times T$ (quartic pole).}
\begin{align*}
T_{(3)}T &= c/2, & T_{(2)}T &= 0, &
T_{(1)}T &= 2T, & T_{(0)}T &= \partial T.
\end{align*}

\emph{$T \times W$ (double pole).}
Since $W$ is primary of weight~$3$, the $TW$ OPE has at most a double pole:
\begin{align*}
T_{(1)}W &= 3W, & T_{(0)}W &= \partial W.
\end{align*}
All $T_{(n)}W = 0$ for $n \geq 2$.

\emph{$W \times T$ (double pole).}
By the vertex algebra skew-symmetry formula
\[
b_{(n)}a = \sum_{j \geq 0}
\frac{(-1)^{n+j+1}}{j!}\,\partial^j(a_{(n+j)}b)\colon
\]
\begin{align*}
W_{(1)}T &= T_{(1)}W = 3W, \\
W_{(0)}T &= -T_{(0)}W + \partial(T_{(1)}W) = -\partial W + 3\partial W = 2\partial W.
\end{align*}
The asymmetry $W_{(0)}T = 2\partial W \neq \partial W = T_{(0)}W$ is a consequence
of the derivative correction in the skew-symmetry formula. All $W_{(n)}T = 0$ for $n \geq 2$.

\emph{$W \times W$ (sixth-order pole).}
\begin{center}
\begin{tabular}{clll}
$n$ & $W_{(n)}W$ & Pole order & Weight \\
\hline
$5$ & $c/3$ & $(z\!-\!w)^{-6}$ & $0$ \\
$4$ & $0$ & $(z\!-\!w)^{-5}$ & -- \\
$3$ & $2T$ & $(z\!-\!w)^{-4}$ & $2$ \\
$2$ & $\partial T$ & $(z\!-\!w)^{-3}$ & $3$ \\
$1$ & $\tfrac{3}{10}\partial^2 T + \tfrac{16}{22+5c}\Lambda$ & $(z\!-\!w)^{-2}$ & $4$ \\
$0$ & $\tfrac{1}{15}\partial^3 T + \tfrac{8}{22+5c}\partial\Lambda$ & $(z\!-\!w)^{-1}$ & $5$
\end{tabular}
\end{center}
where $\Lambda = \normord{TT} - \frac{3}{10}\partial^2 T$ (Definition~\ref{def:lambda-complete}).
The vanishing $W_{(4)}W = 0$ occurs because $\dim \bar{V}_1 = 0$ in the vacuum module
($L_{-1}|0\rangle = 0$), so there is no weight-$1$ state to appear at the fifth-order pole.
\end{computation}


\begin{computation}[\texorpdfstring{$\mathcal{W}_3$}{W_3} bar differential on degree-2 elements]\label{comp:w3-bar-degree2}
\index{bar differential!$\mathcal{W}_3$}
The bar differential $D \colon \bar{B}^2(\mathcal{W}_3) \to \bar{B}^0 \oplus \bar{B}^1$
extracts all singular OPE data from each generator pair.
The four ordered generator pairs give the following values of
$D(a \otimes b \otimes \eta_{12})$.

\emph{(i)} $D(T \otimes T \otimes \eta_{12})$
(from Computation~\ref{comp:virasoro-bar-diff}):
\[
D(T \otimes T \otimes \eta_{12}) = \frac{c}{2}\cdot|0\rangle + 2T + \partial T
\]

\emph{(ii)} $D(T \otimes W \otimes \eta_{12})$:
\[
D(T \otimes W \otimes \eta_{12}) = T_{(1)}W + T_{(0)}W = 3W + \partial W
\]
No vacuum contribution (the $TW$ OPE has no pole beyond order~$2$).

\emph{(iii)} $D(W \otimes T \otimes \eta_{12})$:
\[
D(W \otimes T \otimes \eta_{12}) = W_{(1)}T + W_{(0)}T = 3W + 2\partial W
\]
The asymmetry with~(ii) comes from skew-symmetry
(Computation~\ref{comp:w3-nthproducts}).

\emph{(iv)} $D(W \otimes W \otimes \eta_{12})$:
\begin{align*}
D(W \otimes W \otimes \eta_{12})
&= W_{(5)}W \cdot |0\rangle + W_{(3)}W + W_{(2)}W + W_{(1)}W + W_{(0)}W \\
&= \frac{c}{3}\cdot|0\rangle + 2T + \partial T
 + \Bigl[\frac{3}{10}\partial^2 T + \frac{16}{22+5c}\Lambda\Bigr] \\
&\quad + \Bigl[\frac{1}{15}\partial^3 T + \frac{8}{22+5c}\partial\Lambda\Bigr]
\end{align*}
(The $n=4$ term vanishes: $W_{(4)}W = 0$.)
This is the most complex bar differential in the manuscript: it produces components
in every weight from~$0$ (vacuum) through~$5$ ($\partial\Lambda$).

\emph{Summary of vacuum contributions:}
\begin{center}
\begin{tabular}{lcc}
\textbf{Generator pair} & \textbf{Vacuum term} & \textbf{Max weight in $\bar{B}^1$} \\
\hline
$T \otimes T$ & $c/2$ & $3$ ($\partial T$) \\
$T \otimes W$ & $0$ & $4$ ($\partial W$) \\
$W \otimes T$ & $0$ & $4$ ($2\partial W$) \\
$W \otimes W$ & $c/3$ & $5$ ($\partial\Lambda$)
\end{tabular}
\end{center}
Only the ``diagonal'' pairs $T \otimes T$ and $W \otimes W$ produce vacuum leakage.
\end{computation}


\begin{computation}[\texorpdfstring{$\mathcal{W}_3$}{W_3} curvature and dual level complementarity]\index{curvature!$\mathcal{W}_3$ bar complex}
The $\mathcal{W}_3$ bar complex has vacuum leakage from two independent channels:
\begin{enumerate}
\item \emph{Virasoro sector}: $T_{(3)}T = c/2$ (quartic pole), contributing $m_0^{(T)} = c/2$.
\item \emph{$W$-sector}: $W_{(5)}W = c/3$ (sixth-order pole), contributing $m_0^{(W)} = c/3$.
\end{enumerate}
The ratio $m_0^{(W)}/m_0^{(T)} = 2/3$ is independent of the level and equals the
ratio of central term normalizations in the $WW$ and $TT$ OPEs
($c/3$ vs.\ $c/2$).

Both leakages vanish simultaneously if and only if $c = 0$.
The affine pre-DS curvature is proportional to $k+h^\vee=k+3$,
whereas the reduced $\mathcal{W}_3$ vacuum leakages above are
controlled by the central charge $c(k)$. These are distinct
conditions:
\begin{itemize}
\item $c = 0$ occurs at rational levels $k = -5/3$ and $k = -9/4$
 (from solving $12(k+2)^2 = k+3$, i.e., $12u^2 - u - 1 = 0$ where $u = k+2$).
\item $k = -3$ is the critical level where the Sugawara tensor is undefined and
 the DS central-charge formula has a pole.
\end{itemize}
Similarly, for the Virasoro ($\mathfrak{sl}_2$), $c = 0$ occurs at
$k = -1/2$ and $k = -4/3$ (both rational), while critical level is $k = -2$.
See Theorem~\ref{thm:w-bar-curvature} for the full discussion.

At the dual level $k' = -k-6$, the dual central charge is $c' = c(k') = 2 - 24(k'+2)^2/(k'+3) = 2 + 24(k+4)^2/(k+3)$.
The sum:
\[
c + c' = 4 + 24\cdot\frac{(k+4)^2 - (k+2)^2}{k+3}
 = 4 + 24\cdot\frac{4(k+3)}{k+3} = 100
\]
giving the general complementarity $c(\mathcal{W}^k) + c(\mathcal{W}^{k'})
= 2r + 4h^\vee d = 2(2) + 4(3)(8) = 100$
(Theorem~\ref{thm:central-charge-complementarity}).

At level $k'$, the vacuum leakages are $m_0^{(T)} = c'/2$ and $m_0^{(W)} = c'/3$.
Since $c' = 100 - c$, the sum of curvatures satisfies:
\[
m_0^{(T)}(k) + m_0^{(T)}(k') = \frac{c}{2} + \frac{100-c}{2} = 50
\]
The curvature of one algebra determines the curvature of its dual. When
$c = 0$ (uncurved at $k = -5/3$ or $-9/4$), the dual has $c' = 100$
(strongly curved). The involution $k \mapsto -k-6$ has fixed point $k = -3$
(the critical level), where the Sugawara construction is undefined and the DS construction degenerates;
the complementarity formula $c + c' = 100$ governs the reduced
$\mathcal{W}_3$ curvature trade-off for all non-critical levels but
has no $\mathcal{W}$-algebraic meaning at the critical level itself.

The composite field coefficient at level $k$ is $\alpha(c) = 16/(22+5c)$, and
at dual level $\alpha(c') = 16/(22+5(100-c)) = 16/(522-5c)$.
These differ (the dual algebra has a different OPE structure), but both are
determined by a single parameter $c$. In the large-level limit $k \to \infty$
(where $c \to -\infty$), both $\alpha(c) \to 0^-$ and
$\alpha(c') \to 0^+$: both algebras become classical simultaneously.
\end{computation}


\section{\texorpdfstring{$\widehat{\mathfrak{sl}}_3$ bar complex: degree-3 Serre relations}{sl-3 bar complex: degree-3 Serre relations}}
\index{Serre relations!bar complex encoding}
\index{sl3-hat@$\widehat{\mathfrak{sl}}_3$!degree-3 bar complex}

The degree-2 bar differential for $\widehat{\mathfrak{sl}}_3$ was computed in
Computation~\ref{comp:sl3-bar}. The full degree-3
computation encodes the Serre relations of the $A_2$ root system.

\begin{computation}[Complete degree-3 bar differential for \texorpdfstring{$\widehat{\mathfrak{sl}}_3$}{sl-hat_3}]
\label{comp:sl3-degree3-complete}
\index{bar complex!sl3@$\mathfrak{sl}_3$!degree 3}

\emph{Setup.}
The degree-3 bar complex is
$\bar{B}^3(\widehat{\mathfrak{sl}}_{3,k}) = \mathfrak{sl}_3^{\otimes 3}
\otimes \Omega^2(\overline{\mathrm{Conf}}_3)$.
With $\dim(\mathfrak{sl}_3) = 8$ and
$\dim \Omega^2(\overline{\mathrm{Conf}}_3) = 2$
(cf.\ Computation~\ref{comp:heisenberg-deg3-full}), the total dimension is
$8^3 \times 2 = 1024$.

Use the basis $\{\eta_{12}\wedge\eta_{13},\; \eta_{13}\wedge\eta_{23}\}$
for $\Omega^2(\overline{\mathrm{Conf}}_3)$, with the Arnold relation
$\eta_{12}\wedge\eta_{23} = \eta_{12}\wedge\eta_{13} + \eta_{13}\wedge\eta_{23}$.

The differential $d\colon \bar{B}^3 \to \bar{B}^2$ acts by residues at three
collision divisors $D_{12}$, $D_{13}$, $D_{23}$:
\[
d = \operatorname{Res}_{D_{12}} + \operatorname{Res}_{D_{13}} + \operatorname{Res}_{D_{23}}
\]

\emph{(i) Cartan--root element: $[H_{1,-\ell}|E_{1,-m}|E_{2,-n}]$.}
On $\eta_{12} \wedge \eta_{13}$:
\begin{align*}
\operatorname{Res}_{D_{12}} &= [H_1, E_1]_{-(\ell+m)} \otimes E_{2,-n} \otimes \eta_{13}
= 2[E_{1,-(\ell+m)}|E_{2,-n}] \otimes \eta_{13} \\
\operatorname{Res}_{D_{13}} &= H_{1,-\ell} \otimes [H_1, E_2]_{-(m+n)} \otimes \eta_{12}
\quad\text{(transposition of labels)} \\
&\quad\text{but } [H_1, E_2] = A_{12}\,E_2 = -E_2, \text{ so this gives }
-[H_{1,-\ell}|E_{2,-(m+n)}] \otimes \eta_{12} \\
\operatorname{Res}_{D_{23}} &= H_{1,-\ell} \otimes [E_1, E_2]_{-(m+n)} \otimes \eta_{1?}
= [H_{1,-\ell}|E_{3,-(m+n)}] \otimes \eta_{13}
\end{align*}
where $E_3 = [E_1, E_2]$. The subscripts $\eta_{1?}$ require careful tracking:
after the collision $z_2 = z_3$, the surviving coordinates are $(z_1, z_{23})$,
with form $\eta_{1,23} = \eta_{13}|_{z_2=z_3}$.

Hence:
\begin{equation}d\bigl([H_{1}|E_1|E_2] \otimes \eta_{12}\wedge\eta_{13}\bigr)
= 2[E_1|E_2] \otimes \eta_{13}
- [H_1|E_2] \otimes \eta_{12}
+ [H_1|E_3] \otimes \eta_{13}
\end{equation}
(suppressing mode indices for clarity).

\emph{(ii) Serre relation element:}
\[
S = [E_{1,-\ell}|E_{1,-m}|E_{2,-n}] -
2[E_{1,-\ell}|E_{2,-m}|E_{1,-n}] + [E_{2,-\ell}|E_{1,-m}|E_{1,-n}].
\]

The key computation is the evaluation of $d(S)$ on
$\eta_{12}\wedge\eta_{13}$.

For the first term $[E_1|E_1|E_2]$:
\begin{align*}
\operatorname{Res}_{D_{12}}([E_1|E_1|E_2]) &= [E_1, E_1] \otimes E_2
= 0 \quad ([E_1, E_1] = 0) \\
\operatorname{Res}_{D_{13}}([E_1|E_1|E_2]) &= E_1 \otimes [E_1, E_2]
= [E_1|E_3] \\
\operatorname{Res}_{D_{23}}([E_1|E_1|E_2]) &= E_1 \otimes [E_1, E_2]
= [E_1|E_3]
\end{align*}

For the second term $-2[E_1|E_2|E_1]$:
\begin{align*}
\operatorname{Res}_{D_{12}}(-2[E_1|E_2|E_1]) &= -2[E_3|E_1] \\
\operatorname{Res}_{D_{13}}(-2[E_1|E_2|E_1]) &= 0
\quad ([E_1, E_1] = 0) \\
\operatorname{Res}_{D_{23}}(-2[E_1|E_2|E_1]) &= -2[E_1|E_3]
\end{align*}

For the third term $[E_2|E_1|E_1]$:
\begin{align*}
\operatorname{Res}_{D_{12}}([E_2|E_1|E_1]) &= -[E_3|E_1]
\quad ([E_2, E_1] = -E_3) \\
\operatorname{Res}_{D_{13}}([E_2|E_1|E_1]) &= [E_2|0] = 0
\quad ([E_1, E_1] = 0) \\
\operatorname{Res}_{D_{23}}([E_2|E_1|E_1]) &= [E_2|0] = 0
\end{align*}

Combining with the form $\eta_{12}\wedge\eta_{13}$,
the collision $D_{12}$ produces $\eta_{13}$,
the collision $D_{13}$ produces $\eta_{12}$,
and the collision $D_{23}$ produces $\eta_{13}$:
\begin{align*}
d(S \otimes \eta_{12}\wedge\eta_{13})
&= \bigl(0 - 2[E_3|E_1] - [E_3|E_1]\bigr) \otimes \eta_{13}
\quad\text{(from $D_{12}$)} \\
&\quad + \bigl([E_1|E_3] + 0 + 0\bigr) \otimes \eta_{12}
\quad\text{(from $D_{13}$)} \\
&\quad + \bigl([E_1|E_3] - 2[E_1|E_3] + 0\bigr) \otimes \eta_{13}
\quad\text{(from $D_{23}$)} \\
&= -3[E_3|E_1]\otimes\eta_{13}
+ [E_1|E_3]\otimes\eta_{12}
- [E_1|E_3]\otimes\eta_{13}
\end{align*}

Using the degree-2 differential from Computation~\ref{comp:sl3-bar}:
$d([E_1|E_3]) = [E_1, E_3] = 0$ and $d([E_3|E_1]) = [E_3, E_1] = 0$
(since $[E_1, [E_1, E_2]] = 0$ in $\mathfrak{sl}_3$; this is the
Serre relation for the $A_2$ root system).

Therefore $d(S \otimes \eta_{12}\wedge\eta_{13})$ lies in the kernel of
$d\colon \bar{B}^2 \to \bar{B}^1$. This is a non-trivial \emph{cycle}
in $\bar{B}^2$ produced by the Serre relation combination~$S$. At generic
level~$k$, this cycle is also a \emph{boundary}: the acyclicity
of the bar-twisted complex (Proposition~\ref{prop:km-generic-acyclicity})
implies it bounds a degree-3 chain via the curvature correction. At
critical level $k = -3$, it represents a non-trivial class in
$H_2(\bar{B}(\widehat{\mathfrak{sl}}_{3,-3}))$, corresponding to the
second fundamental invariant of the Feigin--Frenkel center.

\emph{(iii) Dual Serre relation:
$\tilde{S} = [E_{2}|E_{2}|E_{1}] - 2[E_{2}|E_{1}|E_{2}]
+ [E_{1}|E_{2}|E_{2}]$.}

By symmetry of the $A_2$ root system (the Dynkin diagram automorphism
$1 \leftrightarrow 2$), the computation for $\tilde{S}$ is obtained from
that of $S$ by exchanging $E_1 \leftrightarrow E_2$ and
$H_1 \leftrightarrow H_2$. The element $\tilde{S}$ encodes
$(\operatorname{ad} E_2)^2(E_1) = 0$, and $d(\tilde{S})$ is a cycle
in $\bar{B}^2$ by the same mechanism.
\end{computation}

\begin{proposition}[Serre tensors are quadratic syzygies, not the dual algebra;
\ClaimStatusProvedHere]
\label{prop:sl3-serre-cohomology}
\index{Serre relations!cohomological role}
For $\widehat{\mathfrak{sl}}_{3,k}$ at generic level $k \neq -3$:
\begin{enumerate}[label=\textup{(\roman*)}]
\item The tensors $S$ and $\tilde S$ are the two simple-root Serre
syzygies in the quadratic associated graded presentation. Their
images under $d\colon \bar{B}^3\to\bar{B}^2$ are degree-$2$ cycles
because $[E_i,[E_i,E_j]]=0$ in $\mathfrak{sl}_3$.
\item These tensors do not determine the full degree-$3$ bar
cohomology and do not identify the Verdier--Koszul dual algebra
$A^!$. They are relations inside the associated graded presentation
of $A$, while $A^i=H^\star B(A)$ is the dual coalgebra and
$A^!$ is obtained after graded linear/Verdier duality.
\item At critical level $k=-3$, the Feigin--Frenkel centre is generated
by the invariant polynomials of degrees $2$ and $3$ in the Langlands
dual Lie algebra. Its cubic generator is represented by the
Segal--Sugawara invariant, not by the raw Serre tensor itself.
\end{enumerate}
\end{proposition}

\begin{proof}
The displayed computation of $d(S)$ shows that all surviving
degree-$2$ terms are built from $[E_1,E_3]$ or $[E_3,E_1]$, hence
vanish after applying the next bracket differential because
$E_3=[E_1,E_2]$ and $(\operatorname{ad}E_1)^2(E_2)=0$.
The Dynkin involution gives the same statement for $\tilde S$.
This proves that the Serre tensors are syzygies in the
associated graded bar presentation.

The stronger statement that the Serre tensors determine the full
dual algebra is false: the Chevalley--Eilenberg cohomology
$H^3(\mathfrak{sl}_3;\mathbb C)$ is one-dimensional, whereas the
simple-root Serre list has two elements. These are different objects.
The twisted Koszul complex is acyclic at generic level by
Proposition~\ref{prop:km-generic-acyclicity}; the bar cohomology
$A^i=H^\star B(A)$ and the Verdier--Koszul dual algebra $A^!$ require
the dual coalgebra construction, not a count of Serre tensors.

At the critical level, Feigin--Frenkel identifies the centre with the
polynomial algebra on invariant polynomials of degrees $2$ and $3$.
The degree-$3$ invariant is the cubic Segal--Sugawara field. The Serre
tensor is one presentation-level relation whose compatibility is
required by that invariant theory; it is not itself the central
generator.
\end{proof}

\begin{computation}[Dimension table for \texorpdfstring{$\widehat{\mathfrak{sl}}_3$}{sl-hat_3}
bar complex]\index{bar complex!dimension table!sl3@$\mathfrak{sl}_3$}
The dimensions of the bar complex components $\bar{B}^n = \mathfrak{sl}_3^{\otimes n}
\otimes \Omega^{n-1}(\overline{\mathrm{Conf}}_n)$ through degree~5:
\begin{center}
\renewcommand{\arraystretch}{1.2}
\begin{tabular}{c|ccc|c}
$n$ & $\dim(\mathfrak{sl}_3^{\otimes n})$ &
$\dim \Omega^{n-1}(\overline{\mathrm{Conf}}_n)$ &
$\dim \bar{B}^n$ & Generating relations \\
\hline
$0$ & $1$ & $1$ & $1$ & vacuum \\
$1$ & $8$ & $1$ & $8$ & generators $J^a$ \\
$2$ & $64$ & $1$ & $64$ & OPE data \\
$3$ & $512$ & $2$ & $1024$ & Serre ($\times 2$) \\
$4$ & $4096$ & $6$ & $24576$ & higher Serre \\
$5$ & $32768$ & $24$ & $786432$ &
\end{tabular}
\end{center}
(The Poincar\'{e} polynomial of $\mathrm{Conf}_n(\mathbb{C})$ is
$\prod_{j=1}^{n-1}(1+jt)$, giving $\dim \Omega^{n-1} = (n-1)!$.)

\emph{Comparison with $\mathfrak{sl}_2$.}
\begin{center}
\renewcommand{\arraystretch}{1.2}
\begin{tabular}{c|cc|cc}
$n$ & $\dim \bar{B}^n(\widehat{\mathfrak{sl}}_2)$ &
$\dim \bar{B}^n(\widehat{\mathfrak{sl}}_3)$ &
Ratio & $= (8/3)^n \cdot \tfrac{n!_{\mathfrak{sl}_2}}{n!_{\mathfrak{sl}_3}}$ \\
\hline
$0$ & $1$ & $1$ & $1$ & $1$ \\
$1$ & $3$ & $8$ & $2.67$ & $8/3$ \\
$2$ & $9$ & $64$ & $7.11$ & $(8/3)^2$ \\
$3$ & $54$ & $1024$ & $18.96$ & $\approx (8/3)^3$ \\
$4$ & $486$ & $24576$ & $50.6$ & $\approx (8/3)^4$
\end{tabular}
\end{center}
The ratio $\dim \bar{B}^n(\widehat{\mathfrak{sl}}_3) /
\dim \bar{B}^n(\widehat{\mathfrak{sl}}_2) = (8/3)^n$
holds exactly, since both share the form factor $(n-1)!$
and the generator factor scales as $(\dim \mathfrak{g})^n$.
\end{computation}


\subsection{Weight-decomposed modular rank}

\begin{computation}[Modular rank of \texorpdfstring{$\widehat{\mathfrak{sl}}_3$}{sl-hat_3} bar differential;
\ClaimStatusConditional]
\label{comp:sl3-modular-rank}
\index{bar complex!modular rank}
\index{sl3-hat@$\widehat{\mathfrak{sl}}_3$!modular rank}

Decompose the bracket-part of the bar differential
$d_{\mathrm{bracket}}\colon \bar{B}^n \to \bar{B}^{n-1}$
by the $\mathfrak{h}$-weight $\mu \in \mathbb{Z}^2$ and
compute its rank over $\mathbb{Q}$ and $\mathbb{F}_p$ for
small primes. This is feasible because the weight decomposition
breaks the $8^n \times 8^{n-1}$ matrix into blocks of size
$\leq 346 \times 56$ (at weight $(0,0)$).

\emph{Degree~$3 \to 2$} (source $= 512$, target $= 64$).
\begin{center}
\renewcommand{\arraystretch}{1.2}
\begin{tabular}{c|cc}
& $\mathbb{Q}$ & finite-field comparison \\
\hline
$\mathrm{rk}(d \mid_{\mu=(0,0)})$ & $9$ &
must be $\leq 9$ for the same integral matrix \\
$\mathrm{rk}(d)$ (total) & $63$ &
must be $\leq 63$ for the same integral matrix
\end{tabular}
\end{center}
There is a $1$-dimensional rational cokernel at weight
$(0,0)$ for the exact matrix used here: the Killing form element
$\kappa_{\mathrm{Kill}} =
\sum_{a,b} \kappa^{ab} [J_a \mid J_b] \in \bar{B}^2_{(0,0)}$
is $d$-closed by invariance of the trace form. No finite-field rank
can exceed the rank over~$\mathbb{Q}$ for the same integral matrix:
if a minor is nonzero modulo~$p$, that minor is already nonzero over
$\mathbb{Q}$. Thus a rank jump $63\mapsto64$ is a normalization
change, not a modular anomaly theorem. The rational Killing class is the
bar-complex avatar of the genus-$1$ universal MC component
$\Theta^{(1)}_{\widehat{\mathfrak{sl}}_{3,k}}$
(Theorem~\ref{thm:explicit-theta}): the obstruction to
$d^2 = 0$ in bar degree~$2$ is controlled by $\kappa_{\mathrm{Kill}}$,
whose persistence as a rational cocycle is forced by the
non-vanishing of $\Theta$ at generic level.

The degree-$3$ table records only this rational cokernel. A finite-field
comparison is meaningful only after fixing the same integral
Chevalley basis and the same Orlik--Solomon sign normalization; changing
either changes the matrix rather than revealing a new modular class.

\emph{Degree~$4 \to 3$} (source $= 4096$, target $= 512$).
\begin{center}
\renewcommand{\arraystretch}{1.2}
\begin{tabular}{c|ccccc}
& $\mathbb{Q}$ & $\mathbb{F}_2$ & $\mathbb{F}_3$ &
$\mathbb{F}_5$ & $\mathbb{F}_7$ \\
\hline
$\mathrm{rk}(d)$ (total) & $512$ & $512$ & $512$ & $512$ & $512$
\end{tabular}
\end{center}
The total rank equals $\dim \bar{B}^3 = 512 = 8^3$: the differential
$d_{\mathrm{bracket}}\colon \bar{B}^4 \to \bar{B}^3$ is
\emph{surjective} in all characteristics. No modular anomalies
occur at degree~$4$.

\emph{Interpretation.} The degree-$3$ rational deficiency reflects the
Casimir obstruction: the Killing form element
$\kappa_{\mathrm{Kill}} \in \bar{B}^2_{(0,0)}$ is $d$-closed by centrality,
and over~$\mathbb{Q}$ it is not $d$-exact: the rank deficiency
$\operatorname{rk}(d) = 63 < 64$ shows that~$\kappa_{\mathrm{Kill}}$
survives in the degree-$2$ cohomology of this bracket-weight
truncation.
No conclusion about positive characteristic follows from matrices with
changed normalization. The surjectivity at degree~$4$ is a separate
low-degree computation in the bracket-plus-form model; it is not an
argument for concentration of the full bar cohomology.
\end{computation}


\begin{computation}[Chiral bracket rank with Orlik--Solomon forms; \ClaimStatusConditional]
\label{comp:sl3-chiral-bracket-os}
\index{bar complex!chiral bracket rank}
\index{Orlik--Solomon algebra!residue maps}

For comparison with Computation~\ref{comp:sl3-modular-rank}, the
rank of the \emph{full} chiral bracket
differential $d_{\mathrm{bracket}}\colon \bar{B}^n \to \bar{B}^{n-1}$
on the chain groups $\bar{B}^n = \mathfrak{g}^{\otimes n} \otimes
\mathrm{OS}^{n-1}(n)$, incorporating the Orlik--Solomon residue maps
$\operatorname{Res}_{D_{ij}}$ with signs $(-1)^{i+j}$, is computed
from the following matrix model.

\emph{Matrix model.}
The chiral bracket differential is built as a matrix of size
$\dim \bar{B}^n \times \dim \bar{B}^{n-1}$, where
$\dim \bar{B}^n = (\dim \mathfrak{g})^n \cdot (n-1)!$.
For each pair $(i,j)$ with $1 \leq i < j \leq n$, the bracket
$[\cdot, \cdot]_{ij}$ contracts slots $i$ and~$j$ using the
$\mathfrak{sl}_3$ structure constants, while the residue map
$\operatorname{Res}_{D_{ij}}\colon \mathrm{OS}^{n-1}(n) \to
\mathrm{OS}^{n-2}(n-1)$ restricts the differential form. The
sign $(-1)^{i+j}$ accounts for the orientation of the boundary
divisor $D_{ij} = \{z_i = z_j\}$ in the FM compactification.

\emph{Results.}
\begin{enumerate}[label=(\roman*)]
\item \emph{Degree~$2 \to 1$}: $\mathrm{rk} = 8 = \dim \mathfrak{sl}_3$.
 Agrees with both the CE and modular computations.
 (At degree~$2$, $\mathrm{OS}^0(2) = \mathbb{C}$, so the chiral
 and CE brackets coincide.)
\item \emph{Degree~$3 \to 2$}: $\mathrm{rk} = 63$.
 Matches Computation~\ref{comp:sl3-modular-rank} exactly.
 The CE bracket (without OS forms) gives rank~$56$; the chiral
 bracket with OS residue maps has seven additional image directions.
\item \emph{Degree~$4 \to 3$}: $\mathrm{rk} = 1024 = \dim \bar{B}^3$.
 The differential is \emph{surjective} in the chiral computation,
 agreeing with the modular rank result.
\end{enumerate}

Over $\mathbb{Q}$, the displayed degree-$4$ matrix has rank~$1024$.
Reducing the same integral matrix modulo $p=5,7,11,13$ preserves rank
$1024$. This finite-truncation statement concerns the
$24576 \times 1024$ matrix determined by the $\mathfrak{sl}_3$ structure
constants and the Orlik--Solomon residue maps; it does not assert
surjectivity in all bar degrees.
\end{computation}


\subsection{\texorpdfstring{$\widehat{\mathfrak{sl}}_3$ PBW target and dual Hilbert sequence}{sl-3 PBW target and dual Hilbert sequence}}
\index{sl3-hat@$\widehat{\mathfrak{sl}}_3$!spectral sequence}
\index{PBW spectral sequence!sl3@$\mathfrak{sl}_3$}

The direct computation of bar cohomology via matrix ranks
(Computations~\ref{comp:sl3-modular-rank}
and~\ref{comp:sl3-chiral-bracket-os}) becomes unwieldy beyond degree~$4$.
The PBW spectral sequence of Theorem~\ref{thm:pbw-koszulness-criterion}
identifies the acyclicity target of the twisted Koszul complex.

\begin{proposition}[PBW target for \texorpdfstring{$\widehat{\mathfrak{sl}}_3$}{sl-hat_3};
\ClaimStatusConditional]
\label{prop:sl3-pbw-ss}
The PBW filtration on the bar complex of $\widehat{\mathfrak{sl}}_{3,k}$ at
generic level $k \neq -h^\vee = -3$ yields a spectral sequence
\[
E_1^{p,q} =
H^{p+q}\!\left(
B(U(\mathfrak{sl}_3[t^{-1}]t^{-1}))\otimes_{\tau_0}
U(\mathfrak{sl}_3[t^{-1}]t^{-1})
\right)
\;\Longrightarrow\;
H^{p+q}\!\left(\barB(\widehat{\mathfrak{sl}}_{3,k})
\otimes_\tau V_k(\mathfrak{sl}_3)\right).
\]
Thus:
\begin{enumerate}[label=\textup{(\roman*)}]
\item This spectral sequence proves acyclicity of the twisted Koszul
 complex, not a formula for the raw bar cohomology
 $H^\star(\barB(\widehat{\mathfrak{sl}}_{3,k}))$.
\item The invariant tensors of $\mathfrak{sl}_3$ in degrees $2$ and $3$
 control the Casimir and cubic channels in the associated graded, but
 their presence does not imply that finite matrix ranks determine all
 bar cohomology degrees.
\item The low-degree numbers used later for the Koszul-dual Hilbert
 function are the conditional sequence
 \[
 \dim(A^!)_1=8,\quad \dim(A^!)_2=36,\quad
 \dim(A^!)_3=204.
 \]
 They concern the Verdier--Koszul dual algebra $A^!$ through
 $A^i=H^\star B(A)$, not the derived centre and not the twisted
 acyclic complex.
\end{enumerate}
\end{proposition}

\begin{proof}
The PBW filtration on $U(\widehat{\mathfrak{sl}}_3)$ induces a filtration
on the twisted complex
$\barB(\widehat{\mathfrak{sl}}_{3,k})\otimes_\tau V_k$. Its associated
graded object is the ordinary bar resolution of
$U(\mathfrak{sl}_3[t^{-1}]t^{-1})$. Hence the $E_1$ page is
concentrated on the augmentation line. This is the standard
bar-resolution argument and avoids the false inference that
Whitehead's lemmas compute all positive Lie algebra homology with the
loop-algebra coefficient system.

The raw bar cohomology is a different object. When $A$ is chirally
Koszul, $A^i=H^\star B(A)$ is the dual coalgebra and the dual algebra
$A^!$ is obtained by graded linear/Verdier duality. The numerical
sequence recorded in (iii) belongs to that dual Hilbert function and
retains its conditional scope in the table.
\end{proof}

\begin{remark}[PBW convergence target and the loop algebra filtration]
\label{rem:pbw-convergence-target}
The spectral sequence of Proposition~\ref{prop:sl3-pbw-ss} requires
careful identification of its convergence target. The $E_1$ page
has $E_1^{p,0} = S^p(\bar{V})^{\mathfrak{g}}$ where
$\bar{V} = \mathfrak{g}[t^{-1}]t^{-1}$.
Since the adjoint representation has no invariant vectors,
$E_1^{1,0} = \bar{V}^{\mathfrak{g}} = 0$
at \emph{every} weight $H \geq 1$. Yet $\dim H^1(\barB(\widehat{\mathfrak{g}}_k)) = 8$
by direct computation (Computation~\ref{comp:sl3-modular-rank}).

The resolution is that the PBW spectral sequence of
Theorem~\ref{thm:pbw-koszulness-criterion} proves acyclicity
of the \emph{Koszul complex} $K = \barB(A) \otimes_\tau A$,
not the bar complex $\barB(A)$ itself. The $E_1^{p,0}$ terms
are invariants in the symmetric algebra filtration of the
Koszul complex, and the convergence
$E_\infty \Rightarrow H^*(K)$ gives the Koszul criterion
$H^*(K) = \Bbbk$ (acyclicity), not the bar cohomology
$H^*(\barB(A))$ directly.

The conditional Koszul-dual Hilbert dimensions
$[1, 8, 36, 204, 1352, \ldots]$ are instead obtained from
Koszul acyclicity via the Hilbert series inversion
$H_A(t) \cdot H_{A^!}(-t) = 1$, using the Koszul dual Hilbert series
$H_{A^!}(t) = \sum_n \dim(A^!)_n \, t^n$.
The direct spectral sequence approach to $H^*(\barB(A))$
requires the \emph{Chevalley--Eilenberg cohomology of the
loop algebra} $H^*(\mathfrak{g}^-; \Bbbk)$, computed in
\texttt{ce\_cohomology\_loop.py}:
\[
\dim H^n(\mathfrak{g}^-; \Bbbk) =
[1, 8, 20, 70, 120, 125, \ldots],
\]
which satisfies $H^n_{\mathrm{CE}} \leq H^n_{\mathrm{bar}}$
at each degree, with ratios $1.00, 1.80, 2.91, 11.27$ reflecting the
additional chiral OPE structure beyond the Lie bracket alone.
\end{remark}

\begin{computation}[Casimir decomposition of \texorpdfstring{$\mathfrak{sl}_3^{\otimes n}$}{sl_3tensor n};
\ClaimStatusProvedHere]
\label{comp:sl3-casimir-decomp}
\index{Casimir operator!eigenspace decomposition}
\index{sl3@$\mathfrak{sl}_3$!Casimir decomposition}

The diagonal adjoint action of $\mathfrak{sl}_3$ on tensor powers
$\mathfrak{g}^{\otimes n}$ decomposes into irreducible representations
$V(\lambda_1, \lambda_2)$, identified by the scaled quadratic Casimir
$3C_2(V(\lambda_1,\lambda_2)) = 2(\lambda_1^2 + \lambda_2^2 + \lambda_1\lambda_2) + 6(\lambda_1 + \lambda_2)$.
Direct computation
gives the eigenspace multiplicities:
\begin{center}
\begin{tabular}{r|ccccccccc|r}
$3C_2$ & 0 & 18 & 36 & 48 & 72 & 90 & 108 & 120 & 144 & total \\
\hline
$n=1$ & & 8 & & & & & & & & 8 \\
$n=2$ & 1 & 16 & 20 & 27 & & & & & & 64 \\
$n=3$ & 2 & 64 & 80 & 162 & 140 & 64 & & & & 512 \\
$n=4$ & 8 & 256 & 400 & 891 & 1050 & 768 & 112 & 486 & 125 & 4096
\end{tabular}
\end{center}
The $n=2$ row matches the known decomposition
$\operatorname{adj} \otimes \operatorname{adj}
= \mathbf{1} \oplus \mathbf{8} \oplus \mathbf{8}
\oplus \mathbf{10} \oplus \overline{\mathbf{10}} \oplus \mathbf{27}$.
The $n=4$ row (dimension~$4096$) is the representation-theoretic
input for the degree-$4$ PBW analysis.

The Chevalley--Eilenberg bracket differential
$d\colon \mathfrak{g}^{\otimes n} \to \mathfrak{g}^{\otimes (n-1)}$
preserves Casimir eigenspaces by equivariance. Its rank within each eigenspace:
\begin{center}
\begin{tabular}{r|cccccc|r}
$3C_2$ & 0 & 18 & 36 & 48 & 72 & 90 & total \\
\hline
$\operatorname{rank}(d_2|_\lambda)$ & 0 & 8 & 0 & 0 & -- & -- & 8 \\
$\operatorname{rank}(d_3|_\lambda)$ & 1 & 8 & 20 & 27 & 0 & 0 & 56 \\
$\operatorname{rank}(d_4|_\lambda)$ & 1 & 56 & 60 & 135 & 140 & 64 & 456 \\
\end{tabular}
\end{center}
\emph{The CE complex on the full tensor algebra
$T(\mathfrak{g}) = \bigoplus_n \mathfrak{g}^{\otimes n}$ is exact}:
$\operatorname{rank}(d_{n+1}) = \dim \ker(d_n)$ for all $n \geq 2$
(through $n = 5$: ranks $8, 56, 456, 3640$).
The exactness is consistent with the acyclicity of the bar complex of $U(\mathfrak{g})$,
distinct from Chevalley--Eilenberg cohomology on the exterior algebra
$\Lambda^*(\mathfrak{g}^*)$ (where $H^3(\mathfrak{sl}_3) = \mathbb{C}$).

The eigenvalue-$0$ multiplicities (trivial summands in $\mathfrak{g}^{\otimes n}$)
are $0, 1, 2, 8$ for $n = 1, 2, 3, 4$. The $n=2$ invariant is the
Killing form~$C_2$; the two $n=3$ invariants are the antisymmetric
structure constants $f_{abc}$ and the symmetric cubic tensor $d_{abc}$.
The $8$ quartic invariant tensors include products $\kappa_{ab}\kappa_{cd}$
(three pairings), contractions $f_{abe}f_{cde}$ and $d_{abe}d_{cde}$,
and mixed contractions.
\end{computation}

\begin{remark}[No second Hilbert inversion for the \texorpdfstring{$\mathfrak{sl}_3$}{sl3} dual]
\label{rem:sl3-koszul-dual}
\index{Koszul dual!sl3@$\mathfrak{sl}_3$}

The conditional sequence
\[
1,\;8,\;36,\;204,\;1352,\;9892,\ldots
\]
is already the Hilbert sequence of the dual coalgebra
$A^i=H^\star\barB(\widehat{\mathfrak{sl}}_{3,k})$ and, after
graded linear/Verdier duality, of the Koszul-dual algebra $A^!$.
One must not apply the classical inversion
$H_A(t)H_{A^!}(-t)=1$ to this sequence a second time. That inversion
starts from the Hilbert series of the original algebra $A$ and, under
Koszulness, solves for the dual Hilbert series. Applying it to an
already dual series produces the unrelated numbers
$1,8,28,140,\ldots$ and confuses $A$, $B(A)$, $A^i$, and $A^!$.

The value $1352$ in degree~$4$ remains conditional until an
independent degree-$4$ bar computation or a proof of the recurrence
in Conjecture~\ref{conj:sl3-bar-gf} is supplied.
\end{remark}

\begin{computation}[Quadratic relation scan for \texorpdfstring{$\widehat{\mathfrak{sl}}_{3,k}$}{sl-hat_3,k};
\ClaimStatusProvedHere]
\label{comp:sl3-koszul-dual-scan}
\index{Koszul dual!quadratic relation scan}
\index{sl3@$\mathfrak{sl}_3$!Koszul dual scan}

For a quadratic algebra $A = T(V)/(R)$ with $R \subseteq V \otimes V$,
the Koszul dual is $A^! = T(V^*) / (R^\perp)$ and $\dim(A^!)_n$
is computed from the degree-$n$ ideal generated by $R^\perp$ in
$T(V^*)$. The scan tests several natural choices of $R$ for
$V = \mathfrak{g} = \mathfrak{sl}_3$ (dimension~$8$) and compares
the resulting Koszul dual dimensions with the conditional
Koszul-dual Hilbert sequence
$[1, 8, 36, 204, 1352, \ldots]$.

\medskip
\noindent
\begin{center}
\begin{tabular}{l|c|cccc}
Relation space $R$ & $\dim R$ & $(A^!)_0$ & $(A^!)_1$ & $(A^!)_2$ & $(A^!)_3$ \\
\hline
$\Lambda^2(V)$ (classical) & 28 & 1 & 8 & 28 & 56 \\
$\Lambda^2(V) + \langle \kappa \rangle$ & 29 & 1 & 8 & 29 & 64 \\
$\Lambda^2(V) \oplus \operatorname{diag}$ & 36 & 1 & 8 & \textbf{36} & 120 \\
$V \otimes V$ (full) & 64 & 1 & 8 & 64 & 512 \\
\hline
\textit{Dual Hilbert target} & n/a & 1 & 8 & \textit{36} & \textit{204}
\end{tabular}
\end{center}

Here $\langle \kappa \rangle$ is the Killing form element
$\sum_{a,b} \kappa^{ab} e_a \otimes e_b \in S^2(V)$, and
``$\operatorname{diag}$'' denotes $\operatorname{span}\{e_a \otimes e_a
: a = 1, \ldots, 8\}$.

\medskip
\noindent\textbf{Finding.}
The relation space $R = \Lambda^2(V) \oplus \operatorname{diag}$ (dimension~$36$)
is the unique simple choice that reproduces the degree-$2$ target
$\dim(A^!)_2 = 36$. However, it gives $(A^!)_3 = 120 \neq 204$.
The gap
\[
204-120=84
\]
is the measured degree-$3$ failure of the closest quadratic
presentation. It is not obtained by re-inverting the dual Hilbert
series; the numbers $1,8,28,140,\ldots$ belong to that
incorrect second inversion and are not dual Hilbert dimensions here.

No simple quadratic relation space $R \subseteq V \otimes V$ reproduces
the Koszul-dual Hilbert target beyond degree~$2$. The chiral bar complex
involves the complete OPE-pole data in Borcherds' identity, including
the quadratic bracket-and-Killing data. This is the algebraic shadow
of the $R^{(1)}$ barrier (higher residues on the configuration space)
that prevents the bracket differential alone from satisfying $d^2 = 0$.

\emph{The script \texttt{koszul\_dual\_dims.py} computes the table.}
\end{computation}

\begin{remark}[Euler characteristic obstruction to naive bar models]
\label{rem:euler-chi-obstruction}
\index{bar complex!Euler characteristic obstruction}

For any cochain complex $(C^\bullet, d)$ with $d^2 = 0$ and
finite-dimensional weight spaces, the Euler characteristic
$\sum_n (-1)^n \dim C^n_H = \sum_n (-1)^n \dim H^n_H$ holds
at each weight~$H$. This is a necessary condition for
the chain dimensions and cohomology to be mutually consistent.

\medskip
\noindent\emph{Model~A} (bracket complex): $B^n = \mathfrak{g}^{\otimes n}
\otimes \operatorname{OS}^{n-1}(\operatorname{Conf}_n)$, with
$\dim \mathfrak{g}_H = 8$ at each conformal weight $H \geq 1$.
At weight~$2$: $\chi_2 = -8 + 64 = 56$, requiring
$H^2_2 \geq 56 > 36 = H^2$. \emph{Inconsistent}.
This is the correct outcome: $d_{\mathrm{bracket}}^2 \neq 0$ on this
complex.

\medskip
\noindent\emph{Model~B} (full algebra): $B^n = \bar{A}^{\otimes n}
\otimes \operatorname{OS}^{n-1}$, with $\dim \bar{A}_H = \dim V_H$
(vacuum module dimensions $1, 8, 44, 192, \ldots$).
At weight~$2$: $\chi_2 = -44 + 64 = 20$, consistent.
At weight~$3$: $\chi_3 = -192 + 704 - 1024 = -512$, requiring
$H^3_3 \geq 512$, which is incompatible with the degree-$3$ target
$\dim(A^!)_3=204$. \emph{Inconsistent}.

\medskip
Neither model reproduces the conditional Koszul-dual Hilbert sequence
$[1, 8, 36, 204, 1352, \ldots]$. The actual chiral bar
complex involves the full Borcherds product (including higher
residues), and its chain structure differs from both naive
models. This is a \emph{chain-level} manifestation of the
$R^{(1)}$ barrier: the curvature corrections modify the differential
and the effective chain space simultaneously.
\end{remark}


\section{\texorpdfstring{Non-simply-laced example:
$\widehat{\mathfrak{so}}_5$ (type $B_2$)}{Non-simply-laced example:
so-5 (type B-2)}}
\label{sec:B2-details}
\index{non-simply-laced!bar complex}
\index{so5-hat@$\widehat{\mathfrak{so}}_5$!bar complex}
\index{B2@$B_2$|see{$\widehat{\mathfrak{so}}_5$}}

All Kac--Moody computations in the manuscript so far treat
simply-laced types ($A_n$, $D_n$, $E_n$), where roots have a single length
and $h^\vee = \mathsf{h}$ (Coxeter number equals dual Coxeter number).
The next computation treats the bar complex for
$\widehat{\mathfrak{so}}_{5,k}$ (type $B_2$),
the simplest non-simply-laced affine Kac--Moody algebra. The key
difference: \emph{the dual Coxeter number $h^\vee = 3$ is not equal
to the Coxeter number $\mathsf{h} = 4$}, and the level-shifting duality
$k \mapsto -k - 2h^\vee = -k - 6$ uses $h^\vee$, not $h$.

\subsection{Root system and structure constants}

\begin{computation}[\texorpdfstring{$B_2$}{B_2} root system]The Lie algebra $\mathfrak{so}_5 \cong \mathfrak{sp}_4$ has:
\begin{itemize}
\item Rank: $2$
\item Dimension: $\dim(\mathfrak{so}_5) = 10$
\item Simple roots: $\alpha_1$ (long), $\alpha_2$ (short), with
$|\alpha_1|^2 = 2$, $|\alpha_2|^2 = 1$
\item Positive roots: $\Delta_+ = \{\alpha_1,\; \alpha_2,\;
\alpha_1 + \alpha_2,\; \alpha_1 + 2\alpha_2\}$
\item Cartan matrix:
$A = \begin{pmatrix} 2 & -1 \\ -2 & 2 \end{pmatrix}$
\item Dual Coxeter number: $h^\vee = 3$
\item Coxeter number: $\mathsf{h} = 4$
\item Lacing number: $r^\vee = 2$ (ratio of long to short root squared lengths)
\end{itemize}
\end{computation}

\begin{computation}[\texorpdfstring{$\widehat{\mathfrak{so}}_{5,k}$}{so_5,k} OPE]\index{OPE!so5@$\widehat{\mathfrak{so}}_5$}
Generators: Cartan currents $H_1(z), H_2(z)$ and root currents
$E_{\pm\alpha}(z)$ for $\alpha \in \Delta_+$. Basis:
$\{H_1, H_2, E_1, E_2, E_{12}, E_{122}, F_1, F_2, F_{12}, F_{122}\}$
where $E_{12} = [E_1, E_2]$, $E_{122} = [E_{12}, E_2] = [[E_1, E_2], E_2]$,
and similarly for $F$.

\emph{Cartan--Cartan OPE.}
\begin{align*}
H_1(z) H_1(w) &\sim \frac{2k}{(z-w)^2}, &
H_1(z) H_2(w) &\sim \frac{-k}{(z-w)^2}, \\
H_2(z) H_2(w) &\sim \frac{k}{(z-w)^2} &
\end{align*}
(The Cartan matrix for the \emph{invariant form} is
$(\alpha_i^\vee, \alpha_j^\vee)_{\mathrm{inv}} = d_i A_{ij}$ where
$d_1 = 1$, $d_2 = 1/2$ for $B_2$ with long-root normalization
$|\alpha_1|^2 = 2$. The OPE uses the normalized form
$k \cdot (\cdot, \cdot)_{\mathrm{std}}$.)

\emph{Root--root OPE (sample).}
\begin{align*}
E_1(z) F_1(w) &\sim \frac{k}{(z-w)^2} + \frac{H_1(w)}{z-w} \\
E_2(z) F_2(w) &\sim \frac{k/2}{(z-w)^2} + \frac{H_2(w)}{z-w}
\end{align*}
The coefficient $k/2$ for the short root $\alpha_2$ (versus $k$ for the
long root $\alpha_1$) is the crucial non-simply-laced feature: the
double-pole coefficient is $k \cdot |\alpha|^2/2$, which depends on
the root length.

\emph{Serre relations.}
The Serre relations for $B_2$ are:
\begin{align}
(\operatorname{ad} E_1)^2(E_2) &= 0
\quad\text{(same as $A_2$: long root)} \label{eq:serre-B2-1} \\
(\operatorname{ad} E_2)^3(E_1) &= 0
\quad\text{(new: short root, order 3)} \label{eq:serre-B2-2}
\end{align}
The asymmetry $(\operatorname{ad} E_1)^2$ vs.\
$(\operatorname{ad} E_2)^3$ reflects the off-diagonal Cartan entries
$A_{12} = -1$, $A_{21} = -2$.
\end{computation}

\subsection{Bar differential}

\begin{computation}[Degree-2 bar differential for \texorpdfstring{$\widehat{\mathfrak{so}}_{5,k}$}{so_5,k}]
\index{bar complex!so5@$\widehat{\mathfrak{so}}_5$!degree 2}

\emph{Dimensions.}
$\bar{B}^1 = \mathfrak{so}_5 \cong \mathbb{C}^{10}$ and
$\bar{B}^2 = \mathfrak{so}_5^{\otimes 2} \otimes \Omega^1(\overline{\mathrm{Conf}}_2)
= \mathbb{C}^{100}$ (since $\dim \Omega^1 = 1$).

\emph{Differential $d\colon \bar{B}^2 \to \bar{B}^1$.}
On constant form:
\begin{align}
d[H_{1,-m}|E_{1,-n}] &= A_{11}\,[E_{1,-(m+n)}] = 2[E_{1,-(m+n)}]
\\
d[H_{1,-m}|E_{2,-n}] &= A_{12}\,[E_{2,-(m+n)}] = -[E_{2,-(m+n)}]
\\
d[H_{2,-m}|E_{1,-n}] &= A_{21}\,[E_{1,-(m+n)}] = -2[E_{1,-(m+n)}]
\\
d[H_{2,-m}|E_{2,-n}] &= A_{22}\,[E_{2,-(m+n)}] = 2[E_{2,-(m+n)}]
\label{eq:so5-H2E2}
\end{align}

\emph{Level contributions (curvature).}
\begin{align}
d[E_{1,-m}|F_{1,-n}] &= [H_{1,-(m+n)}] + k \cdot m\,\delta_{m+n,0}\,[1]
\label{eq:so5-EF-long} \\
d[E_{2,-m}|F_{2,-n}] &= [H_{2,-(m+n)}] + \tfrac{k}{2}\cdot m\,\delta_{m+n,0}\,[1]
\label{eq:so5-EF-short}
\end{align}
The factor $k/2$ in~\eqref{eq:so5-EF-short} is the short-root
double-pole coefficient. This asymmetry between
\eqref{eq:so5-EF-long} and~\eqref{eq:so5-EF-short} is the
fingerprint of the non-simply-laced structure in the bar complex.

\emph{Non-simple root contributions.}
\begin{align}
d[E_{1,-m}|E_{2,-n}] &= [E_{12,-(m+n)}]
\quad\text{(since $[E_1, E_2] = E_{12}$)} \\
d[E_{12,-m}|E_{2,-n}] &= [E_{122,-(m+n)}]
\quad\text{(since $[E_{12}, E_2] = E_{122}$)} \\
d[E_1|E_{122}] &= 0
\quad\text{(since $[E_1, E_{122}] = 0$ in $\mathfrak{so}_5$)}
\end{align}
\end{computation}

\begin{computation}[Curvature and level-shifting for \texorpdfstring{$\widehat{\mathfrak{so}}_{5,k}$}{so_5,k}]
\label{comp:so5-curvature}
\index{curvature!so5@$\widehat{\mathfrak{so}}_5$}
\index{non-simply-laced!curvature}

The curvature is computed from $m_1^2$:
\begin{equation}
m_0 = \frac{k + h^\vee}{2h^\vee} \cdot \kappa
= \frac{k + 3}{6} \cdot \kappa_{\mathfrak{so}_5}
\end{equation}
where $h^\vee = 3$, while the Coxeter number is $h = 4$. This
vanishes at the critical level
$k = -h^\vee = -3$.

\emph{Affine same-family shadow.}
\begin{equation}\label{eq:so5-level-shift}
k'=-k-2h^\vee=-k-6.
\end{equation}
The shifted affine algebra $\widehat{\mathfrak{so}}_{5,k'}$ records
the same scalar dual lane; it differs from both the raw bar coalgebra
and the derived chiral centre. The shift uses the dual Coxeter number
$h^\vee = 3$, while the Coxeter number is $h = 4$.

\emph{Central charge.}
\begin{equation}
c(\widehat{\mathfrak{so}}_{5,k}) = \frac{k \cdot \dim(\mathfrak{so}_5)}{k + h^\vee}
= \frac{10k}{k + 3}
\end{equation}
The complementarity gives $c + c' = 2\dim(\mathfrak{so}_5) = 20$:
\begin{equation}
c(\widehat{\mathfrak{so}}_{5,k}) + c(\widehat{\mathfrak{so}}_{5,-k-6})
= \frac{10k}{k+3} + \frac{10(-k-6)}{-k-6+3}
= \frac{10k}{k+3} + \frac{-10(k+6)}{-(k+3)}
= \frac{10k + 10(k+6)}{k+3} = 20
\end{equation}
as predicted by Theorem~\ref{thm:central-charge-complementarity}.

\emph{Non-simply-laced subtlety: the asymmetric double pole.}
In a simply-laced algebra, the double-pole coefficient in
$E_\alpha(z) F_\alpha(w) \sim k\delta_{\alpha\beta}/(z-w)^2$
is the same for all roots. For $B_2$:
\begin{center}
\renewcommand{\arraystretch}{1.2}
\begin{tabular}{lcc}
\textbf{Root $\alpha$} & \textbf{$|\alpha|^2$} & \textbf{Double-pole coeff.} \\
\hline
$\alpha_1$ (long) & $2$ & $k$ \\
$\alpha_2$ (short) & $1$ & $k/2$ \\
$\alpha_1+\alpha_2$ (short) & $1$ & $k/2$ \\
$\alpha_1+2\alpha_2$ (long) & $2$ & $k$
\end{tabular}
\end{center}
All long roots contribute $k$; the two short positive roots
$\alpha_2$ and $\alpha_1+\alpha_2$ contribute $k/2$. The curvature computation uses the adjoint
Casimir $C_2^{\mathrm{ad}} = 2h^\vee = 6$ (independent of root lengths),
so the curvature formula $m_0 = (k+3)\kappa/6$ is uniform.
\end{computation}

\begin{computation}[Serre relations for \texorpdfstring{$B_2$}{B_2}: bar complex encoding]
\label{comp:B2-serre}
\index{Serre relations!B2@$B_2$}

The two Serre relations~\eqref{eq:serre-B2-1}--\eqref{eq:serre-B2-2}
encode differently in the bar complex:

\emph{(a) Long-root Serre: $(\operatorname{ad} E_1)^2(E_2) = 0$.}
This is identical in structure to the $A_2$ Serre relation
(Computation~\ref{comp:sl3-degree3-complete}). The bar element is:
\[
S_1 = [E_1|E_1|E_2] - 2[E_1|E_2|E_1] + [E_2|E_1|E_1]
\in \bar{B}^3
\]
and $d(S_1) = 0$ in degree~2: each term $d[E_i|E_j|E_k]$ produces a sum of degree-$2$ elements $[E_i \cdot E_j | E_k] \pm [E_i | E_j \cdot E_k]$, and the alternating coefficients cancel pairwise by the Jacobi identity for $[E_1, [E_1, E_2]]$, exactly as in the $\mathfrak{sl}_3$ computation (Proposition~\ref{prop:sl3-serre-cohomology}).

\emph{(b) Short-root Serre: $(\operatorname{ad} E_2)^3(E_1) = 0$.}
This involves three applications of $\operatorname{ad} E_2$
(versus two for the long root), reflecting $A_{21} = -2$. The bar element
is of degree~4:
\begin{align*}
S_2 &= [E_2|E_2|E_2|E_1]
- 3[E_2|E_2|E_1|E_2]
+ 3[E_2|E_1|E_2|E_2]
- [E_1|E_2|E_2|E_2]
\end{align*}
(the binomial coefficients $\binom{3}{0}, -\binom{3}{1}, \binom{3}{2}, -\binom{3}{3}$ from the Koszul sign rule for the free Lie element $(\operatorname{ad} E_2)^3(E_1)$).

The bar differential $d\colon \bar{B}^4 \to \bar{B}^3$ acts on $S_2$
via residues at six collision divisors $D_{ij}$ ($1 \leq i < j \leq 4$)
on $\overline{\mathrm{Conf}}_4$. Each collision between two copies of $E_2$
gives $[E_2, E_2] = 0$; each collision $E_2 \cdot E_1$ gives $E_{12}$;
each collision $E_{12} \cdot E_2$ gives $E_{122}$. The alternating
signs in $S_2$ ensure all terms cancel:
\begin{equation}
d(S_2 \otimes \omega_4) = 0
\end{equation}
for any $\omega_4 \in \Omega^3(\overline{\mathrm{Conf}}_4)$.

\emph{Structural comparison.}
\begin{center}
\renewcommand{\arraystretch}{1.2}
\begin{tabular}{lcccc}
\textbf{Serre} & \textbf{Root} & \textbf{Bar degree} & \textbf{$A_{ij}$} &
\textbf{Coefficients} \\
\hline
$(\operatorname{ad} E_1)^2(E_2)$ & long & $3$ & $-1$ & $1, -2, 1$ \\
$(\operatorname{ad} E_2)^3(E_1)$ & short & $4$ & $-2$ & $1, -3, 3, -1$
\end{tabular}
\end{center}
In general, for a non-simply-laced algebra with Cartan entry $A_{ji} = -m$
(for the multiple edge this is the local lacing multiplicity), the Serre relation
$(\operatorname{ad} E_j)^{m+1}(E_i) = 0$ produces a bar cycle
of degree $m+2$ with binomial coefficients $\binom{m+1}{s}$.
\end{computation}

\begin{proposition}[Bar complex dimensions for \texorpdfstring{$\widehat{\mathfrak{so}}_{5,k}$}{so_5,k};
\ClaimStatusProvedHere]
\label{prop:so5-bar-dims}
\index{bar complex!dimension table!so5@$\mathfrak{so}_5$}
The dimensions of the bar complex through degree~4:
\begin{center}
\renewcommand{\arraystretch}{1.2}
\begin{tabular}{c|ccc|c}
$n$ & $\dim(\mathfrak{so}_5^{\otimes n})$ &
$\dim \Omega^{n-1}$ & $\dim \bar{B}^n$ & Serre contributions \\
\hline
$0$ & $1$ & $1$ & $1$ & vacuum \\
$1$ & $10$ & $1$ & $10$ & generators \\
$2$ & $100$ & $1$ & $100$ & OPE data (asymmetric poles) \\
$3$ & $1000$ & $2$ & $2000$ & long-root Serre ($\times 1$) \\
$4$ & $10000$ & $6$ & $60000$ & short-root Serre ($\times 1$)
\end{tabular}
\end{center}
The Serre contribution at degree 3 is one-dimensional (only
$(\operatorname{ad} E_1)^2(E_2) = 0$); at degree~4, the short-root
Serre $(\operatorname{ad} E_2)^3(E_1) = 0$ appears. This is in
contrast to $\mathfrak{sl}_3$ (type $A_2$), where both Serre relations
appear at degree~3.
\end{proposition}

\begin{proof}
The generator and form dimensions are $10^n$ and $(n-1)!$ respectively.
The Serre contributions are identified by the computation of
$\ker(d^n)/\operatorname{im}(d^{n+1})$ on the associated graded (where
the bar differential reduces to the Chevalley--Eilenberg differential
of $\mathfrak{so}_5$): the Serre relation
$(\operatorname{ad} E_j)^{1-A_{ji}}(E_i) = 0$ appears in the
$(2 - A_{ji})$-fold tensor power, giving degree $3$ for $A_{12} = -1$
and degree $4$ for $A_{21} = -2$.
Computation~\ref{comp:B2-serre} verifies both explicitly.
\end{proof}


\section{Geometric-algebraic bar comparison}
\label{sec:geometric-algebraic-comparison}
\index{bar complex!geometric vs algebraic}
\index{geometric bar complex!comparison with algebraic}

A central claim of this monograph is that the \emph{geometric} bar
complex (constructed from configuration space integrals on
Fulton--MacPherson compactifications) agrees with the \emph{algebraic}
bar complex (constructed from OPE residues). This agreement is
proved abstractly as
Theorem~\ref{thm:geometric-equals-operadic-bar}
(bar\_cobar\_construction.tex), but the term-by-term residue comparison has
not been exhibited. The comparison is carried out below for
$\widehat{\mathfrak{sl}}_{2,k}$, the simplest algebra where \emph{both}
curvature and bracket data are visible.

\subsection{Setup: two bar complexes}

\begin{computation}[Geometric vs.\ algebraic bar: degree~2 for \texorpdfstring{$\widehat{\mathfrak{sl}}_{2,k}$}{sl-hat_2,k}]
\label{comp:geom-alg-comparison-deg2}
\index{bar complex!geometric!degree 2}

\emph{Algebraic bar complex.}
The degree-2 algebraic bar differential is determined by the OPE:
\begin{equation}\label{eq:alg-bar-def}
d_{\mathrm{alg}}([J^a|J^b])
= \sum_{n \geq 0} J^a_{(n)} J^b
= k\kappa^{ab}\cdot\mathbf{1} + {f^{ab}}_c J^c
\end{equation}
where $J^a_{(1)} J^b = k\kappa^{ab}\cdot\mathbf{1}$ (double pole,
curvature) and $J^a_{(0)} J^b = {f^{ab}}_c J^c$ (simple pole,
bracket). Higher $n$-products vanish for the Kac--Moody OPE.

\emph{Geometric bar complex.}
The degree-2 element $[J^a|J^b] \otimes \eta_{12}$, where
$\eta_{12} = d\log(z_1 - z_2)$, lives in:
\[
\bar{B}^2_{\mathrm{geom}} = \mathfrak{g}^* \otimes \mathfrak{g}^*
\otimes H^0\bigl(\overline{C}_3(\mathbb{P}^1),\,
\Omega^1_{\log}\bigr).
\]
The geometric bar differential is
$d_{\mathrm{geom}} = d_{\mathrm{res}} + d_{\mathrm{dR}}$
(Definition~\ref{def:geometric-bar-definition}), where $d_{\mathrm{res}}$
extracts residues at the boundary divisor $D_{12} = \{z_1 = z_2\}
\subset \overline{C}_3(\mathbb{P}^1)$.

\emph{Collision extraction.}
Write $z_{12} = z_1 - z_2$ as local coordinate near $D_{12}$.
The two-point function on $\overline{C}_3(\mathbb{P}^1)$ is:
\begin{equation}
J^a(z_1)\, J^b(z_2) \cdot \eta_{12}
= \left(\frac{k\kappa^{ab}}{z_{12}^2}
+ \frac{{f^{ab}}_c J^c(z_2)}{z_{12}} + \text{reg}\right)
\cdot \frac{dz_{12}}{z_{12}}
\end{equation}
The geometric bar map uses the chiral-product extraction rule
\eqref{eq:frame-borcherds-extraction}, not the ordinary one-variable
Poincar\'e residue of the displayed meromorphic one-form. The
logarithmic form marks the boundary divisor; the
$\mathcal D$-module pushforward extracts all singular OPE coefficients:
\begin{align}
\partial_{D_{12}}\bigl([J^a|J^b]\otimes\eta_{12}\bigr)
&=
J^a_{(1)}J^b + J^a_{(0)}J^b \notag \\
&=
k\kappa^{ab}\cdot\mathbf{1}+{f^{ab}}_cJ^c.
\label{eq:geom-residue-simple}
\end{align}
The double pole contributes the curvature/vacuum component in
$\bar B^0$; the simple pole contributes the bracket component in
$\bar B^1$. The ordinary-residue calculation
$\operatorname{Res}_{z_{12}=0}(k\kappa^{ab}z_{12}^{-3}\,dz_{12})=0$
computes an ordinary meromorphic residue and discards the chiral
product coefficient. The frame computation
\eqref{eq:frame-borcherds-extraction} is the convention used by the
bar complex.

\emph{Scalar projection.}
For the uniform-weight scalar lane, the trace-form double-pole
projection contributes $\dim(\fg)k/(2h^\vee)$, while the simple-pole
Lie bracket channel contributes $\dim(\fg)/2$; together
\[
\kappa(V_k(\fg))=\frac{\dim(\fg)(k+h^\vee)}{2h^\vee}.
\]
This is the scalar modular characteristic. It is not an assertion that
the full tensor curvature is the scalar multiple
$\kappa(V_k(\fg))\mathbf 1$ in every channel.

\emph{Comparison.}
\begin{center}
\renewcommand{\arraystretch}{1.2}
\begin{tabular}{lcc}
\textbf{Term} & \textbf{Algebraic} & \textbf{Geometric} \\
\hline
Bracket & ${f^{ab}}_c J^c$ (simple pole) &
 $\partial_{D_{12}}[\cdots]_{\bar B^1} = {f^{ab}}_c J^c$ \\
Curvature & $k\kappa^{ab}$ (double pole) &
 $\partial_{D_{12}}[\cdots]_{\bar B^0}=k\kappa^{ab}$ \\
Higher & $0$ (no triple+ poles in KM) &
 $0$ (no higher residues)
\end{tabular}
\end{center}
The bracket and curvature channels match exactly before scalar
projection. The genus-$1$ computation of
Theorem~\ref{thm:sl2-genus1-curvature} propagates their scalar trace to
$\kappa(V_k(\mathfrak{sl}_2))=3(k+2)/4$.
\end{computation}

\subsection{Degree 3: Arnold relations and the Jacobi identity}

\begin{computation}[Geometric bar: degree~3 for \texorpdfstring{$\widehat{\mathfrak{sl}}_{2,k}$}{sl-hat_2,k}]
\label{comp:geom-alg-comparison-deg3}
\index{Arnold relations!Jacobi identity}
\index{Jacobi identity!from Arnold relations}

The degree-3 bar element
$[J^a|J^b|J^c] \otimes \omega_3$
lives in $\mathfrak{g}^{\otimes 3} \otimes
\Omega^2(\overline{C}_4(\mathbb{P}^1))$.

\emph{Arnold relation.}
On $C_3(\mathbb{C})$, the forms $\eta_{ij} = d\log(z_i - z_j)$
satisfy the Arnold relation
(Theorem~\ref{thm:arnold-relations}):
\begin{equation}\label{eq:arnold-deg3}
\eta_{12} \wedge \eta_{23} + \eta_{23} \wedge \eta_{31}
+ \eta_{31} \wedge \eta_{12} = 0
\end{equation}

\emph{Algebraic translation.}
The three terms in~\eqref{eq:arnold-deg3} correspond to the three
cyclic orderings of the collision sequence
$(z_1 = z_2, z_2 = z_3)$, $(z_2 = z_3, z_3 = z_1)$,
$(z_3 = z_1, z_1 = z_2)$.
The residue at each collision extracts a bracket:
\begin{align}
\operatorname{Res}_{D_{12}}\operatorname{Res}_{D_{23}}
[\cdots \otimes \eta_{12} \wedge \eta_{23}]
&\leftrightarrow [J^a, J^b]_{(0)} J^c
= {f^{ab}}_d\, J^d_{(0)} J^c
= {f^{ab}}_d\, {f^{dc}}_e J^e \\
\operatorname{Res}_{D_{23}}\operatorname{Res}_{D_{31}}
[\cdots \otimes \eta_{23} \wedge \eta_{31}]
&\leftrightarrow {f^{bc}}_d\, {f^{da}}_e J^e \\
\operatorname{Res}_{D_{31}}\operatorname{Res}_{D_{12}}
[\cdots \otimes \eta_{31} \wedge \eta_{12}]
&\leftrightarrow {f^{ca}}_d\, {f^{db}}_e J^e
\end{align}
The Arnold relation~\eqref{eq:arnold-deg3} forces:
\begin{equation}\label{eq:jacobi-from-arnold}
{f^{ab}}_d\, {f^{dc}}_e + {f^{bc}}_d\, {f^{da}}_e
+ {f^{ca}}_d\, {f^{db}}_e = 0
\end{equation}
which is precisely the \textbf{Jacobi identity} for
$\mathfrak{sl}_2$ (and, more generally, for any Lie algebra).

\emph{Significance.}
This is the geometric content of $\dzero^2 = 0$ at degree~3:
the Arnold relations on
$\overline{C}_4(\mathbb{P}^1)$ \emph{are} the Jacobi identity
for the structure constants extracted by residues. The
implication is bidirectional:
\begin{enumerate}
\item \emph{Arnold $\Rightarrow$ Jacobi}: the bar differential
 squares to zero because the forms satisfy the Arnold relations.
\item \emph{Jacobi $\Rightarrow$ Arnold}: given any Lie algebra,
 the structure constants $f^{ab}_c$ define a bar differential
 for which $\dzero^2 = 0$ is equivalent to the Arnold relation via
 the residue dictionary.
\end{enumerate}
This is Theorem~\ref{thm:geometric-equals-operadic-bar} made
completely explicit.
\end{computation}

\subsection{\texorpdfstring{Explicit matrices for $\mathfrak{sl}_2$}{Explicit matrices for sl2}}

\begin{computation}[Degree-2 bar differential matrix for \texorpdfstring{$\widehat{\mathfrak{sl}}_{2,k}$}{sl-hat_2,k}]
\label{comp:sl2-bar-matrix}
\index{bar complex!sl2@$\widehat{\mathfrak{sl}}_2$!differential matrix}

Using the basis $\{e, f, h\}$ for $\mathfrak{sl}_2$ with
$[e,f] = h$, $[h,e] = 2e$, $[h,f] = -2f$, the degree-2
bar differential $d\colon \bar{B}^2 \to \bar{B}^1$ on constant forms
is the $3 \times 9$ matrix
\[
\text{rows } = \{e,f,h\}, \qquad
\text{columns } =
\{[e|e], [e|f], [e|h], [f|e], [f|f], [f|h], [h|e], [h|f], [h|h]\}.
\]
With this ordering, one has:
\begin{equation}\label{eq:sl2-bar-matrix}
d = \begin{pmatrix}
0 & 0 & -2 & 0 & 0 & 0 & 2 & 0 & 0 \\
0 & 0 & 0 & 0 & 0 & 2 & 0 & -2 & 0 \\
0 & 1 & 0 & -1 & 0 & 0 & 0 & 0 & 0
\end{pmatrix}
\end{equation}

The entries are the structure constants
$d([J^a|J^b]) = {f^{ab}}_c J^c$:
\begin{center}
\renewcommand{\arraystretch}{1.2}
\begin{tabular}{lcl}
$d([e|f])$ & $=$ & $h$ \\
$d([f|e])$ & $=$ & $-h$ \\
$d([h|e])$ & $=$ & $2e$ \\
$d([e|h])$ & $=$ & $-2e$ \\
$d([h|f])$ & $=$ & $-2f$ \\
$d([f|h])$ & $=$ & $2f$ \\
$d([e|e]) = d([f|f]) = d([h|h])$ & $=$ & $0$
\end{tabular}
\end{center}

The curvature terms $k\kappa^{ab}\cdot\mathbf{1}$
\textup{(}mapping to $\bar{B}^0 = \mathbb{C}$\textup{)} are:
\begin{align}
d([e|f]) &\ni k \cdot 1 \cdot \mathbf{1}, &
d([f|e]) &\ni k \cdot 1 \cdot \mathbf{1}, &
d([h|h]) &\ni 2k \cdot \mathbf{1}
\end{align}

(using $(e,f) = 1$, $(h,h) = 2$ in the standard normalization).

\emph{Rank computation.}
The matrix~\eqref{eq:sl2-bar-matrix} has rank~$3$
(six non-zero entries, and $e, f, h$ are all hit).
Therefore $\dim \ker(d\colon \bar{B}^2 \to \bar{B}^1) = 9 - 3 = 6$
on constant forms. These 6 kernel elements are:
\begin{enumerate}
\item $[e|e],\; [f|f],\; [h|h]$ (diagonal, $3$ elements)
\item $[e|f] + [f|e]$ (symmetric part, killed by antisymmetry of bracket)
\item $[e|h] + [h|e]$ and $[f|h] + [h|f]$
 (kernel of Cartan action: $d([e|h] + [h|e]) = -2e + 2e = 0$)
\end{enumerate}

\emph{Geometric interpretation.}
Each kernel element corresponds to a relation in the quotient
$R^\perp \subset V \otimes V$ (the orthogonal complement of the
relation space $R = \operatorname{span}\{d(\bar{B}^2 \to \bar{B}^1)\}$). The quadratic dual is
$\bar{B}^{\mathrm{ch}}(\widehat{\mathfrak{sl}}_{2,k})^i
= \operatorname{sub}(T^c(V), R)$, the maximal sub-coalgebra of
$T^c(V)$ cogenerated by~$R$. The Verdier--Koszul dual algebra is
obtained from this coalgebra by graded linear/Verdier duality. The
Feigin--Frenkel shifted level
$-k-4$ records the same scalar modular characteristic as the dual
lane; it is not obtained by declaring the raw bar coalgebra equal to
the affine algebra $\widehat{\mathfrak{sl}}_{2,-k-4}$
(cf.\ Theorem~\ref{thm:sl2-koszul-dual}).
\end{computation}

\begin{remark}[Vacuum module decomposition and PBW comparison for
\texorpdfstring{$\widehat{\mathfrak{sl}}_{2,k}$}{sl-hat_2,k}]\label{rem:sl2-vacuum-module-pbw}
The bar complex computations above can be compared with
the representation theory of the vacuum module~$V_k$.
The PBW basis of~$V_k$ gives weight-space dimensions
$\dim(V_k)_h = p(h,3)$, the number of $3$-colored partitions of~$h$
(matching $\prod_{n \ge 1}(1-q^n)^{-3}$).
Under the zero-mode adjoint action of~$\mathfrak{sl}_2$,
the weight-$2$ space decomposes as
\begin{equation}\label{eq:sl2-weight2-decomp}
(V_k)_2 \;\cong\; V_5 \oplus V_3 \oplus V_1
\qquad (\text{spins } 2, 1, 0)
\end{equation}
with Casimir eigenvalues $\{12, 4, 0\}$ respectively.
The $\mathfrak{g}$-invariant dimensions at weights $h = 0, 1,
\ldots, 9$ are
\[
1,\; 0,\; 1,\; 1,\; 3,\; 3,\; 8,\; 9,\; 19,\; 25,
\]
consistent with the adjoint-invariant counting
$\dim (V_k)_h^{\mathfrak{g}} = [q^h]\,\prod_{n \ge 1}
(1-q^n)^{-1}$ for~$\mathfrak{sl}_2$.
The Sugawara construction
$T = \frac{1}{2(k+2)}\sum_a {:}J^a J^a{:}$
acts as~$L_0 = h$ on each weight-$h$ subspace, matching the
conformal-weight grading of the bar complex.
\end{remark}

\begin{proposition}[PBW \texorpdfstring{$E_2$}{E_2} from vacuum module data;
\ClaimStatusProvedHere]\label{prop:pbw-e2-from-vacuum-module}
\index{PBW spectral sequence!E2 from vacuum module}
\index{MC1!algorithmic check}
Let $\mathfrak{g}$ be a simple Lie algebra and
$\widehat{\mathfrak{g}}_k$ the associated Kac--Moody algebra
with vacuum module~$V_k$.
The $E_2$ page of the PBW spectral sequence for
$\barB(\widehat{\mathfrak{g}}_k)$ is:
\begin{equation}\label{eq:pbw-e2-vacuum}
E_2^{p,q}
\;\cong\;
H^p\!\bigl(\mathfrak{g};\,
(V_k)_q^{\mathfrak{g}\text{-}\mathrm{triv}}\bigr),
\end{equation}
where $(V_k)_q^{\mathfrak{g}\text{-}\mathrm{triv}}$
denotes the direct sum of trivial $\mathfrak{g}$-isotypic
components of the weight-$q$ subspace of~$V_k$.
In particular, $E_2^{0,q} =
\bigl((V_k)_q\bigr)^{\mathfrak{g}}$, with dimensions
computable from the Casimir eigenvalue decomposition
of~$(V_k)_q$.

MC1 concentration ($E_2^{p,q} = 0$ for $q \neq 0$) therefore
reduces to a representation-theoretic computation within the
vacuum module: check that the $\mathfrak{g}$-invariants of each
weight space carry trivial Lie algebra cohomology.
\end{proposition}

\begin{proof}
The PBW filtration on
$\barB^n = \mathfrak{g}^{\otimes n} \otimes \mathrm{OS}^{n-1}$
is by conformal weight. The associated graded at weight~$q$ is
$(V_k)_q$ as a $\mathfrak{g}$-module (via the zero-mode adjoint
action), and the $E_1$ differential is the Chevalley--Eilenberg
differential $d_{\mathrm{CE}} \colon
\Lambda^p(\mathfrak{g}) \otimes (V_k)_q \to
\Lambda^{p+1}(\mathfrak{g}) \otimes (V_k)_q$.
Taking cohomology gives~\eqref{eq:pbw-e2-vacuum}.
By complete reducibility (Weyl's theorem), the $\mathfrak{g}$-module
$(V_k)_q$ decomposes into irreducibles, and
$H^*(\mathfrak{g}; W) = 0$ for every nontrivial irreducible~$W$
(Whitehead's lemma). Hence only the trivial isotypic component
survives on~$E_2$.
\end{proof}

\begin{computation}[Degree-3 bar differential and curvature for \texorpdfstring{$\widehat{\mathfrak{sl}}_{2,k}$}{sl-hat_2,k}]
\label{comp:sl2-bar-deg3-curvature}
\ClaimStatusProvedHere
\index{bar complex!sl2@$\widehat{\mathfrak{sl}}_2$!degree~3}
\index{curvature!associator}

The degree-3 \emph{associative} bar differential
$d\colon \bar{B}^3 \to \bar{B}^2$ for $\mathfrak{sl}_2$ and
its square are computed below; $d^2 \neq 0$, with $d^2$ given by
the associator.

On the associative bar, $d([a|b|c]) = [m(a,b)|c] - [a|m(b,c)]$
where $m(a,b) = [a,b]$ is the Lie bracket. The $9 \times 27$ matrix
has rank~$8$. Representative entries:
\begin{center}
\renewcommand{\arraystretch}{1.2}
\begin{tabular}{lcl}
$d([e|f|e])$ & $=$ & $[h|e] + [e|h]$ \\
$d([e|f|h])$ & $=$ & $-2[e|f] + [h|h]$ \\
$d([h|e|h])$ & $=$ & $2[e|h] + 2[h|e]$ \\
$d([h|e|f])$ & $=$ & $2[e|f] - [h|h]$ \\
$d([e|h|e])$ & $=$ & $-4[e|e]$ \\
$d([f|h|f])$ & $=$ & $4[f|f]$
\end{tabular}
\end{center}

\emph{Curvature: $d^2 \neq 0$.}
Since the Lie bracket is not associative, $d^2 \neq 0$:
\begin{equation}\label{eq:sl2-assoc-bar-d2}
d^2([a|b|c]) = [[a,b],c] - [a,[b,c]]
\end{equation}
which is the \textbf{associator} of the Lie bracket.
The 12 nonzero instances are:
\begin{center}
\renewcommand{\arraystretch}{1.2}
\begin{tabular}{lclclcl}
$d^2([e|e|f])$ & $=$ & $2e$ &\qquad&
 $d^2([f|e|e])$ & $=$ & $-2e$ \\
$d^2([e|h|h])$ & $=$ & $4e$ &&
 $d^2([h|h|e])$ & $=$ & $-4e$ \\
$d^2([e|f|h])$ & $=$ & $-2h$ &&
 $d^2([f|e|h])$ & $=$ & $-2h$ \\
$d^2([h|e|f])$ & $=$ & $2h$ &&
 $d^2([h|f|e])$ & $=$ & $2h$ \\
$d^2([e|f|f])$ & $=$ & $-2f$ &&
 $d^2([f|f|e])$ & $=$ & $2f$ \\
$d^2([f|h|h])$ & $=$ & $4f$ &&
 $d^2([h|h|f])$ & $=$ & $-4f$
\end{tabular}
\end{center}

\emph{Connection to the chiral bar.}
In the chiral bar complex
$\bar{B}^3 = \mathfrak{g}^{\otimes 3} \otimes \operatorname{OS}^2(C_3)$,
the differential sums over all pairwise collisions with
Orlik--Solomon form residues. The Arnold relation
(Computation~\ref{comp:geom-alg-comparison-deg3}) is
exactly the cancellation needed to convert the
associator~\eqref{eq:sl2-assoc-bar-d2} into the Jacobi identity,
yielding $d^2 = 0$ on the Chevalley--Eilenberg complex
$C^*(\mathfrak{g}; \mathbb{C}) = \Lambda^*(\mathfrak{g}^*)$
(cf.~Appendix~\ref{app:signs}).

This verifies the \emph{classical} CE differential on
$\Lambda^n(\mathfrak{g}^*)$: $d^2 = 0$ follows from the Jacobi identity,
and the cohomology matches $H^*(\mathfrak{g}; \mathbb{C})$.
The full chiral bar differential on
$\mathfrak{g}^{\otimes n} \otimes \operatorname{OS}^{n-1}(C_n)$
requires the Borcherds identity
(Proposition~\ref{prop:pole-decomposition}).
\end{computation}

\begin{computation}[SDR and formality for \texorpdfstring{$\mathfrak{sl}_2$}{sl_2}]
\label{comp:sl2-ce-sdr}
\ClaimStatusProvedHere
\index{SDR (strong deformation retract)!sl2@$\mathfrak{sl}_2$}
\index{formality!sl2@$\mathfrak{sl}_2$}

The Chevalley--Eilenberg complex
$C^*(\mathfrak{sl}_2; \mathbb{C})
= \Lambda^*(\mathfrak{sl}_2^*)$
has differentials:
\begin{equation}\label{eq:sl2-ce-diff}
d_0 = 0, \quad
d_1 = \operatorname{diag}(-2, 1, -2), \quad
d_2 = 0.
\end{equation}
One verifies $d^2 = 0$ and
$H^*(\mathfrak{sl}_2; \mathbb{C}) = \{1, 0, 0, 1\}
= \Lambda(x_3)$ with $|x_3| = 3$.

Since $d_1$ is an isomorphism (rank~3), the contracting homotopy
for the strong deformation retract
$(p, \iota, h)$ of
$C^*(\mathfrak{sl}_2)$
onto $H^*(\mathfrak{sl}_2)$
has the explicit form:
\begin{equation}\label{eq:sl2-sdr}
h_1 = 0, \quad
h_2 = d_1^{-1} = \operatorname{diag}(-\tfrac{1}{2}, 1, -\tfrac{1}{2}), \quad
h_3 = 0,
\end{equation}
with $\iota_k = p_k = \operatorname{id}$ for $k \in \{0, 3\}$
and $\iota_k = p_k = 0$ for $k \in \{1, 2\}$.

\emph{Formality.}
The homotopy transfer theorem applied to this SDR shows that
$\mathfrak{sl}_2$ is \textbf{formal}: all transferred $L_\infty$
operations $l_n$ vanish for $n \geq 3$, since the homotopy $h$
is concentrated in degree~2 and the cohomology gap
($H^1 = H^2 = 0$) prevents the tree formula from producing
nonzero contributions. This is consistent with the general
result that semisimple Lie algebras are formal as dg~Lie algebras
(cf.\ Deligne--Griffiths--Morgan--Sullivan~\cite{DGMS75}).

The SDR conditions
($dh + hd = \operatorname{id} - \iota \circ p$,
$h^2 = 0$, $p \circ h = 0$, $h \circ \iota = 0$)
are satisfied by direct matrix computation on the
$10$-dimensional complex $\Lambda^*(\mathfrak{sl}_2^*)$.
\end{computation}

\begin{computation}[CE cohomology of
 \texorpdfstring{$\mathfrak{sl}_2 \otimes t^{-1}\mathbb{C}{[t^{-1}]}$}{sl_2 tensor t^-1C[t^-1]}]
\label{comp:sl2-ce-verification}
\ClaimStatusProvedHere
\index{Chevalley--Eilenberg cohomology!loop algebra}
\index{bar complex!PBW spectral sequence!CE differential}

The loop algebra
$\mathfrak{g}_- = \mathfrak{sl}_2 \otimes t^{-1}\mathbb{C}[t^{-1}]$
has generators $(a, n)$ for $a \in \{e, h, f\}$ and $n \geq 1$,
with bracket
$[(a, m), (b, n)] = ([a, b], m + n)$
(no central extension for $m, n \geq 1$).
The CE complex $\Lambda^*(\mathfrak{g}_-^*)$,
graded by conformal weight $H = \sum n_i$,
has differential $d_{\mathrm{CE}}$ preserving~$H$.

\emph{$d^2 = 0$ check.}
Checked at all degrees $p \leq 4$ and weights $H \leq 5$.

\emph{Chain group dimensions.}
$\dim \Lambda^p_H$ at low weights:
\begin{center}
\renewcommand{\arraystretch}{1.1}
\begin{tabular}{c|ccccc}
$H \backslash p$ & 0 & 1 & 2 & 3 & 4 \\
\hline
1 & 0 & 3 & 0 & 0 & 0 \\
2 & 0 & 3 & 3 & 0 & 0 \\
3 & 0 & 3 & 9 & 1 & 0 \\
4 & 0 & 3 & 12 & 9 & 0
\end{tabular}
\end{center}

\emph{Differential ranks.}
$\mathrm{rk}(d^{1,2}_{\mathrm{CE}}) = 3$,
$\mathrm{rk}(d^{2,3}_{\mathrm{CE}}) = 1$,
$\mathrm{rk}(d^{1,3}_{\mathrm{CE}}) = 3$,
$\mathrm{rk}(d^{1,4}_{\mathrm{CE}}) = 3$.

\emph{Cohomology.}
$H^1_{\mathrm{CE}} = 3$ (all at weight~$1$),
$H^2_{\mathrm{CE}} = 5$ (all at weight~$3$:
$\ker = 9 - 1 = 8$, $\mathrm{im} = 3$, $H^2_3 = 5$),
$H^2_H = 0$ for $H \neq 3$ (checked through $H = 6$).
This is the $E_2$ page of the PBW spectral
sequence on $\barBgeom(\widehat{\mathfrak{sl}}_2)$
(Corollary~\ref{cor:bar-cohomology-koszul-dual}),
giving $\dim H^n(\barBgeom(\widehat{\mathfrak{sl}}_2))
= H^n_{\mathrm{CE}}(\mathfrak{g}_-, \mathbb{C})$.

\emph{Matrix source.}
The \texttt{LoopAlgebraCE} computation records the matrices used for
these ranks.
\end{computation}

\begin{remark}[The comparison theorem in action]
\label{rem:comparison-summary}
\index{bar complex!comparison theorem}
Computations~\ref{comp:geom-alg-comparison-deg2}--\ref{comp:sl2-bar-deg3-curvature}
realize Theorem~\ref{thm:geometric-equals-operadic-bar} for
$\widehat{\mathfrak{sl}}_{2,k}$ at degrees~2 and~3. The dictionary is:
\begin{center}
\renewcommand{\arraystretch}{1.3}
\begin{tabular}{lll}
\textbf{Geometric} & \textbf{Algebraic} & \textbf{Identity} \\
\hline
$\partial_{D_{ij}}$ &
 $n$-th product $a_{(n)}b$ & OPE extraction \\
Arnold relation & Jacobi identity & $d^2 = 0$ \\
$B$-cycle period & Curvature $m_0$ & Genus-1 contribution \\
FM boundary faces & Stasheff associahedron & $A_\infty$ structure \\
\end{tabular}
\end{center}
The comparison shows that the geometric language of
Parts~\ref{part:bar-complex}--\ref{part:standard-landscape} (configuration space integrals, FM
compactifications, Verdier duality) and the algebraic language
of Part~\ref{part:physics-bridges} (OPE residues, screening operators,
spectral sequences) are \emph{identical}: the same computation carried out in two
different coordinate systems.
At the scalar projection, the unifying object is the shadow obstruction tower
$\Theta_\cA^{\leq r}$
(Definition~\ref{def:shadow-postnikov-tower}), whose
scalar level $\Theta_\cA^{\leq 2} = \kappa(\cA)$
(Theorem~\ref{thm:explicit-theta})
controls the genus family:
geometrically, $\kappa$ is the curvature of
$\barB^{(g)}(\cA) \to \overline{\mathcal{M}}_g$;
algebraically, it is the Killing class
$\kappa \in H^2_{\mathrm{cyc}}(\fg,\fg)$ paired with
the Hodge class $\lambda_g$. The table above is the
dictionary for translating between the two coordinate
systems at the scalar level.
\end{remark}

\section{\texorpdfstring{Finite horizontal BGG spine inside
$\widehat{\mathfrak{sl}}_2$}{Finite horizontal BGG spine inside sl-2-hat}}
\label{sec:bgg-sl2-bar}
\index{BGG resolution!from bar complex!sl2@$\mathfrak{sl}_2$}
\index{bar complex!BGG resolution!sl2@$\mathfrak{sl}_2$}

The constant-current sector of the
$\widehat{\mathfrak{sl}}_{2,k}$ bar complex contains the finite
horizontal $\mathfrak{sl}_2$ BGG complex. The raw bar cohomology is
the dual coalgebra $A^i=H^\star B(A)$; it is not a Verma module.
Verma modules occur on the module bar spectral sequence of
Theorem~\ref{thm:bgg-from-bar}. The finite horizontal comparison
identifies the matrix entries of
Computation~\ref{comp:sl2-bar-matrix} with the singular vector
$f\cdot v_0\in M(0)$.
At genus $g \geq 1$, the scalar projection of the bar complex carries
curvature
$m_0 = \kappa(\widehat{\mathfrak{sl}}_{2,k}) \cdot \lambda_g$
from the open-sector cyclic trace composed with clutching,
controlled by the universal MC class $\Theta$ of
Theorem~\ref{thm:explicit-theta}. Extending the finite horizontal
spine to a curved affine BGG statement requires the module-bar
screening normalization and the convergence statement named in
Theorem~\ref{thm:bgg-from-bar}.

\subsection{Setup and weight decomposition}

Fix $\mathfrak{g} = \mathfrak{sl}_2$ with standard basis
$\{e, f, h\}$, Killing form $(h,h) = 2$, $(e,f) = 1$, and Cartan
subalgebra $\mathfrak{h} = \mathbb{C} h$. The affine Weyl group is
$\mathcal{W} = \mathbb{Z} \rtimes \mathbb{Z}/2$, generated by the
finite reflection $s$ and translations $t_{n\alpha^\vee}$.

The bar complex $\bar{B}^n(\widehat{\mathfrak{sl}}_{2,k})$ carries a
natural $\mathfrak{h}$-grading inherited from the $\mathfrak{h}$-action
on~$\mathfrak{g}$. At each bar degree~$n$, it decomposes into weight
spaces:
\begin{equation}\label{eq:bar-weight-decomposition}
\bar{B}^n = \bigoplus_{\mu \in \mathfrak{h}^*} \bar{B}^n_\mu.
\end{equation}

\begin{computation}[Weight decomposition at degrees~1 and~2;
\ClaimStatusProvedHere]
\label{comp:bgg-weight-decomp}
\index{bar complex!weight decomposition}

\emph{Degree~1.} $\bar{B}^1 = \mathfrak{g}^* \cong \mathfrak{g}$
(via the Killing form). The $\mathfrak{h}$-weights are:
\begin{equation}\label{eq:degree1-weights}
\bar{B}^1_{+\alpha} = \mathbb{C} \cdot e, \qquad
\bar{B}^1_0 = \mathbb{C} \cdot h, \qquad
\bar{B}^1_{-\alpha} = \mathbb{C} \cdot f,
\end{equation}
where $\alpha$ is the positive root with $\alpha(h) = 2$.

\emph{Degree~2.} From Computation~\ref{comp:sl2-bar-matrix},
$\bar{B}^2$ has dimension~$9$ with basis
$\{[J^a \mid J^b]\}_{a,b \in \{e,f,h\}}$. The
$\mathfrak{h}$-weights are:
\begin{center}
\renewcommand{\arraystretch}{1.2}
\begin{tabular}{cl}
\textbf{Weight} & \textbf{Basis elements} \\
\hline
$+2\alpha$ & $[e \mid e]$ \\
$+\alpha$ & $[e \mid h],\; [h \mid e]$ \\
$0$ & $[e \mid f],\; [f \mid e],\; [h \mid h]$ \\
$-\alpha$ & $[f \mid h],\; [h \mid f]$ \\
$-2\alpha$ & $[f \mid f]$ \\
\end{tabular}
\end{center}
The weight multiplicities $(1, 2, 3, 2, 1)$ are precisely the dimensions
of $(\mathfrak{sl}_2^{\otimes 2})_\mu$.
\end{computation}

\subsection{Bar matrix and the BGG map}

\begin{computation}[Bar differential as BGG differential;
\ClaimStatusProvedHere]
\label{comp:bgg-differential}
\index{BGG resolution!bar differential}

The bar differential $d: \bar{B}^2 \to \bar{B}^1$ from the $3 \times 9$
matrix of Computation~\ref{comp:sl2-bar-matrix} acts by:
\begin{align}
d([e \mid f]) &= h, &
d([f \mid e]) &= -h, \label{eq:bgg-diff-bracket}\\
d([h \mid e]) &= 2e, &
d([e \mid h]) &= -2e, \\
d([h \mid f]) &= -2f, &
d([f \mid h]) &= 2f, \\
d([e \mid e]) &= d([f \mid f]) = d([h \mid h]) &= 0. \label{eq:bgg-diff-zero}
\end{align}
These are precisely the structure constants of $\mathfrak{sl}_2$:
the map $d([J^a \mid J^b]) = [J^a, J^b]$ is the Lie bracket.
The kernel $\ker d \subset \bar{B}^2$ has dimension~$6$, spanned by:
\begin{enumerate}[label=\textup{(\alph*)}]
\item \emph{Symmetric tensors}: $[e \mid e]$, $[f \mid f]$, $[h \mid h]$
 (killed because the bracket is skew-symmetric);
\item \emph{Cartan relations}: $[e \mid f] + [f \mid e]$
 (maps to $h + (-h) = 0$);
\item \emph{Weight combinations}: $[e \mid h] + [h \mid e]$,
 $[f \mid h] + [h \mid f]$.
\end{enumerate}
These kernel elements are the \emph{quadratic dual relations}
$R^\perp$ of the $\mathfrak{sl}_2$ OPE. They determine the
degree-$2$ piece of the dual coalgebra $A^i=H^\star B(A)$; the
quadratic part of $A^!$ is obtained only after graded
linear/Verdier duality, and this calculation does not compute
$Z^{\mathrm{der}}_{\mathrm{ch}}(A)$.
\end{computation}

\subsection{Affine lift obligation}

\begin{theorem}[Affine Bar--BGG lift target for
\texorpdfstring{$L(\Lambda_0)$}{L(Lambda0)};
\ClaimStatusConditional]
\label{thm:bgg-sl2-bar-explicit}
\index{BGG resolution!explicit sl2@explicit $\mathfrak{sl}_2$}

For $\widehat{\mathfrak{sl}}_{2,k}$ on the finite-type
module-bar lane of Theorem~\textup{\ref{thm:bgg-from-bar}}, the
affine BGG lift is the following conditional target. Its required
module-bar $E_1$ page for the vacuum module $L(\Lambda_0)$ is
\begin{equation}\label{eq:bgg-E1}
E_1^n = \bigoplus_{\substack{w \in \mathcal{W} \\ \ell(w) = n}}
\mathcal{M}(w \cdot \Lambda_0)
\end{equation}
with $\mathcal{W}$ the affine Weyl group and with $d_1$ given by
normalized screening operator integrals. The finite horizontal
subcomplex is the two-term classical BGG complex
\[
0\to M(-2)\xrightarrow{\,f\,}M(0)\to L(0)\to0.
\]
The full affine assertion requires three inputs not supplied by the
$3\times9$ matrix alone:
\begin{enumerate}[label=\textup{(\roman*)}]
\item identify the affine Weyl summands in the module-bar filtration
without replacing $A^i$ by a Verma module;
\item prove that the screening integrals have the Kac--BGG
normalization and satisfy the affine braid relations;
\item prove convergence to the classical affine BGG resolution and
exclude additional bar-comodule summands.
\end{enumerate}

On the horizontal finite summand, the differential is the screening
integral
\begin{equation}\label{eq:screening-sl2}
S_{\alpha_i} = \oint_{\gamma_i} V_{\alpha_i}(z)\, dz
\end{equation}
integrated over cycles $\gamma_i \subset \overline{C}_2(X)$, where
$V_{\alpha_i}$ is the vertex operator corresponding to a simple
root. The bar differential
$d([J^a \mid J^b]) = [J^a, J^b]$ from
Computation~\textup{\ref{comp:bgg-differential}} is the \emph{algebraic
finite-current shadow} of these screening integrals.
\end{theorem}

\begin{proof}
The finite horizontal assertion is immediate from
Computation~\ref{comp:bgg-differential}: the entry
$d([e\mid f])=h$ gives
$e(fv_0)=h v_0=0$, and $d([h\mid f])=-2f$ gives
$h(fv_0)=-2fv_0$. Thus $fv_0$ is the singular vector of weight
$-2$ in $M(0)$, and the classical finite BGG complex follows from
\cite{BGG76} or \cite[\S10.7]{Kac}.

For the affine lift, Theorem~\ref{thm:bgg-from-bar} supplies the
conditional spectral-sequence target, while
Theorem~\ref{thm:screening-bar} supplies the current-valued
screening operator only after the Wakimoto/bar--BRST comparison and
the curvature convention are fixed. The $3\times9$ matrix proves the
horizontal singular-vector calculation; it does not prove the full
affine Weyl indexing or the absence of extra bar-comodule summands.
The required character check is the Weyl--Kac identity
(Proposition~\ref{prop:weyl-kac-sl2-bar}):
\[
\operatorname{ch}(L(\Lambda_0))
= \sum_{w \in \mathcal{W}} (-1)^{\ell(w)}
 \operatorname{ch}(\mathcal{M}(w \cdot \Lambda_0))
= \frac{\vartheta_3(q)}{\eta(q)^3}.
\]
The denominator $\eta(q)^{-3}$ arises from the 3~bar generators
$\dim \bar{B}^1 = \dim \mathfrak{sl}_2 = 3$, and the numerator
$\vartheta_3$ from the Weyl group sum on $\bar{B}^n$.
\end{proof}

\begin{remark}[Bar data behind BGG]
\label{rem:bar-adds-to-bgg}
The classical resolution~\eqref{eq:bgg-from-bar} and the finite
two-term complex~\eqref{eq:bgg-sl2-classical} are representation
theory. The bar complex records three additional pieces of data:
\begin{enumerate}[label=\textup{(\alph*)}]
\item A \emph{geometric origin}: the finite-current matrix entries are residues
 of propagators on Fulton--MacPherson compactifications, not
 abstract intertwining operators;
\item \emph{Higher genus corrections}: at genus $g \geq 1$, the scalar
 projection of the bar complex carries curvature
 $m_0^{(g)} = \kappa(\cA) \cdot \lambda_g$ (Theorem~\ref{thm:genus-universality}), which deforms the BGG
 differential on the scalar lane;
\item \emph{Koszul duality}: the bar resolution for
 $\widehat{\mathfrak{sl}}_{2,k}$ maps first to a bar-comodule; the
 rewrite as a dual-level module uses the finite-type graded-dual lane
 of Proposition~\ref{prop:vacuum-verma-koszul}.
\end{enumerate}
\end{remark}


\section{\texorpdfstring{Non-simply-laced example: $\widehat{G}_2$
(type $G_2$)}{Non-simply-laced example: G-2 (type G-2)}}
\label{sec:G2-details}
\index{non-simply-laced!bar complex!G2@$G_2$}
\index{G2@$G_2$!bar complex|textbf}

After the $B_2$ computation of \S\ref{sec:B2-details},
the most exceptional non-simply-laced simple Lie algebra is
$G_2$. This algebra has lacing number $r^\vee = 3$
(the ratio of long to short root squared lengths), the largest
among simple Lie algebras. Its Coxeter and dual Coxeter numbers
differ: $h = 6 \neq h^\vee = 4$.

\subsection{Root system and structure constants}

\begin{computation}[\texorpdfstring{$G_2$}{G_2} root system]\label{comp:G2-roots}
\index{G2@$G_2$!root system}
The Lie algebra $\mathfrak{g}_2$ has:
\begin{itemize}
\item Rank: $2$
\item Dimension: $\dim(\mathfrak{g}_2) = 14$
\item Simple roots: $\alpha_1$ (short), $\alpha_2$ (long), with
$|\alpha_1|^2 = 2/3$, $|\alpha_2|^2 = 2$
(normalization: $(\theta|\theta) = 2$ where
$\theta = 3\alpha_1 + 2\alpha_2$ is the highest root)
\item Lacing number: $r^\vee = |\alpha_2|^2/|\alpha_1|^2 = 3$
\item Positive roots ($6$ total):
\begin{align*}
&\alpha_1 \text{ (short)}, \quad
\alpha_2 \text{ (long)}, \quad
\alpha_1 + \alpha_2 \text{ (short)}, \\
&2\alpha_1 + \alpha_2 \text{ (short)}, \quad
3\alpha_1 + \alpha_2 \text{ (long)}, \quad
3\alpha_1 + 2\alpha_2 \text{ (long)}
\end{align*}
\item Cartan matrix:
$A = \begin{pmatrix} 2 & -3 \\ -1 & 2 \end{pmatrix}$
\item Dual Coxeter number: $h^\vee = 4$
\item Coxeter number: $\mathsf{h} = 6$
\end{itemize}
The three short roots $\{\alpha_1,\; \alpha_1+\alpha_2,\;
2\alpha_1+\alpha_2\}$ form a root system of type $A_2$,
while the three long roots
$\{\alpha_2,\; 3\alpha_1+\alpha_2,\; 3\alpha_1+2\alpha_2\}$
also form a root system of type $A_2$.
\end{computation}

\begin{computation}[\texorpdfstring{$G_2$}{G_2} structure constants]\label{comp:G2-structure}
\index{G2@$G_2$!structure constants}
Use the Chevalley basis
$\{H_1, H_2, E_\alpha, F_\alpha\}_{\alpha \in \Delta_+}$
and write $E_i = E_{\alpha_i}$ for simple roots:
\begin{align}
[E_1, E_2] &= N_{12}E_{12} \quad\text{(root $\alpha_1+\alpha_2$, short)} \\
[E_1, E_{12}] &= N_{1,12}E_{112} \quad\text{(root $2\alpha_1+\alpha_2$, short)} \\
[E_1, E_{112}] &= N_{1,112}E_{1112} \quad\text{(root $3\alpha_1+\alpha_2$, long)} \\
[E_2, E_{112}] &= 0
 \quad\text{($2\alpha_1+2\alpha_2$ is not a root)} \\
[E_1, E_{1112}] &= 0 \\
[E_2, E_{1112}] &= N_{2,1112}E_{11122}
 \quad\text{(root $3\alpha_1+2\alpha_2$, long)}
\end{align}
where the nonzero $N_{\alpha,\beta}$ are the Chevalley constants.
The table uses only the root-string facts: $2\alpha_1+2\alpha_2$
is absent, while $3\alpha_1+2\alpha_2$ is present
(see~\cite{Kac}, \S7.8).

The Serre relations are:
\begin{align}
(\operatorname{ad} E_2)^2(E_1) &= 0
 \quad\text{(long root, since $A_{21} = -1$)}
 \label{eq:serre-G2-long} \\
(\operatorname{ad} E_1)^4(E_2) &= 0
 \quad\text{(short root, since $A_{12} = -3$)}
 \label{eq:serre-G2-short}
\end{align}
Note: $1 - A_{12} = 1 - (-3) = 4$, so the short-root Serre
relation involves \emph{four} applications of $\operatorname{ad} E_1$.
\end{computation}

\subsection{Affine OPE and bar differential}

\begin{computation}[\texorpdfstring{$\widehat{G}_{2,k}$}{G_2,k} OPE]\label{comp:G2-ope}
\index{OPE!G2@$\widehat{G}_2$}
The affine vertex algebra $V_k(\mathfrak{g}_2)$ has $14$ generating
currents. The OPEs are:
\begin{align}
H_i(z)\, H_j(w) &\sim \frac{k\, d_i A_{ij}}{(z-w)^2}
 \label{eq:G2-HH} \\
E_\alpha(z)\, F_\alpha(w) &\sim
 \frac{k |\alpha|^2/2}{(z-w)^2}
 + \frac{H_\alpha(w)}{z-w}
 \label{eq:G2-EF}
\end{align}
where $d_1 = 1/3$, $d_2 = 1$ are the symmetrizing constants
($d_i A_{ij}$ is the symmetrized Cartan matrix).
The double-pole coefficients depend on root length:
\begin{center}
\renewcommand{\arraystretch}{1.2}
\begin{tabular}{lccc}
\textbf{Root type} & \textbf{$|\alpha|^2$} &
\textbf{Double-pole coeff.\ in $E_\alpha F_\alpha$ OPE} &
\textbf{Count} \\
\hline
Short & $2/3$ & $k/3$ & $3$ roots \\
Long & $2$ & $k$ & $3$ roots
\end{tabular}
\end{center}
The ratio $3:1$ between long and short double-pole coefficients
(compared to $2:1$ for $B_2$) is the fingerprint of the
lacing number $r^\vee = 3$.
\end{computation}

\begin{computation}[Degree-2 bar differential for \texorpdfstring{$\widehat{G}_{2,k}$}{G_2,k}]
\label{comp:G2-bar-deg2}
\index{bar complex!G2@$\widehat{G}_2$!degree 2}

\emph{Dimensions.}
$\bar{B}^1 = \mathfrak{g}_2^* \cong \mathbb{C}^{14}$,
$\bar{B}^2 = \mathfrak{g}_2^{\otimes 2} \otimes
\Omega^1(\overline{\mathrm{Conf}}_2) = \mathbb{C}^{196}$
(since $\dim \Omega^1 = 1$).

\emph{Bracket terms (simple pole).}
\begin{align}
d[H_{1}|E_1] &= A_{11}\, E_1 = 2E_1 \\
d[H_{1}|E_2] &= A_{12}\, E_2 = -3E_2 \\
d[H_{2}|E_1] &= A_{21}\, E_1 = -E_1 \\
d[H_{2}|E_2] &= A_{22}\, E_2 = 2E_2 \\
d[E_1|E_2] &= N_{12}E_{12} \\
d[E_1|E_{12}] &= N_{1,12}E_{112} \\
d[E_1|E_{112}] &= N_{1,112}E_{1112}
\end{align}
The asymmetry $A_{12} = -3$ (vs.\ $A_{21} = -1$) produces
markedly different structures on the two sides.

\emph{Curvature terms (double pole).}
\begin{align}
d[E_1|F_1] &\ni \frac{k}{3}\cdot\mathbf{1}
 \quad\text{(short root: $|\alpha_1|^2/2 = 1/3$)} \\
d[E_2|F_2] &\ni k\cdot\mathbf{1}
 \quad\text{(long root: $|\alpha_2|^2/2 = 1$)} \\
d[E_{12}|F_{12}] &\ni \frac{k}{3}\cdot\mathbf{1}
 \quad\text{(short root $\alpha_1+\alpha_2$)} \\
d[E_{112}|F_{112}] &\ni \frac{k}{3}\cdot\mathbf{1}
 \quad\text{(short root $2\alpha_1+\alpha_2$)} \\
d[E_{1112}|F_{1112}] &\ni k\cdot\mathbf{1}
 \quad\text{(long root $3\alpha_1+\alpha_2$)} \\
d[E_{11122}|F_{11122}] &\ni k\cdot\mathbf{1}
 \quad\text{(long root $3\alpha_1+2\alpha_2$)}
\end{align}
\end{computation}

\subsection{Curvature and level-shifting}

\begin{computation}[Curvature and affine shadow for \texorpdfstring{$\widehat{G}_{2,k}$}{G_2,k}]
\label{comp:G2-curvature}
\index{curvature!G2@$\widehat{G}_2$}

\emph{Curvature.}
\begin{equation}
m_0 = \frac{k + h^\vee}{2h^\vee}\cdot\kappa
= \frac{k+4}{8}\cdot\kappa_{\mathfrak{g}_2}
\end{equation}
This vanishes at the critical level $k = -h^\vee = -4$.

\emph{Affine same-family shadow.}
\begin{equation}\label{eq:G2-level-shift}
k'=-k-2h^\vee=-k-8.
\end{equation}
The shifted affine algebra $\widehat{G}_{2,k'}$ records the same
scalar dual lane; it is not the raw bar coalgebra and not the derived
chiral centre. The shift uses $h^\vee = 4$.
Note: the Coxeter number $\mathsf{h} = 6$ does \emph{not} appear in the
level shift. This is the critical distinction for non-simply-laced
algebras: the Feigin--Frenkel level shift uses the dual Coxeter number $h^\vee$, not the Coxeter number $\mathsf{h}$.

\emph{Central charge.}
\begin{equation}
c(\widehat{G}_{2,k}) = \frac{k \cdot 14}{k + 4}
\end{equation}
Complementarity:
\begin{equation}
c(\widehat{G}_{2,k}) + c(\widehat{G}_{2,-k-8})
= \frac{14k}{k+4} + \frac{14(-k-8)}{-k-8+4}
= \frac{14k}{k+4} + \frac{-14(k+8)}{-(k+4)}
= \frac{14k + 14(k+8)}{k+4} = 28
\end{equation}
giving $c + c' = 2\dim(\mathfrak{g}_2) = 28$
(Theorem~\ref{thm:central-charge-complementarity}).

\emph{Obstruction coefficient.}
From the genus universality formula
(Theorem~\ref{thm:genus-universality}):
\begin{equation}\label{eq:G2-kappa}
\kappa(\widehat{G}_{2,k}) = \frac{(k+4) \cdot 14}{2 \cdot 4}
= \frac{7(k+4)}{4}
\end{equation}
(using $t = k + h^\vee = k + 4$, $d = 14$, $h^\vee = 4$).
At $k = 1$ (the minimal non-critical level):
$\kappa = 7 \cdot 5/4 = 35/4$.
\end{computation}

\subsection{Genus expansion}

\begin{computation}[\texorpdfstring{$\widehat{G}_{2,k}$}{G_2,k} genus expansion]
\label{comp:G2-genus}
\index{G2@$G_2$!genus expansion}
\index{genus expansion!G2@$G_2$}
The following is the uniform-weight scalar projection. By genus
universality:
\begin{equation}
F_g(\widehat{G}_{2,k}) = \kappa(\widehat{G}_{2,k}) \cdot \lambda_g^{FP} \quad = \frac{7(k+4)}{4} \cdot
\frac{2^{2g-1}-1}{2^{2g-1}} \cdot \frac{|B_{2g}|}{(2g)!}
\end{equation}

\emph{Explicit values at $k = 1$} ($\kappa = 35/4$):
\begin{center}
\renewcommand{\arraystretch}{1.4}
\begin{tabular}{c|c|c}
$g$ & $F_g$ (exact) & $F_g$ (decimal) \\
\hline
$1$ & $35/96$ & $3.65 \times 10^{-1}$ \\
$2$ & $49/4608$ & $1.06 \times 10^{-2}$ \\
$3$ & $31/110592$ & $2.80 \times 10^{-4}$ \\
$4$ & $127/17694720$ & $7.18 \times 10^{-6}$ \\
$5$ & $511/2802843648$ & $1.82 \times 10^{-7}$ \\
$6$ & $1414477/306070526361600$ & $4.62 \times 10^{-9}$ \\
$7$ & $8191/69958977454080$ & $1.17 \times 10^{-10}$ \\
$8$ & $16931177/5708652560252928000$ & $2.97 \times 10^{-12}$
\end{tabular}
\end{center}

\emph{Complementarity.}
\begin{equation}
\kappa(\widehat{G}_{2,k}) + \kappa(\widehat{G}_{2,-k-8})
= \frac{7(k+4)}{4} + \frac{7(-k-4)}{4} = 0
\end{equation}
so $F_g + F_g' = 0$ at every genus: exact cancellation, as for all
Kac--Moody algebras
(Theorem~\ref{thm:complementarity-root-datum}).
\end{computation}

\subsection{Serre relations in the bar complex}

\begin{computation}[Serre relations for \texorpdfstring{$G_2$}{G_2}: bar complex encoding]
\label{comp:G2-serre}
\index{Serre relations!G2@$G_2$}

\emph{(a) Long-root Serre:
$(\operatorname{ad} E_2)^2(E_1) = 0$.}
Identical in structure to the $B_2$ long-root Serre
(Computation~\ref{comp:B2-serre}(a)).
Bar degree: $3$. Element:
\[
S_{\mathrm{long}} = [E_2|E_2|E_1] - 2[E_2|E_1|E_2]
+ [E_1|E_2|E_2] \in \bar{B}^3
\]

\emph{(b) Short-root Serre: $(\operatorname{ad} E_1)^4(E_2) = 0$.}
This is the most extreme Serre relation among simple Lie algebras,
reflecting $A_{12} = -3$. Bar degree: $\mathbf{5}$. Element:
\begin{align}
S_{\mathrm{short}} &= [E_1|E_1|E_1|E_1|E_2]
- 4[E_1|E_1|E_1|E_2|E_1]
+ 6[E_1|E_1|E_2|E_1|E_1] \notag \\
&\quad\; - 4[E_1|E_2|E_1|E_1|E_1]
+ [E_2|E_1|E_1|E_1|E_1]
\in \bar{B}^5
\end{align}
with binomial coefficients
$\binom{4}{0}, -\binom{4}{1}, \binom{4}{2}, -\binom{4}{3},
\binom{4}{4}$
$= 1, -4, 6, -4, 1$.

\emph{Structural comparison.}
\begin{center}
\renewcommand{\arraystretch}{1.2}
\begin{tabular}{lccccc}
\textbf{Type} & \textbf{$r^\vee$} &
\textbf{Short Serre} & \textbf{Bar degree} &
\textbf{Long Serre} & \textbf{Bar degree} \\
\hline
$A_n$ & $1$ & (same) & $3$ & (same) & $3$ \\
$B_2$ & $2$ & $(\operatorname{ad} E_2)^3(E_1)$ & $4$ &
 $(\operatorname{ad} E_1)^2(E_2)$ & $3$ \\
$\mathbf{G_2}$ & $\mathbf{3}$ &
 $\mathbf{(\operatorname{ad} E_1)^4(E_2)}$ & $\mathbf{5}$ &
 $(\operatorname{ad} E_2)^2(E_1)$ & $3$
\end{tabular}
\end{center}
The general pattern: a Serre relation with Cartan entry $A_{ji} = -m$
produces a bar cycle of degree $m+2$, with binomial coefficients
from $\binom{m+1}{s}$. For $G_2$, the maximum bar degree needed
to encode all Serre relations is $5$ (the highest among all rank-2
Lie algebras).
\end{computation}

\begin{proposition}[Bar complex dimensions for \texorpdfstring{$\widehat{G}_{2,k}$}{G_2,k};
\ClaimStatusProvedHere]
\label{prop:G2-bar-dims}
\index{bar complex!dimension table!G2@$G_2$}
The dimensions of the bar complex through degree~$5$:
\begin{center}
\renewcommand{\arraystretch}{1.2}
\begin{tabular}{c|ccc|c}
$n$ & $\dim(\mathfrak{g}_2^{\otimes n})$ &
$\dim \Omega^{n-1}$ & $\dim \bar{B}^n$ &
Notes \\
\hline
$0$ & $1$ & $1$ & $1$ & vacuum \\
$1$ & $14$ & $1$ & $14$ & generators \\
$2$ & $196$ & $1$ & $196$ & OPE (asymmetric poles) \\
$3$ & $2744$ & $2$ & $5488$ & long-root Serre ($\times 1$) \\
$4$ & $38416$ & $6$ & $230496$ & higher Serre \\
$5$ & $537824$ & $24$ & $12\,907\,776$ & short-root Serre ($\times 1$)
\end{tabular}
\end{center}
\end{proposition}

\begin{proof}
The generator dimensions are $14^n$; the form dimensions
are $(n-1)!$ (Poincar\'{e} polynomial of $C_n(\mathbb{C})$).
The Serre contribution analysis follows
Computation~\ref{comp:G2-serre}: the long-root Serre (degree $3$)
contributes by the same mechanism as for $B_2$
(Computation~\ref{comp:B2-serre}(a)), while the short-root
Serre (degree $5$) uses the $G_2$-specific fourth-order relation.
\end{proof}

\begin{remark}[The \texorpdfstring{$G_2$}{G_2} computation in context]\label{rem:G2-context}
\index{non-simply-laced!comparison table}
The $B_2$ (\S\ref{sec:B2-details}) and $G_2$ computations together
exhibit the three possible lacing numbers
$r^\vee \in \{1, 2, 3\}$:
\begin{center}
\renewcommand{\arraystretch}{1.2}
\begin{tabular}{lcccccc}
\textbf{Type} & $\dim$ & $h$ & $h^\vee$ & $r^\vee$ &
\textbf{Level shift} & \textbf{Max Serre deg.} \\
\hline
$A_1$ & $3$ & $2$ & $2$ & $1$ & $-k-4$ & $3$ \\
$A_2$ & $8$ & $3$ & $3$ & $1$ & $-k-6$ & $3$ \\
$B_2 \cong C_2$ & $10$ & $4$ & $3$ & $2$ & $-k-6$ & $4$ \\
$G_2$ & $14$ & $6$ & $4$ & $3$ & $-k-8$ & $5$ \\
$F_4$ & $52$ & $12$ & $9$ & $2$ & $-k-18$ & $4$
\end{tabular}
\end{center}
As the lacing number increases, the bar complex requires higher
degrees to encode all the algebraic structure. The universal
formulas (Theorems~\ref{thm:universal-kac-moody-koszul}
and~\ref{thm:genus-universality}) hold uniformly, but the
\emph{complexity} of the bar complex grows with $r^\vee$.
\end{remark}

\begin{remark}[Complete shadow data for exceptional and non-simply-laced types]
\label{rem:exceptional-nsl-shadow-summary}
\index{shadow obstruction tower!exceptional types}
\index{shadow obstruction tower!non-simply-laced types}
\index{E6@$E_6$!shadow data}
\index{E7@$E_7$!shadow data}
\index{F4@$F_4$!shadow data}

The following table collects the shadow obstruction tower data for all exceptional ($E_6, E_7, E_8$) and non-simply-laced ($B_2, G_2, F_4$) non-critical universal affine Kac--Moody algebras. Every family is class L (shadow depth $r_{\max} = 3$), with cubic shadow $C$ from the Lie bracket, quartic shadow $S_4 = 0$ by the Jacobi identity, and critical discriminant $\Delta = 8\kappa S_4 = 0$. The r-matrix has a simple pole at $z = 0$ ($d\log$ extraction absorbs one power), and complementarity $\kappa + \kappa' = 0$ holds for all Kac--Moody families.
\begin{center}
\renewcommand{\arraystretch}{1.3}
\begin{tabular}{lcccccccl}
\toprule
$\mathfrak{g}$ & $d$ & $r$ & $h$ & $h^\vee$ & $r^\vee$ & $\kappa(\widehat{\mathfrak{g}}_k)$ & $c + c'$ & Class \\
\midrule
\multicolumn{9}{c}{\textit{Exceptional simply-laced ($h = h^\vee$, $r^\vee = 1$)}} \\
$E_6$ & $78$ & $6$ & $12$ & $12$ & $1$ & $\frac{13(k+12)}{4}$ & $156$ & L \\[2pt]
$E_7$ & $133$ & $7$ & $18$ & $18$ & $1$ & $\frac{133(k+18)}{36}$ & $266$ & L \\[2pt]
$E_8$ & $248$ & $8$ & $30$ & $30$ & $1$ & $\frac{62(k+30)}{15}$ & $496$ & L \\
\midrule
\multicolumn{9}{c}{\textit{Non-simply-laced ($h \neq h^\vee$)}} \\
$B_2 \cong C_2$ & $10$ & $2$ & $4$ & $3$ & $2$ & $\frac{5(k+3)}{3}$ & $20$ & L \\[2pt]
$G_2$ & $14$ & $2$ & $6$ & $4$ & $3$ & $\frac{7(k+4)}{4}$ & $28$ & L \\[2pt]
$F_4$ & $52$ & $4$ & $12$ & $9$ & $2$ & $\frac{26(k+9)}{9}$ & $104$ & L \\
\bottomrule
\end{tabular}
\end{center}
The kappa formula uses $h^\vee$ (not $h$) throughout: $\kappa = d(k+h^\vee)/(2h^\vee)$. For $E_6$ at $k=1$: $\kappa = 13 \cdot 13/4 = 169/4$. For $G_2$ at $k=1$: $\kappa = 7 \cdot 5/4 = 35/4$. For $F_4$ at $k=1$: $\kappa = 26 \cdot 10/9 = 260/9$. Every affine Kac--Moody algebra, regardless of type or lacing number, lies in class L.
\end{remark}


\section{Virasoro bar complex: degree 3}
\label{sec:virasoro-degree3-details}
\index{Virasoro algebra!bar complex!degree 3}

The degree-2 Virasoro bar differential was computed in
Computation~\ref{comp:virasoro-bar-diff}. Degree~3 introduces
the composite field $\Lambda = \normord{TT} -
\frac{3}{10}\partial^2 T$.

\begin{computation}[Degree-3 bar differential for \texorpdfstring{$\mathrm{Vir}_c$}{Vir_c}]
\label{comp:virasoro-degree3}

\emph{Degree-3 elements.}
The bar complex $\bar{B}^3(\mathrm{Vir}_c) = \bar{V}^{\otimes 3}
\otimes \Omega^2(\overline{\mathrm{Conf}}_3)$, where
$\bar{V} = V/\mathbb{C}|0\rangle$ is the augmentation ideal
of the Virasoro vacuum module. Low-weight elements include:
\begin{center}
\renewcommand{\arraystretch}{1.2}
\begin{tabular}{cll}
\textbf{Weight} & \textbf{Elements} & \textbf{Form basis} \\
\hline
$(2,2,2)$ & $[T|T|T]$ & $\eta_{12}\wedge\eta_{13},\;\eta_{13}\wedge\eta_{23}$ \\
$(2,2,3)$ & $[T|T|\partial T]$ & same \\
$(2,3,2)$ & $[T|\partial T|T]$ & same \\
$(3,2,2)$ & $[\partial T|T|T]$ & same
\end{tabular}
\end{center}

\emph{The element $[T|T|T] \otimes \eta_{12}\wedge\eta_{13}$.}
The differential $d\colon \bar{B}^3 \to \bar{B}^2$ acts by three
collision residues.

\emph{Collision $D_{12}$ ($z_1 = z_2$).}
The residue extracts the $T(z_1)T(z_2)$ OPE data:
\begin{align*}
\operatorname{Res}_{D_{12}}\bigl([T|T|T] \otimes \eta_{12}\wedge\eta_{13}\bigr)
&= \bigl(T_{(3)}T \cdot |0\rangle + T_{(1)}T + T_{(0)}T\bigr) \otimes T
\otimes \eta_{13} \\
&= \bigl(\tfrac{c}{2}|0\rangle + 2T + \partial T\bigr)
\otimes T \otimes \eta_{13}
\end{align*}
This gives three terms in $\bar{B}^2$:
\begin{align*}
&\frac{c}{2}\,[T] \otimes \eta_{13}
\quad\text{(from vacuum $\times T$: the term $|0\rangle \otimes T$ reduces)} \\
&+ 2\,[T|T] \otimes \eta_{13}
\quad\text{(from $2T \otimes T$)} \\
&+ [\partial T|T] \otimes \eta_{13}
\quad\text{(from $\partial T \otimes T$)}
\end{align*}

\emph{Collision $D_{13}$ ($z_1 = z_3$).}
Similarly, with the surviving form $\eta_{12}|_{z_1=z_3} = \eta_{23}$:
\begin{align*}
\operatorname{Res}_{D_{13}} &= T \otimes \bigl(\tfrac{c}{2}|0\rangle + 2T + \partial T\bigr)
\otimes \eta_{23} \\
&= \frac{c}{2}\,[T] \otimes \eta_{23}
+ 2\,[T|T] \otimes \eta_{23}
+ [T|\partial T] \otimes \eta_{23}
\end{align*}
(Note: the form changes from $\eta_{13}$ to $\eta_{23}$ because after
$z_1 = z_3$, the surviving pair is $(z_{13}, z_2)$ with $\eta_{z_{13},z_2}$.)

\emph{Collision $D_{23}$ ($z_2 = z_3$).}
With surviving form $\eta_{12}|_{z_2=z_3} = \eta_{12}$:
\begin{align*}
\operatorname{Res}_{D_{23}} &= T \otimes \bigl(\tfrac{c}{2}|0\rangle + 2T + \partial T\bigr)
\otimes \eta_{12} \\
&= \frac{c}{2}\,[T] \otimes \eta_{12}
+ 2\,[T|T] \otimes \eta_{12}
+ [T|\partial T] \otimes \eta_{12}
\end{align*}

\emph{Assembling.}
\begin{align}
d\bigl([T|T|T] \otimes \eta_{12}\wedge\eta_{13}\bigr)
&= \frac{c}{2}\bigl([T] \otimes \eta_{13}
+ [T] \otimes \eta_{23}
+ [T] \otimes \eta_{12}\bigr) \notag \\
&\quad + 2\bigl([T|T] \otimes \eta_{13}
+ [T|T] \otimes \eta_{23}
+ [T|T] \otimes \eta_{12}\bigr) \notag \\
&\quad + [\partial T|T] \otimes \eta_{13}
+ [T|\partial T] \otimes (\eta_{23} + \eta_{12})
\label{eq:vir-deg3-TTT}
\end{align}

The degree-1 terms
\[
\frac c2[T]\otimes(\eta_{12}+\eta_{13}+\eta_{23})
\]
do \emph{not} vanish by an Arnold relation. The Arnold relation on
three points is the quadratic Orlik--Solomon identity
\[
\eta_{12}\wedge\eta_{23}+\eta_{23}\wedge\eta_{31}
+\eta_{31}\wedge\eta_{12}=0,
\]
not a linear identity among the one-forms $\eta_{12},\eta_{13},\eta_{23}$.
Thus the quartic Virasoro pole propagates into degree~3 as a genuine
curvature insertion. It cancels only in the curved $A_\infty$ identity
$D^2=[m_0,-]$ after the $m_0=c/2$ component is included; it is not
zero in the ordinary residue differential.

The degree-$3$ differential is therefore:
\begin{align}
d\bigl([T|T|T] \otimes \eta_{12}\wedge\eta_{13}\bigr)
&= \frac{c}{2}[T]\otimes(\eta_{12}+\eta_{13}+\eta_{23}) \notag \\
&\quad + 2[T|T] \otimes (\eta_{12} + \eta_{13} + \eta_{23}) \notag \\
&\quad + [\partial T|T] \otimes \eta_{13}
+ [T|\partial T] \otimes (\eta_{12} + \eta_{23}) \notag \\
\label{eq:vir-deg3-TTT-simplified}
\end{align}
subject only to the quadratic Arnold identities among wedge products.
\end{computation}

\begin{proposition}[Virasoro curvature survives the degree-\texorpdfstring{$3$}{3} residue;
\ClaimStatusProvedHere]
\label{prop:arnold-virasoro-deg3}
\index{Arnold relation!Virasoro cancellation}
For the Virasoro bar complex at any central charge~$c$:
\begin{enumerate}[label=\textup{(\roman*)}]
\item The vacuum leakage from the quartic pole $T_{(3)}T=c/2$
survives in degree~$3$ as
$\frac c2[T]\otimes(\eta_{12}+\eta_{13}+\eta_{23})$.
\item The Arnold relation cancels quadratic wedge expressions, not
linear sums of logarithmic one-forms. Hence finite checks in degrees
$3,4,5$ cannot prove vanishing of Virasoro curvature in every degree.
\item The correct mechanism in every degree is the curved identity
$D^2=[m_0,-]$ with $m_0=c/2$, together with the higher Virasoro
shadow tower.
\end{enumerate}
\end{proposition}

\begin{proof}
The computation above gives (i). For (ii), Arnol'd's relation in the
Orlik--Solomon algebra of $\Conf_3(\mathbb C)$ is
$\eta_{ij}\wedge\eta_{jk}+\eta_{jk}\wedge\eta_{ki}
+\eta_{ki}\wedge\eta_{ij}=0$. Since it is homogeneous of degree~$2$,
it cannot imply a linear dependence among the degree-$1$ generators
$\eta_{ij}$. Assertion (iii) is precisely the curved bar identity for
the Virasoro OPE: the scalar term $T_{(3)}T=c/2$ is the curvature
component, while the double and simple poles supply the weight and
translation components.
\end{proof}


\section{Yangian bar complex details}
\label{sec:yangian-details}

\begin{computation}[Yangian \texorpdfstring{$Y(\mathfrak{sl}_2)$}{Y(sl_2)} relations]\label{comp:yangian-sl2}
Generators: $e^{(r)}, f^{(r)}, h^{(r)}$ for $r \geq 0$.

\emph{Level~0} (= $\mathfrak{sl}_2$):
\[
[h^{(0)}, e^{(0)}] = 2e^{(0)}, \quad [h^{(0)}, f^{(0)}] = -2f^{(0)}, \quad [e^{(0)}, f^{(0)}] = h^{(0)}
\]

\emph{Level~1.}
\begin{align*}
[h^{(0)}, e^{(1)}] &= 2e^{(1)} \\
[h^{(1)}, e^{(0)}] &= 2e^{(1)} \\
[e^{(0)}, f^{(1)}] &= h^{(1)} \\
[e^{(1)}, f^{(0)}] &= h^{(1)} \\
[e^{(1)}, f^{(1)}] - [e^{(0)}, f^{(2)}] &= \frac{1}{4}\{h^{(0)}, h^{(1)}\}
\end{align*}

\emph{Yangian current.}
\[
e(u) = \sum_{r \geq 0} e^{(r)} u^{-r-1}
\]

\emph{RTT presentation.}
\[
R(u-v) T_1(u) T_2(v) = T_2(v) T_1(u) R(u-v)
\]
with $R(u) = 1 - \frac{\hbar P}{u}$ and $T(u) = \begin{pmatrix} k_+(u) & e(u) \\ f(u) & k_-(u) \end{pmatrix}$ (matching the sign convention of Definition~\ref{def:yangian-rtt}).
\end{computation}

\begin{computation}[Yangian bar complex]\label{comp:yangian-bar-full}
\emph{Degree~1 generators.} $[e^{(r)}_{-n}], [f^{(r)}_{-n}], [h^{(r)}_{-n}]$ for $r \geq 0$, $n > 0$.

\emph{Degree~2 differential with R-matrix.}

The Yangian $R$-matrix $R(u) = 1 - \hbar P/u$ (where $P$ is the permutation and $u$ is the spectral parameter; cf.\ Definition~\ref{def:yangian-rtt}) deforms the bar differential. At the level of the Yangian generating series $T^a(u) = \sum_{r \geq 0} T^{a,(r)} u^{-r-1}$, the deformed differential on degree-2 elements is:
\[
d^R[a^{(r)}|b^{(s)}] = \sum_{t=0}^{r+s} [a^{(r)}, b^{(s)}]^{(t)} + \text{(R-matrix corrections at level } r+s+1 \text{)}
\]
where the Yangian brackets $[a^{(r)}, b^{(s)}]^{(t)}$ encode the RTT relations.

The \emph{loop modes} $a^{(r)}_{-n}$ (combining Yangian level $r$ with affine mode $-n$) carry two independent gradings. The bar differential must respect both: the Yangian level $r$ governs the $R$-matrix corrections, while the affine mode $n$ governs the Lie bracket. For the simplest case ($r = s = 0$):
\[d^R[e^{(0)}_{-m}|f^{(0)}_{-n}] = [h^{(0)}_{-(m+n)}] + O(\text{level } \geq 1)\]
where the leading term is the standard affine Lie algebra bracket and higher-level corrections involve $T^{(r)}$ for $r \geq 1$.

\emph{Cohomology.} Related to Yangian cohomology $H^*(Y(\mathfrak{sl}_2))$.
\end{computation}


\subsection{\texorpdfstring{$q$-deformed computations}{q-deformed computations}}\label{sec:q-deformed-details}

\begin{computation}[\texorpdfstring{$q$}{q}-Heisenberg commutator expansion]\label{comp:q-comm}
The $q$-commutator:
\[
[a_m, a_n]_q = q^{\mathrm{sgn}(n-m)/2} a_m a_n - q^{\mathrm{sgn}(m-n)/2} a_n a_m
\]

\emph{For $m < n$.}
\[
[a_m, a_n]_q = q^{1/2} a_m a_n - q^{-1/2} a_n a_m
\]

\emph{For $m > n$.}
\[
[a_m, a_n]_q = q^{-1/2} a_m a_n - q^{1/2} a_n a_m = -[a_n, a_m]_q
\]
(antisymmetry preserved).

\emph{For $m = -n$.}
\[
[a_m, a_{-m}]_q = [m]_q
\]
where $[m]_q = \frac{q^m - q^{-m}}{q - q^{-1}}$.

\emph{Expansion as $q \to 1$.}
\[
[a_m, a_n]_q = [a_m, a_n] + \frac{\ln q}{2}\mathrm{sgn}(n-m)\{a_m, a_n\} + O((\ln q)^2)
\]
The first-order correction involves the anticommutator.
\end{computation}


\subsection{Roots of unity phenomena}

\begin{computation}[\texorpdfstring{$q$}{q}-Heisenberg at roots of unity]\label{comp:q-root-unity}
Let $q = e^{2\pi i/N}$ be a primitive $N$-th root of unity.

\emph{Central elements.} The element $a_0^N$ becomes central:
\[
[a_0^N, a_m] = 0 \quad \text{for all } m
\]
(since $q^{mN} = 1$).

\emph{New relations.}
\[
a_m^N = 0 \quad \text{for } m \neq 0 \text{ (in certain quotients)}
\]

\emph{Bar complex modification.} New generators:
\[
[a_0^N], [a_1^N], \ldots
\]
with modified differential reflecting the truncation.

\emph{New homology.} Classes corresponding to the restricted Lie algebra structure.
\end{computation}


\section{Higher genus formulas}
\label{sec:higher-genus-details}

\subsection{Genus 1 (torus) formulas}

\begin{computation}[Heisenberg on torus]\label{comp:heisenberg-torus}
On the torus $E_\tau = \bC/(\bZ + \tau\bZ)$:

\emph{Green's function.}
\[
G(z, w | \tau) = -\log|\theta_1(z-w|\tau)| + \frac{\pi(\mathrm{Im}(z-w))^2}{\mathrm{Im}(\tau)}
\]
where $\theta_1$ is the odd Jacobi theta function.

\emph{OPE on torus.}
\[
J(z) J(w) = \kappa\,\wp(z-w|\tau) + \text{regular}
\]
where $\wp$ is the Weierstrass $\wp$-function with double pole at $z = w$.

\emph{Mode expansion.}
\[
J(z) = \sum_{n \in \bZ} a_n e^{2\pi i n z}
\]

\emph{Commutation relations.}
\[
[a_m, a_n] = \kappa\, m\,\delta_{m+n,0}
\]
The commutation relations are independent of $\tau$; the genus-1 data enters through the propagator $\wp(z|\tau) - \frac{\pi^2 E_2(\tau)}{3}$, not through the algebra structure.

\emph{Chiral partition function} (character of the Fock space):
\[
Z_{\mathrm{chiral}}(\tau) = \mathrm{tr}_{\mathcal{F}}(q^{L_0 - 1/24}) = \frac{1}{\eta(\tau)} = q^{-1/24} \prod_{n=1}^\infty \frac{1}{1 - q^n}
\]
where $q = e^{2\pi i \tau}$. The full (non-chiral) partition function is $Z(\tau) = |\eta(\tau)|^{-2}$.
\end{computation}


\begin{computation}[Bar complex on higher genus]\label{comp:bar-genus-g}
For genus $g$ curve $\Sigma_g$:

\emph{Configuration space.} $\Conf_n(\Sigma_g)$ with compactification
$\overline{\Conf}_n(\Sigma_g)$.

\emph{Cohomology.}
$H^*(\Conf_n(\Sigma_g))$ is computed by the Cohen--Taylor (or Totaro) spectral sequence, converging from $E_1 = H^*(\Sigma_g^n) \otimes H^*(\text{OS}_n)$ (the Orlik--Solomon algebra of logarithmic forms $\eta_{ij}$ modulo Arnold relations) together with the mixed relations $\eta_{ij} \cdot [\text{diagonal}]_{ij} = 0$. In particular, $H^1(\Sigma_g) \cong \mathbb{C}^{2g}$ contributes new generators.

\emph{New bar complex generators.} For each handle ($A_i$, $B_i$ cycle):
\[
[p_i], [q_i] \quad (i = 1, \ldots, g)
\]
representing the zero-modes from the period integrals.

\emph{Modified differential.}
\[
d[p_i | q_j] = k\delta_{ij} \cdot [\text{fundamental class}]
\]
reflecting the symplectic pairing $\int_{A_i} \omega \cdot \int_{B_j} \omega = \delta_{ij}$.

\emph{Homology.}
\[
H_n(\B(\cH_g)) = \wedge^n(V^* \oplus \bC^{2g})
\]
with the additional $2g$ generators from the handles.
\end{computation}


\section{Structure theorems from bar complex computations}
\label{sec:bar-structure-theorems}
\index{bar complex!structure theorems}

The preceding computations imply the following structural statements.

\begin{proposition}[Heisenberg bar complex: adjacent residues and central class; \ClaimStatusProvedHere]
\label{prop:heisenberg-maximal-form-cycles}
\index{Heisenberg algebra!central class}
For the Heisenberg algebra $\cH_k$ on $\mathbb P^1$, the
maximal-form chain
\[
[a|\cdots|a]\otimes
\eta_{12}\wedge\eta_{23}\wedge\cdots\wedge\eta_{n-1,n}
\]
has nonzero residue differential in general. Its adjacent collision
terms are the alternating contractions by the double-pole coefficient
$a_{(1)}a=k\mathbf 1$, while the non-adjacent collision terms vanish
on this chain basis by form-squaring. Consequently
\[
H^m(\barB^{\mathrm{ch}}_{\mathrm{geom}}(\cH_k))
=
\begin{cases}
\mathbb C & m=0,\\
0 & m=1,\\
\mathbb C\cdot c_k & m=2,\\
0 & m>2,
\end{cases}
\]
where $c_k$ is the level central class of
Theorem~\ref{thm:frame-heisenberg-bar}.
\end{proposition}

\begin{proof}
The bar differential has residue and de~Rham components. The de~Rham
component vanishes on the displayed logarithmic chain because each
$\eta_{ij}=d\log(z_i-z_j)$ is closed. The residue component is not
zero: the Heisenberg OPE
\[
a(z)a(w)\sim \frac{k}{(z-w)^2}
\]
has $a_{(0)}a=0$ and $a_{(1)}a=k\mathbf 1$. Thus there is no bracket
residue, but there is a central double-pole residue. On the chain
basis, exactly the adjacent divisors $D_{12},D_{23},\ldots,D_{n-1,n}$
survive; the non-adjacent divisors repeat one logarithmic factor after
restriction and vanish. The surviving terms form the alternating
boundary of the chain complex, and the Arnold relation gives
$d_{\mathrm{res}}^2=0$.

The cohomology statement is Theorem~\ref{thm:frame-heisenberg-bar}.
It concerns the raw geometric bar complex. The dual coalgebra is
$\cH_k^i=H^\star B(\cH_k)$, and the Verdier--Koszul dual algebra is
$\cH_k^!\simeq(\mathrm{Sym}^{\mathrm{ch}}(V^*),m_0=-k\omega)$.
\end{proof}

\begin{remark}\label{rem:maximal-form-mechanism}
The absence of a simple pole is specific to the Heisenberg and kills
only the bracket channel. It does not kill the whole residue
differential. For the Virasoro
(quartic pole $T_{(3)}T = c/2$), the maximal-form bar differential
also carries curvature, now with scalar $\kappa(\mathrm{Vir}_c)=c/2$
(Computation~\ref{comp:virasoro-curvature}).
For affine Kac--Moody algebras ($\widehat{\fg}_k$ with double pole $\kappa_{ij}$
\emph{and} simple pole $f^{ab}{}_c J^c$), the simple pole contribution
$a_{(0)}b \neq 0$ gives nonzero residues, so the maximal-form
elements are not automatically cycles.
\end{remark}


\begin{proposition}[Universal Kac--Moody acyclicity; critical centre separated; \ClaimStatusProvedHere]
\label{prop:km-generic-acyclicity}
\index{Kac--Moody algebra!universal affine acyclicity}
For any simple Lie algebra $\fg$ and any level $k$, the bar-twisted
complex of the universal affine vacuum module is acyclic:
\[
H_n(\B(\hat{\fg}_k) \otimes_\tau V_k) = \begin{cases}
\bC & n = 0, \\
0 & n > 0.
\end{cases}
\]
At the critical level $k = -h^\vee$, the Feigin--Frenkel centre appears
in the completed critical centre/Hochschild theory; it is not a pole or
failure of the universal bar-cobar resolution.
\end{proposition}

\begin{proof}
Define the PBW filtration $F_p$ by the number of generators applied
to the vacuum. The associated graded is
$\mathrm{gr}_F(\B(\hat{\fg}_k) \otimes_\tau V_k) \cong
B(U(\fg[t^{-1}]t^{-1})) \otimes_{\tau_0}
U(\fg[t^{-1}]t^{-1})$.
This is the augmented bar resolution of the enveloping algebra of
the negative loop algebra. Its homology is $\bC$ in degree~$0$ and
vanishes in positive degree. The spectral sequence therefore has
$E_1=\bC$ on the augmentation line and no positive homological
classes, giving acyclicity of the twisted complex for the universal
affine module. The proof uses PBW freeness of the induced vacuum
module and therefore does not divide by $k+h^\vee$.

At $k = -h^\vee$, what degenerates is the Sugawara construction and
the completed centre, not the PBW bar resolution. The
Feigin--Frenkel centre gives new central/Hochschild classes in the
critical theory. For non-critical admissible simple quotients, the
acyclicity statement must be rechecked after quotienting; it is not a
formal consequence of the universal computation. Computation~\ref{comp:km-acyclic}
gives the argument explicitly for $\fg = \mathfrak{sl}_2$.
\end{proof}


\begin{proposition}[\texorpdfstring{$\mathcal{W}_3$}{W_3} vacuum leakage dichotomy; \ClaimStatusProvedHere]
\label{prop:w3-vacuum-dichotomy}
\index{vacuum leakage!$\mathcal{W}_3$ dichotomy}
\index{W3 algebra@$\mathcal{W}_3$!vacuum leakage}
In the degree-$2$ bar differential of the $\mathcal{W}_3$ algebra,
only the diagonal generator pairs produce vacuum leakage:
\begin{center}
\renewcommand{\arraystretch}{1.2}
\begin{tabular}{lcc}
\textbf{Pair} & \textbf{Vacuum leakage} & \textbf{Mechanism} \\
\hline
$T \otimes T$ & $c/2$ & quartic pole $T_{(3)}T = c/2$ \\
$W \otimes W$ & $c/3$ & sixth-order pole $W_{(5)}W = c/3$ \\
$T \otimes W$ & $0$ & pole order $\leq 2$; need order $\geq 5$ \\
$W \otimes T$ & $0$ & pole order $\leq 2$; need order $\geq 5$
\end{tabular}
\end{center}
Consequently, the curvature of the $\mathcal{W}_3$ bar complex decomposes
into two independent channels, both proportional to~$c$ and both vanishing
simultaneously if and only if $c = 0$.
The ratio $m_0^{(W)}/m_0^{(T)} = 2/3$ is independent of the level~$k$.
\end{proposition}

\begin{proof}
A vacuum contribution from $D(a \otimes b \otimes \eta_{12})
= \sum_{n \geq 0} a_{(n)}b$ requires $a_{(n_0)}b \in \bC \cdot |0\rangle$
for some~$n_0$. The conformal weight condition forces
$n_0 = \mathrm{wt}(a) + \mathrm{wt}(b) - 1$.

\emph{Diagonal pairs.}
For $T \otimes T$: $n_0 = 2 + 2 - 1 = 3$, and $T_{(3)}T = c/2 \neq 0$.
For $W \otimes W$: $n_0 = 3 + 3 - 1 = 5$, and $W_{(5)}W = c/3 \neq 0$.

\emph{Off-diagonal pairs.}
For $T \otimes W$: $n_0 = 2 + 3 - 1 = 4$. But $W$ is primary of
weight~$3$ under the Virasoro, so $T_{(n)}W = 0$ for $n \geq 2$
(the $TW$ OPE has pole order at most~$2$). In particular $T_{(4)}W = 0$.
For $W \otimes T$: $n_0 = 4$, and $W_{(n)}T = 0$ for $n \geq 2$ by
skew-symmetry (Computation~\ref{comp:w3-nthproducts}).

The curvature ratio follows from the OPE normalizations:
$(c/3)/(c/2) = 2/3$, independent of~$k$.
Computation~\ref{comp:w3-bar-degree2} gives the complete residue calculation.
\end{proof}

\begin{remark}[Generalization to \texorpdfstring{$\mathcal{W}_N$}{W_N}]\label{rem:wn-vacuum-channels}
For $\mathcal{W}_N$ with generators $W^{(d_i)}$ of weights $d_1 = 2 < d_2 < \cdots < d_r$
(where $r = N-1$), the vacuum leakage dichotomy generalizes: the pair
$W^{(d_i)} \otimes W^{(d_j)}$ produces vacuum leakage if and only if
$i = j$ (i.e., the pair is diagonal).
This follows from the primary field condition: if $d_i \neq d_j$,
then the OPE pole order is at most $\max(d_i, d_j)$, while the vacuum
condition requires pole order $d_i + d_j - 1 > \max(d_i, d_j)$ when
$\min(d_i, d_j) \geq 2$. Thus the curvature of any $\mathcal{W}_N$
bar complex decomposes into $r$ independent channels, one per generator.
\end{remark}


%% ======================================================================
%% r-MATRIX POLE STRUCTURE ACROSS THE STANDARD LANDSCAPE
%% ======================================================================

\section{Collision-residue \texorpdfstring{$r$}{r}-matrix: pole structure}
\label{sec:rmatrix-pole-structure}
\index{r-matrix@$r$-matrix!pole structure|textbf}
\index{pole absorption!$d\log$ measure}

The bar construction extracts collision residues along
$d\log(z_i - z_j)$, not along $dz/(z-w)$. The $d\log$ measure absorbs
one power of $(z-w)$: a pole of order~$n$ in the OPE becomes a pole of
order~$n-1$ in the collision residue $r(z) = \operatorname{Res}^{\mathrm{coll}}_{0,2}(\Theta_\cA)$.
In particular, simple poles ($z^{-1}$) in the OPE produce
regular ($z^0$) terms in the $r$-matrix and drop.

Table~\ref{tab:rmatrix-pole-landscape} summarises the pole structure
for the standard families.

\begin{table}[ht]
\centering
\renewcommand{\arraystretch}{1.15}
\caption{Collision-residue $r$-matrix pole structure across the
standard landscape. The rule ``OPE pole $z^{-n} \to$
$r$-matrix pole $z^{-(n-1)}$'' is the $d\log(z_i-z_j)$ collision
residue rule.}
\label{tab:rmatrix-pole-landscape}
\begin{tabular}{lcccc}
\toprule
\textbf{Family / channel}
& \textbf{$h$}
& \textbf{OPE poles}
& \textbf{$r$-matrix poles}
& \textbf{$r^{\mathrm{coll}}(z)$} \\
\midrule
$\cH_k$ ($JJ$)
 & $1$ & $\{2\}$ & $\{1\}$ & $k/z$ \\
$\hat{\mathfrak{sl}}_2$ (diag.)
 & $1$ & $\{2\}$ & $\{1\}$ & $k\Omega/z$ \\
$\hat{\mathfrak{sl}}_2$ (off-diag.)
 & $1$ & $\{1\}$ & $\varnothing$ & $0$ \\
$\hat{\mathfrak{sl}}_3$ (diag.)
 & $1$ & $\{2\}$ & $\{1\}$ & $k\Omega/z$ \\
$\hat{\mathfrak{sl}}_3$ (off-diag.)
 & $1$ & $\{1\}$ & $\varnothing$ & $0$ \\
$\mathrm{Vir}_c$ ($TT$)
 & $2$ & $\{4,2,1\}$ & $\{3,1\}$ & $(c/2)/z^3 + 2T/z$ \\
$\cW_3$ ($TT$)
 & $2$ & $\{4,2,1\}$ & $\{3,1\}$ & $(c/2)/z^3 + 2T/z$ \\
$\cW_3$ ($TW$)
 & $2,3$ & $\{2,1\}$ & $\{1\}$ & $3W/z$ \\
$\cW_3$ ($WW$)
 & $3$ & $\{6,4,3,2,1\}$ & $\{5,3,2,1\}$
 & $(c/3)/z^5 + 2T/z^3 + \cdots$ \\
$\beta\gamma$ ($\beta\gamma$)
 & $0,1$ & $\{1\}$ & $\varnothing$ & $0$ \\
$\beta\gamma$ (diag.)
 & n/a & $\varnothing$ & $\varnothing$ & $0$ \\
$bc$ ($\psi_i\psi_i$)
 & $\tfrac{1}{2}$ & $\{1\}$ & $\varnothing$ & $0$ \\
\bottomrule
\end{tabular}
\end{table}

\begin{remark}[Even-order poles in the $\cW_3$ $WW$ $r$-matrix]
\label{rem:w3-ww-even-poles}
\index{W-algebra@$\cW$-algebra!$r$-matrix even poles}
The $WW$ channel of $\cW_3$ has an $r$-matrix pole at $z^{-2}$
(from the $\partial T/(z{-}w)^3$ term in the $WW$ OPE, shifted by~$1$).
This does \emph{not} violate the bosonic parity constraint
(Proposition~\ref{prop:bosonic-parity-constraint}): that constraint
applies to same-generator channels between generators of the
\emph{same} integer weight.
Here $W$ has conformal weight~$3$ (odd), and the $WW$ OPE pole
pattern $\{6, 4, 3, 2, 1\}$ includes the odd-order pole $z^{-3}$,
which shifts to the even-order $z^{-2}$ in the $r$-matrix.
For multi-generator algebras, even-order $r$-matrix poles are
generically present in cross-weight channels.
\end{remark}

\begin{remark}[Pole hierarchy and conformal weight]
\label{rem:pole-hierarchy}
\index{r-matrix@$r$-matrix!pole hierarchy}
The maximal $r$-matrix pole order for the same-generator channel
follows the pattern $\max\text{pole}(r) = 2h - 1$: the Heisenberg
($h=1$) gives $z^{-1}$, the Virasoro ($h=2$) gives $z^{-3}$, and
$\cW_3$ ($h=3$) gives $z^{-5}$.
Systems whose OPE has at most a simple pole ($\beta\gamma$, $bc$)
produce an entirely regular $r$-matrix: the bar-complex propagator
extracts no singular collision residue. This is consistent with
shadow class~C (contact) for $\beta\gamma$ and the fact that the
leading bar interaction is quartic.
\end{remark}

\section{Free fermion bar complex}
\label{sec:fermion-bar-details}
\index{free fermion!bar complex!detailed computation}

This section gives a degree-by-degree computation of the bar complex
for the \emph{two-generator} free fermion system, the $bc$ system
at conformal weights $h_b = h_c = \tfrac{1}{2}$.
This supplements the single-generator analysis of
\S\ref{sec:free-fermion} and illuminates the key difference between
fermionic and bosonic bar complexes.

\subsection{Setup}

\begin{computation}[Two-generator free fermion OPE]
\label{comp:fermion-two-gen-ope}
\index{free fermion!two-generator}
Let $\mathcal{F}_2$ denote the free fermion chiral algebra generated
by two fermionic fields $\psi_1(z)$, $\psi_2(z)$ of conformal
weight $h = \tfrac{1}{2}$, with OPE:
\begin{equation}\label{eq:fermion-two-gen-ope}
\psi_i(z)\,\psi_j(w) = \frac{\delta_{ij}}{z - w} + \text{regular},
\qquad i, j \in \{1, 2\}.
\end{equation}
The $n$-th products are:
\begin{center}
\renewcommand{\arraystretch}{1.15}
\begin{tabular}{clll}
$n$ & $(\psi_i)_{(n)}\psi_j$ & Pole order & Type \\
\hline
$0$ & $\delta_{ij}\cdot |0\rangle$ & $(z-w)^{-1}$ & simple pole \\
$\geq 1$ & $0$ & none & regular
\end{tabular}
\end{center}
The augmentation ideal $\bar{V} = V/\bC|0\rangle$ has basis given by
normally ordered monomials in the modes
$\psi_{i,-n-1/2}$ for $n \geq 0$ (Neveu--Schwarz sector) or
$\psi_{i,-n}$ for $n \geq 1$ (Ramond sector).
The computations below use the NS sector throughout.
\end{computation}

\subsection{Degree-by-degree computation}

\begin{computation}[Degree 1]
\label{comp:fermion-deg1}
The degree-1 component is
\[
\B^1(\mathcal{F}_2) = \bar{V},
\]
spanned by all monomials in $\psi_{1,-n-1/2}$ and $\psi_{2,-n-1/2}$
(for $n \geq 0$) applied to the vacuum, modulo the vacuum itself.
At conformal weight $h$, one has
\[
\dim \B^1_h(\mathcal{F}_2) = p_{\mathrm{ferm}}(h;\,2)
\]
where $p_{\mathrm{ferm}}(h;\,r)$ counts the number of states of weight $h$
in the NS vacuum module of $r$ free fermions. The generating function is
\begin{equation}\label{eq:fermion-char}
\sum_{h \geq 0} p_{\mathrm{ferm}}(h;\,2)\,q^h
= \prod_{n=0}^{\infty}(1 + q^{n+1/2})^2 - 1
= 2q^{1/2} + 3q + 4q^{3/2} + 7q^2 + \cdots
\end{equation}
(the $-1$ removes the vacuum).
\end{computation}

\begin{computation}[Degree 2: fermionic OPE residues]
\label{comp:fermion-deg2}
\index{bar differential!free fermion}
Elements of $\B^2(\mathcal{F}_2) = \bar{V}^{\otimes 2}
\otimes \Omega^1(\overline{\Conf}_2)$ have the form
$[\phi_a \mid \phi_b] \otimes \eta_{12}$, where $\phi_a, \phi_b \in \bar{V}$.
The bar differential acts by the residue at the unique collision divisor $D_{12}$:
\begin{equation}\label{eq:fermion-bar-d2}
d\bigl([\phi_a \mid \phi_b] \otimes \eta_{12}\bigr)
= \sum_{n \geq 0} (\phi_a)_{(n)}\phi_b.
\end{equation}

\emph{Diagonal pairs.}
For the simplest generators $\psi_{i,-1/2}|0\rangle$ (the weight-$\frac{1}{2}$
states), the degree-2 differential gives:
\begin{align}
d\bigl([\psi_i \mid \psi_j] \otimes \eta_{12}\bigr)
&= (\psi_i)_{(0)}\psi_j = \delta_{ij}\cdot |0\rangle.
\label{eq:fermion-d2-simple}
\end{align}
In particular:
\begin{itemize}
\item $d([\psi_1 \mid \psi_1] \otimes \eta_{12}) = |0\rangle$
 \quad (maps to vacuum, nonzero!),
\item $d([\psi_2 \mid \psi_2] \otimes \eta_{12}) = |0\rangle$
 \quad (same),
\item $d([\psi_1 \mid \psi_2] \otimes \eta_{12}) = 0$
 \quad (off-diagonal: $\delta_{12} = 0$).
\end{itemize}

This is the key difference from the $\beta\gamma$ system: for the
bosonic generators $\beta$, $\gamma$ with $\beta(z)\gamma(w) = 1/(z-w)$,
one has $\beta_{(0)}\beta = 0$ and $\gamma_{(0)}\gamma = 0$
(the self-OPE is regular). The fermionic self-OPE
$\psi_i(z)\psi_i(w) = 1/(z-w)$ is \emph{singular}, producing a
nonzero differential on diagonal elements.
\end{computation}

\begin{computation}[Degree 2: Koszul signs]
\label{comp:fermion-deg2-signs}
\index{Koszul sign!fermion bar complex}
Since $\psi_i$ has odd parity, the bar complex carries a
sign from the desuspension $s^{-1}$:
the element $[s^{-1}\psi_i \mid s^{-1}\psi_j]$ in the
desuspended tensor coalgebra has parity $|s^{-1}\psi_i| = |\psi_i| - 1 = 0$
(even). Thus the symmetric coalgebra structure on $\B(\mathcal{F}_2)$
is the \emph{symmetric} (not exterior) coalgebra:
\[
\B(\mathcal{F}_2) \cong \mathrm{Sym}^c(s^{-1}\bar{V})
\]
as a graded coalgebra (before the differential).
Explicitly, the transposition acts by
\[
\tau\bigl([s^{-1}\psi_i \mid s^{-1}\psi_j]\bigr)
= (-1)^{|s^{-1}\psi_i|\,|s^{-1}\psi_j|}
 [s^{-1}\psi_j \mid s^{-1}\psi_i]
= +\,[s^{-1}\psi_j \mid s^{-1}\psi_i],
\]
i.e., the desuspended elements commute. This is dual to the
fact that $\mathcal{F}_2^! \cong \mathrm{Sym}^{\mathrm{ch}}(V^*)$
(the symmetric chiral algebra on $V^*$), consistent with
$\Lie^! = \Com$ at the operadic level.
\end{computation}

\begin{computation}[Degree 3: iterated residues]
\label{comp:fermion-deg3}
Degree-3 elements have the form
$[\phi_a \mid \phi_b \mid \phi_c] \otimes \omega$
with $\omega \in \Omega^2(\overline{\Conf}_3)$.
The differential $d\colon \B^3 \to \B^2$ acts by three collision residues.

\emph{The element $[\psi_1 \mid \psi_1 \mid \psi_2]
\otimes \eta_{12} \wedge \eta_{13}$.}

At $D_{12}$ ($z_1 = z_2$), with the $\psi_1$--$\psi_1$ OPE:
\begin{align*}
\Res_{D_{12}} &= (\psi_1)_{(0)}\psi_1 \otimes [\psi_2]
\otimes \eta_{13}\big|_{z_1 = z_2}
= |0\rangle \otimes [\psi_2] \otimes \eta_{23}
= [\psi_2] \otimes \eta_{23}.
\end{align*}

At $D_{13}$ ($z_1 = z_3$), with the $\psi_1$--$\psi_2$ OPE:
\begin{align*}
\Res_{D_{13}} &= (\psi_1)_{(0)}\psi_2 \otimes [\psi_1]
\otimes \eta_{12}\big|_{z_1 = z_3}
= 0
\end{align*}
(off-diagonal: $\delta_{12} = 0$).

At $D_{23}$ ($z_2 = z_3$), with the $\psi_1$--$\psi_2$ OPE:
\begin{align*}
\Res_{D_{23}} &= [\psi_1] \otimes (\psi_1)_{(0)}\psi_2
\otimes \eta_{12}\big|_{z_2 = z_3}
= 0
\end{align*}
(again off-diagonal).

Therefore:
\begin{equation}\label{eq:fermion-deg3-ex}
d\bigl([\psi_1 \mid \psi_1 \mid \psi_2]
\otimes \eta_{12} \wedge \eta_{13}\bigr)
= [\psi_2] \otimes \eta_{23}.
\end{equation}

By contrast, the element $[\psi_1 \mid \psi_2 \mid \psi_1]
\otimes \eta_{12} \wedge \eta_{13}$ has nonzero residues at
$D_{13}$ (from $\psi_1$--$\psi_1$) and $D_{23}$ (from $\psi_2$--$\psi_1$,
which vanishes). Collecting all contributions:
\begin{align*}
d\bigl([\psi_1 \mid \psi_2 \mid \psi_1]
\otimes \eta_{12} \wedge \eta_{13}\bigr)
&= 0 + [\psi_2] \otimes \eta_{12} + 0
= [\psi_2] \otimes \eta_{12}.
\end{align*}
(The $D_{13}$ residue uses $(\psi_1)_{(0)}\psi_1 = |0\rangle$,
collapsing two slots and leaving $[\psi_2]$.)
\end{computation}

\begin{computation}[Dimension table through degree 5]
\label{comp:fermion-dim-table}
\index{free fermion!dimension table}
The total dimension $\dim \B^n(\mathcal{F}_2)$ for the two-generator
free fermion in the NS sector, restricting to elements built
from the lowest modes $\psi_{i,-1/2}|0\rangle$ (weight $\tfrac{1}{2}$)
in the ordered tensor model used for the bar tables,
is given by the following table.

\begin{center}
\renewcommand{\arraystretch}{1.2}
\begin{tabular}{c|ccccc|l}
Degree $n$ & 1 & 2 & 3 & 4 & 5 & Generating principle \\
\hline
$\dim \B^n$ (lowest modes only) & 2 & 4 & 16 & 96 & 768
 & $2^n\,(n{-}1)!$ \\
$\dim \Omega^{n-1}(\overline{\Conf}_n)$ & 1 & 1 & 2 & 6 & 24
 & Poincar\'{e} of $\Conf_n(\bC)$ \\
$\operatorname{rk} d_n$ & n/a & $\leq 2$ & $\leq 4$ & $\leq 16$ & $\leq 96$
& target dimension bound
\end{tabular}
\end{center}

At all weights, the generating function for
$\sum_n \dim \B^n_h(\mathcal{F}_2) \cdot t^n \cdot q^h$ is given
by Proposition~\ref{prop:universal-dim-formula} below, with
$r = 2$ fermionic generators of weight $\frac{1}{2}$.
The symmetric-coalgebra description in
Proposition~\ref{prop:fermion-bar-symmetric} is the coinvariant
Com-shadow of the ordered tensor bar coalgebra; the ordered chain table
above remains the $E_1$ bar chain space.
\end{computation}

\begin{proposition}[Free fermion symmetric bar shadow: coalgebra structure;
\ClaimStatusProvedHere]
\label{prop:fermion-bar-symmetric}
\index{free fermion!symmetric coalgebra}
For the two-generator free fermion $\mathcal{F}_2$, the symmetric
coinvariant shadow of the ordered bar complex has the structure of a
\emph{symmetric} coalgebra on the desuspension:
\begin{equation}\label{eq:fermion-bar-symcoalg}
\B(\mathcal{F}_2)_{\Sigma} \cong \mathrm{Sym}^c(s^{-1}\bar{V})
\end{equation}
as a graded coalgebra (forgetting the differential), where
$\bar{V} = V/\bC|0\rangle$ is the augmentation ideal.
This is dual to the commutative algebra structure of the
Koszul dual $\mathcal{F}_2^! \cong
\mathrm{Sym}^{\mathrm{ch}}(V^*)$, under the operadic duality
$\Lie^! = \Com$.
\end{proposition}

\begin{proof}
The desuspension $s^{-1}$ shifts parity:
$|s^{-1}\phi| = |\phi| - 1$. Since all generators $\psi_i$ of
$\mathcal{F}_2$ are odd ($|\psi_i| = 1$), their desuspensions
$s^{-1}\psi_i$ are even ($|s^{-1}\psi_i| = 0$). The ordered bar
coalgebra is
\[
\B_{E_1}(\mathcal{F}_2)=T^c(s^{-1}\bar V)
\]
with deconcatenation coproduct. Passing to the Com-operadic
$\Sigma_n$-coinvariant shadow sends tensor words to unordered monomials
in the even generators $s^{-1}\psi_i$, hence gives
$\mathrm{Sym}^c(s^{-1}\bar V)$. No quotient description is used for the
ordered $E_1$ bar coalgebra.

Duality: $\mathrm{Sym}^c(s^{-1}\bar{V})^*
\cong \mathrm{Sym}(s\bar{V}^*) = \mathrm{Sym}^{\mathrm{ch}}(V^*)$,
which is a commutative (chiral) algebra. Since $\mathcal{F}_2$
is the chiral Lie algebra $\Lambda^{\mathrm{ch}}(V)$ on the
two-dimensional odd space $V = \bC\psi_1 \oplus \bC\psi_2$,
and $\Lie^! = \Com$, the identification
$\mathcal{F}_2^! \cong \mathrm{Sym}^{\mathrm{ch}}(V^*)$
follows from operadic Koszul duality
(Theorem~\ref{thm:single-fermion-boson-duality}).
\end{proof}


\section{\texorpdfstring{Lattice VOA bar complex: $E_8$ root lattice}{Lattice VOA bar complex: E8 root lattice}}
\label{sec:E8-lattice-bar}
\index{E8 root lattice@$E_8$ root lattice!bar complex}
\index{lattice vertex algebra!bar complex!$E_8$}

The $E_8$ analysis of
Proposition~\ref{prop:lattice:bar-E8} with explicit degree-by-degree
computations of the bar differential tracks the coupling between
the Heisenberg (Cartan) sector and the vertex operator (root) sectors.

\subsection{Generator content}

\begin{computation}[\texorpdfstring{$E_8$}{E_8} lattice VOA generators]
\label{comp:E8-generators}
\index{E8 root lattice@$E_8$ root lattice!generators}
The $E_8$ lattice vertex algebra $V_{E_8}$ at level $k = 1$ has the
following strong generators:
\begin{itemize}
\item \emph{Cartan currents}: $J^i(z)$ for $i = 1, \ldots, 8$
 (the rank of $E_8$), each of conformal weight~$1$.
\item \emph{Vertex operators}: $e^\alpha(z)$ for each root
 $\alpha \in \Phi(E_8)$, with $|\Phi(E_8)| = 240$.
 Each $e^\alpha$ has conformal weight
 $h_\alpha = \langle \alpha, \alpha \rangle / 2 = 1$
 (since all roots of $E_8$ satisfy $\langle \alpha, \alpha \rangle = 2$).
\end{itemize}
The total generator count at weight~$1$ is
$8 + 240 = 248 = \dim(E_8)$, consistent with the Frenkel--Kac--Segal
theorem (Theorem~\ref{thm:lattice:frenkel-kac}):
$V_{E_8} \cong L(\widehat{\mathfrak e}_8,1)$ as vertex algebras,
where $L(\widehat{\mathfrak e}_8,1)$ is the simple level-$1$
quotient of the universal affine vacuum module
$V_1(\mathfrak e_8)$.

\emph{Root system data.}
\begin{center}
\renewcommand{\arraystretch}{1.15}
\begin{tabular}{cccccc}
$\dim(\mathfrak{e}_8)$ & $\operatorname{rank}$ & $|\Phi|$ & $h$
 & $h^\vee$ & $c = 248k/(k+30)$ at $k=1$ \\
\hline
248 & 8 & 240 & 30 & 30 & $248/31$
\end{tabular}
\end{center}
\end{computation}

\subsection{OPE structure and bar differential}

\begin{computation}[\texorpdfstring{$E_8$}{E_8} bar differential: degree 2]
\label{comp:E8-bar-deg2}
\index{bar differential!$E_8$}
The degree-2 bar differential $d\colon \B^2(V_{E_8}) \to \B^1(V_{E_8})$
encodes the full OPE data of $V_{E_8}$. The $248^2 = 61{,}504$
potential degree-2 elements $[\phi_a \mid \phi_b] \otimes \eta_{12}$
(where $\phi_a, \phi_b$ are weight-1 generators) decompose into
three types:

\emph{Type I: Cartan--Cartan ($J^i \otimes J^j$, $8^2 = 64$ elements).}
\begin{align}
d\bigl([J^i \mid J^j] \otimes \eta_{12}\bigr)
&= (J^i)_{(1)}J^j \cdot |0\rangle + (J^i)_{(0)}J^j \notag \\
&= k\,\delta^{ij} \cdot |0\rangle + 0
\label{eq:E8-cartan-cartan}
\end{align}
since the Heisenberg OPE $J^i(z) J^j(w) = k\,\delta^{ij}/(z-w)^2$
has only a double pole (no simple pole for commuting Cartan generators).
At $k = 1$, the $8$ diagonal elements $[J^i \mid J^i]$ each
map to the vacuum, producing \emph{curvature}.

\emph{Type II: Cartan--root and root--Cartan
($J^i \otimes e^\alpha$ and $e^\alpha \otimes J^i$,
$2\cdot 8 \cdot 240 = 3840$ ordered elements).}
\begin{align}
d\bigl([J^i \mid e^\alpha] \otimes \eta_{12}\bigr)
&= (J^i)_{(0)}e^\alpha = \alpha_i \, e^\alpha
\label{eq:E8-cartan-root}
\end{align}
where $\alpha_i = \langle \alpha, \epsilon_i \rangle$ is the $i$-th
component of $\alpha$ in the Cartan basis. This is a simple-pole
contribution mapping to $\B^1$. Similarly,
$d([e^\alpha \mid J^i] \otimes \eta_{12}) = -\alpha_i\, e^\alpha$
by skew-symmetry.

\emph{Type III: Root--root ($e^\alpha \otimes e^\beta$,
$240^2 = 57{,}600$ elements).}
The differential depends on
$\langle \alpha, \beta \rangle \in \{-2, -1, 0, 1, 2\}$:
\begin{center}
\renewcommand{\arraystretch}{1.2}
\begin{tabular}{cl}
$\langle \alpha, \beta \rangle$ & $d([e^\alpha \mid e^\beta] \otimes \eta_{12})$ \\
\hline
$-2$ \;($\beta = -\alpha$) &
 $J_\alpha + k \cdot |0\rangle$ \\
$-1$ \;($\alpha + \beta \in \Phi$) &
 $\varepsilon(\alpha,\beta)\, e^{\alpha+\beta}$ \\
$0$ \;($\alpha + \beta \notin \Phi \cup \{0\}$) & $0$ \\
$+1$ & $0$ \;(no singular OPE) \\
$+2$ \;($\beta = \alpha$) & $0$ \;($e^\alpha e^\alpha$ is regular)
\end{tabular}
\end{center}
Here $J_\alpha = \sum_i \alpha_i J^i$ is the Cartan current
in the $\alpha$-direction, and $\varepsilon(\alpha,\beta) \in \{\pm 1\}$
is the FLM $2$-cocycle (Definition~\ref{def:lattice:2-cocycle}).

\emph{Counting nonzero differentials.}
The number of pairs $(\alpha, \beta) \in \Phi \times \Phi$ with
$\langle \alpha, \beta \rangle = -1$ equals $240 \times 56 = 13{,}440$
(each root has exactly $56$ neighbors at inner product $-1$ in $E_8$,
cf.\ Proposition~\ref{prop:lattice:bar-E8}).
The number with $\langle \alpha, \beta \rangle = -2$
equals $240$ (the ordered pairs $(\alpha, -\alpha)$). The ordered
Cartan--root channels contribute $3840$ possible simple-pole channels,
with coefficient $\alpha_i$ in the chosen Cartan basis. Thus the
singular channels outside the pure Cartan--Cartan block number
\[
3840 + 13{,}440 + 240 = 17{,}520.
\]
If one also counts the eight diagonal Cartan--Cartan curvature maps
\eqref{eq:E8-cartan-cartan}, the full ordered degree-$2$ list has
$17{,}528$ nonzero channel entries in a generic orthonormal Cartan
basis. In a coordinate basis adapted to the $E_8$ lattice, the
formula \eqref{eq:E8-cartan-root} is the invariant datum and those
entries with $\alpha_i=0$ are simply zero.
\end{computation}

\begin{computation}[Curvature of the \texorpdfstring{$E_8$}{E_8} bar complex]
\label{comp:E8-curvature}
\index{curvature!$E_8$ lattice}
There are two $E_8$ rows, and the bar table keeps them separate.
For the universal affine algebra $V_1(\mathfrak e_8)$ before the
Frenkel--Kac--Segal quotient, the tensor curvature is
\begin{equation}\label{eq:E8-curvature}
m_0 = \frac{k + h^\vee}{2h^\vee}\,\kappa
= \frac{1 + 30}{2 \cdot 30}\,\kappa
= \frac{31}{60}\,\kappa,
\end{equation}
where $\kappa$ is the Killing form, viewed as an element of
$\mathrm{Sym}^2(\mathfrak{h}^*) \subset \B^0$.
Its scalar modular characteristic is
\[
\kappa(V_1(\mathfrak e_8))
=\frac{248(1+30)}{2\cdot30}
=\frac{1922}{15}
=\frac{62(31)}{15}.
\]
The simple lattice VOA $V_{E_8}\simeq L_1(\mathfrak e_8)$ is a
different object. The FKS quotient identifies the root currents with
composites of the eight Heisenberg generators, and the lattice row in
the standard census has
\[
\kappa(V_{E_8})=\operatorname{rank}(E_8)=8.
\]
Thus \eqref{eq:E8-curvature} belongs to the universal affine
same-family shadow. The lattice bar-complex entries below use the
rank-$8$ lattice value and the unimodular self-duality of $E_8$.
\end{computation}

\subsection{Koszul acyclicity and self-duality}

\begin{proposition}[\texorpdfstring{$E_8$}{E_8} affine pre-quotient Koszul acyclicity;
\ClaimStatusProvedHere]
\label{prop:E8-koszul-acyclic}
\index{Koszul!acyclicity!$E_8$}
For the universal affine algebra $V_1(\mathfrak e_8)$, the twisted
bar complex
$\B(V_1(\mathfrak e_8)) \otimes_\tau V_1(\mathfrak e_8)$ is acyclic:
\[
H_n\bigl(\B(V_1(\mathfrak e_8)) \otimes_\tau V_1(\mathfrak e_8)\bigr) = \begin{cases}
\bC & n = 0, \\
0 & n > 0.
\end{cases}
\]
The statement is not a lattice self-duality theorem for $V_{E_8}$.
\end{proposition}

\begin{proof}
The argument parallels Proposition~\ref{prop:km-generic-acyclicity}.
Introduce the PBW filtration $F_p$ by generator count. The
associated graded is $\mathrm{gr}_F(\B(V_1(\mathfrak e_8)) \otimes_\tau V_1(\mathfrak e_8))
\cong B(U(\mathfrak{e}_8[t^{-1}]t^{-1}))\otimes_{\tau_0}
U(\mathfrak{e}_8[t^{-1}]t^{-1})$.
This associated graded complex is the augmented bar resolution of
the enveloping algebra of the negative loop algebra, hence has
homology $\bC$ in degree~$0$ and no positive homology. The spectral
sequence collapses on the augmentation line.
\end{proof}

\begin{computation}[Affine level shift versus lattice self-duality]
\label{comp:E8-koszul-dual}
\index{E8 root lattice@$E_8$ root lattice!Koszul dual}
For the universal affine algebra $V_k(\mathfrak e_8)$, the scalar
Verdier--Koszul lane carries the Feigin--Frenkel level shift
(Corollary~\ref{cor:level-shifting-part1})
\[
k' = -k - 2h^\vee = -1 - 60 = -61.
\]
This shifted-level object is not the lattice self-dual VOA
$V_{E_8}$ itself. It is the affine same-family shadow whose scalar
invariants match the Verdier--Koszul dual lane of the universal
affine algebra.

The central charge of the shifted affine shadow is:
\[
c' = \frac{248 \cdot (-61)}{-61 + 30} = \frac{248 \cdot 61}{31} = \frac{15128}{31},
\]
and the complementarity sum is
$c + c' = \frac{248}{31} + \frac{15128}{31} = \frac{15376}{31} = 496$.
Equivalently,
$\frac{\dim(\mathfrak{e}_8) \cdot k}{k+h^\vee} +
\frac{\dim(\mathfrak{e}_8) \cdot k'}{k'+h^\vee}
= 248\bigl(\tfrac{1}{31} + \tfrac{61}{31}\bigr)
= 248 \cdot \tfrac{62}{31} = 496$,
consistent with the general formula
$c + c' = 2\dim(\mathfrak{g})$ when $k' = -k - 2h^\vee$
(which gives $\frac{d \cdot k}{k + h^\vee} + \frac{d(-k-2h^\vee)}{-k-h^\vee}
= d\bigl(\frac{k}{k+h^\vee} + \frac{k+2h^\vee}{k+h^\vee}\bigr)
= d \cdot \frac{2k + 2h^\vee}{k+h^\vee} = 2d$).

\emph{Lattice self-duality.}
As a lattice vertex algebra, $V_{E_8}$ is self-dual by
Theorem~\ref{thm:lattice:unimodular-self-dual} (since $E_8$ is
even unimodular). This is a different statement from the affine
level shift above: lattice self-duality acts on the lattice
coalgebra/cocycle data, while $k\leftrightarrow -k-2h^\vee$ acts on
the universal affine same-family shadow. Its scalar row is
$\kappa(V_{E_8})=8$, not $\kappa(V_1(\mathfrak e_8))=1922/15$.
\end{computation}


\section{Virasoro bar complex through degree 5}
\label{sec:virasoro-deg45}
\index{Virasoro algebra!bar complex!degrees 4 and 5}

The Virasoro bar complex computation from degree~3
(Computation~\ref{comp:virasoro-degree3}, Proposition~\ref{prop:arnold-virasoro-deg3})
extends to degrees~4 and~5 as follows.

\subsection{Form spaces}

\begin{computation}[Poincar\'{e} polynomials of \texorpdfstring{$\Conf_n(\bC)$}{Conf_n(C)}]
\label{comp:poincare-config}
\index{configuration space!Poincar\'e polynomial}
The Poincar\'{e} polynomial of the ordered configuration space
$\Conf_n(\bC) = \{(z_1, \ldots, z_n) \in \bC^n : z_i \neq z_j\}$
is (Arnol'd, \cite{Arnold69}):
\begin{equation}\label{eq:poincare-config}
P_t(\Conf_n(\bC)) = \prod_{j=1}^{n-1}(1 + jt).
\end{equation}
The relevant values are:
\begin{center}
\renewcommand{\arraystretch}{1.15}
{\small
\begin{tabular}{c|l|cccccccc}
$n$ & $P_t$ & $b_0$ & $b_1$ & $b_2$ & $b_3$ & $b_4$ & $b_5$ & $b_6$ & $b_7$ \\
\hline
2 & $1 + t$ & 1 & 1 \\
3 & $(1+t)(1+2t)$ & 1 & 3 & 2 \\
4 & $(1+t)\cdots(1+3t)$ & 1 & 6 & 11 & 6 \\
5 & $(1+t)\cdots(1+4t)$ & 1 & 10 & 35 & 50 & 24 \\
6 & $(1+t)\cdots(1+5t)$ & 1 & 15 & 85 & 225 & 274 & 120 \\
7 & $(1+t)\cdots(1+6t)$ & 1 & 21 & 175 & 735 & 1624 & 1764 & 720 \\
8 & $(1+t)\cdots(1+7t)$ & 1 & 28 & 322 & 1960 & 6769 & 13132 & 13068 & 5040
\end{tabular}
}
\end{center}
The top Betti number is $b_{n-1} = (n-1)!$ in each case.
For the bar complex $\B^n(\mathrm{Vir}_c) = \bar{V}^{\otimes n}
\otimes \Omega^{n-1}(\overline{\Conf}_n)$, the form factor
contributes $\dim \Omega^{n-1}(\overline{\Conf}_n) = (n-1)!$
when all insertions are identical.
\end{computation}

\subsection{Degree 4}

\begin{computation}[Degree-4 Virasoro bar complex]
\label{comp:virasoro-deg4-full}
\index{Virasoro algebra!bar complex!degree 4}

\emph{Elements.}
$\B^4(\mathrm{Vir}_c) = \bar{V}^{\otimes 4}
\otimes \Omega^3(\overline{\Conf}_4)$, where
$\dim \Omega^3(\overline{\Conf}_4) = 6$
(the top Betti number of $\Conf_4(\bC)$).

A basis for $\Omega^3(\overline{\Conf}_4)$:
\begin{equation}\label{eq:form-basis-4}
\eta_{12} \wedge \eta_{13} \wedge \eta_{14}, \quad
\eta_{12} \wedge \eta_{13} \wedge \eta_{24}, \quad
\eta_{12} \wedge \eta_{13} \wedge \eta_{34}, \quad
\eta_{12} \wedge \eta_{23} \wedge \eta_{34}, \quad
\eta_{12} \wedge \eta_{24} \wedge \eta_{34}, \quad
\eta_{13} \wedge \eta_{23} \wedge \eta_{34}.
\end{equation}
(These $6$ forms among $\binom{6}{3} = 20$ possible triple wedges
are the survivors after applying the Arnold relations
$\eta_{ij} \wedge \eta_{jk} + \eta_{jk} \wedge \eta_{ki}
+ \eta_{ki} \wedge \eta_{ij} = 0$ for each triple $\{i,j,k\}$.)

\emph{The element $[T\mid T\mid T\mid T]
\otimes \eta_{12} \wedge \eta_{13} \wedge \eta_{14}$.}
The differential $d\colon \B^4 \to \B^3$ acts by six collision
residues, one for each divisor $D_{ij}$ with $1 \leq i < j \leq 4$.

\emph{Collision $D_{12}$:}
\begin{align*}
\Res_{D_{12}} &= \bigl(T_{(3)}T \cdot |0\rangle + T_{(1)}T
 + T_{(0)}T\bigr) \otimes [T \mid T]
 \otimes \eta_{13} \wedge \eta_{14}\big|_{z_1 = z_2} \\
&= \frac{c}{2}\,[T \mid T] \otimes \eta_{13} \wedge \eta_{14}
 + 2\,[T \mid T \mid T] \otimes \eta_{13} \wedge \eta_{14}
 + [\partial T \mid T \mid T] \otimes \eta_{13} \wedge \eta_{14}.
\end{align*}

\emph{Collision $D_{13}$:}
After $z_1 = z_3$, the surviving form indices relabel.
The $\eta_{12}$ factor survives as $\eta_{12}$, and $\eta_{14}$
becomes $\eta_{14}$:
\begin{align*}
\Res_{D_{13}} &= \frac{c}{2}\,[T \mid T] \otimes \eta_{12} \wedge \eta_{14}
 + 2\,[T \mid T \mid T] \otimes \eta_{12} \wedge \eta_{14}
 + [T \mid \partial T \mid T] \otimes \eta_{12} \wedge \eta_{14}.
\end{align*}
(The slot structure reflects that positions $1$ and $3$ merged,
with position~$2$ and position~$4$ surviving.)

The remaining four collisions $D_{14}$, $D_{23}$, $D_{24}$, $D_{34}$
are analogous. Only the vacuum contributions enter the curvature
propagation from degree~2.

\emph{Vacuum contributions.}
Write
\[
V_{ij} := \frac{c}{2} [T \mid T] \otimes \eta_{ij},
\]
where $\eta_{ij}$ is the $2$-form obtained by restricting to
$D_{ij}$ and deleting the factors involving the collided indices.

Summing over all six divisors gives
\begin{align}
\sum_{1 \leq i < j \leq 4} V_{ij}
&= \frac{c}{2} \sum_{i < j} [T \mid T] \otimes \eta_{ij}.
\label{eq:vir-deg4-vacuum}
\end{align}

Each surviving $2$-form $\eta_{ij}$ lives in $\Omega^2(\overline{\Conf}_3)$.
There is no linear Arnold identity of the form
$\sum\eta_{ij}=0$. The Arnold identities in
$\Omega^\bullet(\Conf_4)$ are quadratic wedge identities. Therefore
the sum~\eqref{eq:vir-deg4-vacuum} is not killed by configuration
space alone. It is the degree-$4$ curvature channel of the Virasoro
bar complex and must be paired with the $m_0=c/2$ term in the curved
$A_\infty$ identity. Degree~$4$ does not provide an independent
vanishing theorem for Virasoro vacuum leakage.
\end{computation}

\subsection{Degree 5}

\begin{computation}[Degree-5 bar complex for \texorpdfstring{$\mathrm{Vir}_c$}{Vir_c}]
\label{comp:virasoro-deg5}
\index{Virasoro algebra!bar complex!degree 5}

\emph{Form space.}
$\dim \Omega^4(\overline{\Conf}_5) = 24 = 4!$, consistent with
$P_t(\Conf_5(\bC)) = 1 + 10t + 35t^2 + 50t^3 + 24t^4$.

\emph{Differential structure.}
The differential $d\colon \B^5 \to \B^4$ involves $\binom{5}{2} = 10$
collision divisors. For the all-$T$ element
$[T|T|T|T|T] \otimes \omega_{\mathrm{top}}$
(where $\omega_{\mathrm{top}} \in \Omega^4(\overline{\Conf}_5)$),
each collision $D_{ij}$ produces:
\begin{itemize}
\item A vacuum term $\frac{c}{2}\,[T|T|T] \otimes \omega^{(3)}_{ij}$
 (dropping to degree~3 in the augmentation ideal part),
\item A conformal weight term
 $2\,[T|T|T|T] \otimes \omega^{(4)}_{ij}$,
\item A translation term
 $[\ldots|\partial T|\ldots|T|T] \otimes \omega^{(4)}_{ij}$,
\end{itemize}
where $\omega^{(k)}_{ij} \in \Omega^{k-1}(\overline{\Conf}_k)$ is the
form obtained by restricting $\omega_{\mathrm{top}}$ to the stratum
$z_i = z_j$ and projecting to the configuration of $k$ surviving points.

The vacuum contributions do not vanish by a fixed-slot linear Arnold
sum. The ten boundary terms are constrained by the Orlik--Solomon
quadratic relations and by the curved equation $D^2=[m_0,-]$; they are
not zero as ordinary residue terms. Degree~$5$ therefore gives another
instance of Virasoro curvature propagation, not a cancellation
mechanism in every degree.
\end{computation}

\subsection{Graded dimension table}

\begin{computation}[Virasoro bar complex dimensions through degree 5]
\label{comp:virasoro-dim-table}
\index{Virasoro algebra!bar complex!dimension table}
Let $\B^n_h(\mathrm{Vir}_c)$ denote the subspace of bar degree~$n$
and conformal weight~$h$. The conformal weight of an element
$[L_{-a_1}| \cdots |L_{-a_n}] \otimes \omega$ is
$h = a_1 + \cdots + a_n$ (with $a_i \geq 2$), so
$h \geq 2n$. For weight-$h$ elements involving only the
generator $T = L_{-2}|0\rangle$ and its descendants, the count is:
\begin{center}
\renewcommand{\arraystretch}{1.15}
{\small
\begin{tabular}{c|ccccccccccccccc}
$n \backslash h$ & 2 & 3 & 4 & 5 & 6 & 7 & 8 & 9 & 10 & 11 & 12 & 13 & 14 & 15 & 16 \\
\hline
1 & 1 & 1 & 2 & 2 & 4 & 4 & 7 & 8 & 12 & 14 & 21 & 24 & 34 & 41 & 55 \\
2 & & & 1 & 2 & 5 & 8 & 16 & 24 & 42 & 62 & 100 & 144 & 222 & 314 & 465 \\
3 & & & & & 2 & 6 & 18 & 38 & 84 & 156 & 298 & 516 & 906 & 1492 & 2478 \\
4 & & & & & & & 6 & 24 & 84 & 216 & 534 & 1152 & 2424 & 4728 & 9048 \\
5 & & & & & & & & & 24 & 120 & 480 & 1440 & 3960 & 9624 & 22200 \\
6 & & & & & & & & & & & 120 & 720 & 3240 & 11040 & 33480 \\
7 & & & & & & & & & & & & & 720 & 5040 & 25200 \\
8 & & & & & & & & & & & & & & & 5040
\end{tabular}
}
\end{center}
The row $n = 1$ lists $\dim \bar{V}_h = p_{\geq 2}(h)$
(partitions into parts $\geq 2$). The row $n = 2$ lists
$\sum_{a+b=h} p_{\geq 2}(a) \cdot p_{\geq 2}(b)$
(the form space $\Omega^1(\overline{\Conf}_2)$ is
$1$-dimensional).
The row $n = k$ has entries
\[
\Bigl(\sum_{a_1 + \cdots + a_k = h}
\prod_i p_{\geq 2}(a_i)\Bigr) \cdot (k-1)!
\]
for $h \geq 2k$, where the factor
$(k-1)! = \dim \Omega^{k-1}(\overline{\Conf}_k)$
counts the top-degree forms.

In particular:
\begin{itemize}
\item At $h = 2$: $\dim \B^1_2 = 1$ (the single state $T$).
\item At $h = 4$: $\dim \B^2_4 = 1$ (the element $[T|T] \otimes \eta_{12}$).
\item At $h = 6$: $\dim \B^3_6 = 2$ (two independent $3$-forms on $\Conf_3$).
\item At $h = 8$: $\dim \B^4_8 = 6$ (six independent $3$-forms on $\Conf_4$).
\item At $h = 10$: $\dim \B^5_{10} = 24$ (twenty-four $4$-forms on $\Conf_5$).
\end{itemize}
\end{computation}

\begin{proposition}[Virasoro bar cohomology and Koszul-dual Hilbert function; \ClaimStatusConditional]
\label{prop:virasoro-koszul-acyclic}
\index{Virasoro algebra!bar cohomology}
\index{Koszul!Virasoro}
For the universal Virasoro algebra $\mathrm{Vir}_c$ at generic
$c$ outside the minimal-model/Kac-determinant divisor, the PBW
filtration on $\B(\mathrm{Vir}_c)$ yields a spectral sequence whose
$E_2$-collapse computes the dual coalgebra:
\begin{equation}\label{eq:virasoro-koszul-acyclic}
H^n(\B(\mathrm{Vir}_c)) \cong \bigl((\mathrm{Vir}_c)^i\bigr)_n
\end{equation}
where $(\mathrm{Vir}_c)^i := H^*(\B(\mathrm{Vir}_c))$ is the dual
coalgebra, and $\mathrm{Vir}_c^! = (((\mathrm{Vir}_c)^i)^\vee)$ is the
dual algebra. The displayed dimensions are therefore the Hilbert
function of $A^!$, not a statement about
$Z^{\mathrm{der}}_{\mathrm{ch}}(\mathrm{Vir}_c)$ and not an
all-central-charge theorem.
The Koszul-dual Hilbert dimensions are as listed in
Table~\textup{\ref{tab:bar-dimensions}}:
$\dim(\mathrm{Vir}_c^!)_n =
1, 2, 5, 12, 30, 76, 196, 512, 1353, 3610$ for
$n = 1, \ldots, 10$,
reflecting the associahedron structure of
$\overline{M}_{0,n+2}$.
\end{proposition}

\begin{proof}
Filter $\B(\mathrm{Vir}_c)$ by the PBW degree (total number
of mode applications). The associated graded bar complex is
$\mathrm{gr}\,\B(\mathrm{Vir}_c)
\cong \B(\mathrm{Sym}(L_{-2}, L_{-3}, \ldots))$,
the bar construction of the polynomial algebra on generators
$L_{-n}$ for $n \geq 2$.

The bar complex of $\mathrm{Sym}(V)$ for $V = \bigoplus_{n \geq 2}
\bC \cdot L_{-n}$ is quasi-isomorphic to the Koszul dual
$\mathrm{Sym}(V)^! = \Lambda(V^*)$ (the exterior coalgebra on dual
generators). In particular, $H^n(\B(\mathrm{Sym}(V)))$ is nonzero
in positive degrees: it computes the exterior powers $\Lambda^n(V^*)$.

The PBW filtration gives a spectral sequence with
$E_1^{p,q} = H_p(\B(\mathrm{gr}\,\mathrm{Vir}_c))_q
\cong \Lambda^p(V^*)_q$.
The spectral sequence collapses at $E_2$ by the chiral
Koszulness of generic $\mathrm{Vir}_c$
(Theorem~\ref{thm:virasoro-chiral-koszul}, proved via
the PBW criterion of Theorem~\ref{thm:pbw-koszulness-criterion}),
and the bar cohomology is the dual coalgebra
$(\mathrm{Vir}_c)^i$, concentrated on the diagonal of the bigrading.
The dimensions
\[
1,\; 2,\; 5,\; 12,\; 30,\; 76,\; 196,\; 512,\; 1353,\; 3610
\]
(the Motzkin differences $M(n{+}1) - M(n)$) encode the associahedron
combinatorics of $\overline{M}_{0,n+2}$
(Theorem~\ref{thm:motzkin-path-model}).
\end{proof}


\section{Universal dimension formulas}
\label{sec:universal-dim-formulas}
\index{bar complex!universal dimension formula}
\index{dimension formula!bar complex}

The preceding computations exhibit a common pattern: the graded
dimensions of the bar complex are determined by the
conformal weights and statistics of the strong generators.
The following proposition gives the resulting dimension formula.

\begin{proposition}[Free PBW bar dimension envelope;
\ClaimStatusProvedHere]
\label{prop:universal-dim-formula}
\index{PBW filtration!bar complex}
Let $A$ be a vertex algebra with strong generators
$\phi_1, \ldots, \phi_r$ of conformal weights $h_1, \ldots, h_r$
and parities $\epsilon_i \in \{0, 1\}$, and suppose that the PBW
associated graded is freely generated by the corresponding modes.
Then the generating function for the dimensions of the free PBW bar
envelope $\B^{n,\mathrm{PBW}}_h(A)$ is:
\begin{align}
\chi_{\bar V}(q)
&=
\prod_{i=1}^{r}
 \prod_{m \geq h_i}
 \frac{(1 + q^m)^{\epsilon_i}}{(1 - q^m)^{1-\epsilon_i}}
-1, \label{eq:universal-state-char}\\
\sum_{n \geq 1,\, h \geq 0} \dim \B^{n,\mathrm{PBW}}_h(A) \cdot q^h \, t^n
&=
\sum_{n\geq1} (n-1)!\,t^n\,\chi_{\bar V}(q)^n .
\label{eq:universal-dim}
\end{align}
The factor $(n-1)!$ is the top Orlik--Solomon form dimension
$\dim\Omega^{n-1}(\Conf_n(\mathbb C))$. The variable~$t$ records bar
tensor slots, not PBW occupation number inside a single state.
For a general strongly generated vertex algebra the same expression is
only the free PBW envelope; null vectors and normally ordered
relations quotient it.
\end{proposition}

\begin{proof}
The bar complex $\B^n(A)$ is built from the augmentation ideal
$\bar{V} = V / \bC|0\rangle$. As a graded vector space (ignoring
the bar differential), $\B^n(A) = \bar{V}^{\otimes n} \otimes
\Omega^{n-1}(\overline{\Conf}_n)$. For the purpose of
counting dimensions, the PBW filtration is used: $\bar{V}$ is
spanned by normally ordered monomials in the modes
$(\phi_i)_{-m}$ for $m \geq h_i$ applied to the vacuum.
Under the free PBW hypothesis the character of $\bar{V}$ is
\[
\chi_{\bar{V}}(q) = \prod_{i:\,\epsilon_i = 0}
 \prod_{m \geq h_i} \frac{1}{1 - q^m}
\;\cdot\; \prod_{i:\,\epsilon_i = 1}
 \prod_{m \geq h_i} (1 + q^m) \;-\; 1.
\]
The $n$-fold tensor product $\bar{V}^{\otimes n}$ at weight~$h$
has dimension $\sum_{h_1 + \cdots + h_n = h}
\prod_{j=1}^n \dim \bar{V}_{h_j}$ (using the ordered tensor product).
The form space $\Omega^{n-1}(\overline{\Conf}_n)$ contributes
an additional factor equal to the $(n-1)$-th Betti number of
$\Conf_n(\bC)$.

At the level of the free PBW envelope, each mode
$(\phi_i)_{-m}|0\rangle$ contributes independently to the
\emph{state} character. A bosonic mode contributes $(1-q^m)^{-1}$,
and a fermionic mode contributes $(1+q^m)$. The $n$-fold ordered
tensor product contributes $\chi_{\bar V}(q)^n$, and the top
configuration-form factor contributes $(n-1)!$. This gives
\eqref{eq:universal-dim}. The $-1$ in
\eqref{eq:universal-state-char} removes the vacuum from $\bar V$.
If the associated graded has relations, the actual bar dimensions are
obtained by quotienting this envelope by the induced relation ideal.
\end{proof}

\begin{computation}[Heisenberg dimension formula]\label{comp:dim-formula-heisenberg}
The Heisenberg algebra $\mathcal{H}$ has $r = 1$ bosonic generator
of weight $h_1 = 1$. Formula~\eqref{eq:universal-state-char} gives
$\chi_{\bar V}(q)=\prod_{m\ge1}(1-q^m)^{-1}-1$, hence
Formula~\eqref{eq:universal-dim} gives:
\[
\sum_{n,h} \dim \B^n_h(\mathcal{H}) \cdot q^h t^n
= \sum_{n\ge1}(n-1)!t^n
\left(\prod_{m\ge1}\frac1{1-q^m}-1\right)^n .
\]
At bar degree $n = 1$: $\dim \B^1_h = 1$ for all $h \geq 1$
(the modes $a_{-h}|0\rangle$). At $n = 2$ the coefficient is that of
$\chi_{\bar V}(q)^2$, since $\Omega^1(\Conf_2)$ is one-dimensional.
This matches
Computation~\ref{comp:heisenberg-deg3-full}.
\end{computation}

\begin{computation}[Virasoro dimension formula]\label{comp:dim-formula-virasoro}
The Virasoro has $r = 1$ bosonic generator of weight $h_1 = 2$:
\[
\sum_{n,h} \dim \B^n_h \cdot q^h t^n
= \sum_{n\ge1}(n-1)!t^n
\left(\prod_{m\ge2}\frac1{1-q^m}-1\right)^n .
\]
At $n = 1$:
$\dim \B^1_h = p_{\geq 2}(h)$, the number of partitions into parts
$\geq 2$ (matching Computation~\ref{comp:virasoro-vacuum}).
At $n = 2$, $h = 4$: a single element $[T|T]$,
giving $\dim \B^2_4 = 1$.
At $n = 3$, $h = 6$: two elements, giving $\dim \B^3_6 = 2$.
These agree with the table of Computation~\ref{comp:virasoro-dim-table}.
\end{computation}

\begin{computation}[\texorpdfstring{$\mathcal{W}_3$}{W_3} dimension formula]\label{comp:dim-formula-w3}
\index{W3 algebra@$\mathcal{W}_3$!bar complex dimension}
The $\mathcal{W}_3$ algebra has $r = 2$ bosonic generators of weights
$h_1 = 2$ (stress tensor $T$), $h_2 = 3$ ($W$-current $W$):
\[
\sum_{n,h} \dim \B^n_h \cdot q^h t^n
= \sum_{n\ge1}(n-1)!t^n
\left(
\prod_{m \geq 2} \frac{1}{1 - q^m}
\cdot \prod_{m \geq 3} \frac{1}{1 - q^m} - 1
\right)^n .
\]
At $n = 1$: $\dim \bar{V}_h$ counts partitions of $h$ into parts
$\geq 2$ from the $T$-sector and parts $\geq 3$ from the $W$-sector.
Explicitly: $\dim \bar{V}_2 = 1$ ($T$), $\dim \bar{V}_3 = 2$
($\partial T, W$), $\dim \bar{V}_4 = 4$
($\partial^2 T, \normord{TT}, \partial W, L_{-2}W$).
At $n = 2$, $h = 4$: a single element $[T|T]$ (the only way to
partition weight $4$ into two parts $\geq 2$ from either sector),
giving $\dim \B^2_4 = 1$.
At $n = 2$, $h = 5$: the elements $[T|W]$ and $[W|T]$ give
$\dim \B^2_5 = 2$ (in addition to $[\partial T | T]$ and $[T|\partial T]$,
giving total $\dim \B^2_5 = 4$).
\end{computation}

\begin{computation}[Two-generator free fermion dimension formula]
\label{comp:dim-formula-fermion}
The two-generator free fermion $\mathcal{F}_2$ has $r = 2$ fermionic
generators of weight $h_1 = h_2 = \tfrac{1}{2}$ (in the NS sector):
\[
\sum_{n,h} \dim \B^n_h \cdot q^h t^n
= \sum_{n\ge1}(n-1)!t^n
\left(\prod_{j=0}^{\infty}(1+q^{j+1/2})^2-1\right)^n .
\]
At $n = 1$: $\dim \B^1_{1/2} = 2$ ($\psi_1$, $\psi_2$),
$\dim \B^1_1 = 1$ ($\psi_1 \psi_2$-descendant), matching
\eqref{eq:fermion-char}.
The dimensions match Computation~\ref{comp:fermion-dim-table}.
\end{computation}

\begin{remark}[PBW approximation]\label{rem:PBW-approximation}
\index{PBW filtration!limitations}
Formula~\eqref{eq:universal-dim} counts the
\emph{PBW-level} dimensions of the bar complex: it gives
the dimension of $\B^n_h(A)$ as a graded vector space, before
taking cohomology. The bar differential $d$ reduces these
dimensions; the cohomological answer requires computing
$\ker d / \operatorname{im} d$ explicitly at each bidegree.
The PBW formula is exact for the total dimension of the bar complex,
but the bar cohomology depends on the OPE structure constants
(cf.\ the curvature-dependent analysis of
Computations~\ref{comp:virasoro-bar-diff}--\ref{comp:virasoro-curvature}).
In the Koszul case, the conformal-weight-graded bar cohomology
identifies with the dual coalgebra $A^i$ and hence, after
graded linear/Verdier duality, with the Hilbert function of $A^!$
(Propositions~\ref{prop:E8-koszul-acyclic}
and~\ref{prop:virasoro-koszul-acyclic}). The PBW formula gives
chain-level dimensions before the differential is imposed; it does
not compute total genus-$0$ bar cohomology.
\end{remark}

\begin{computation}[Koszul-dual Hilbert generating functions across standard families; \ClaimStatusConditional]
\label{comp:bar-cohomology-gfs}
\index{Koszul dual!Hilbert generating function}
\index{Motzkin numbers!Virasoro Koszul dual}
\index{partition function!Heisenberg Koszul dual}
The PBW filtration (Proposition~\ref{prop:universal-dim-formula})
plus the family-specific Koszul dual computation maps the
bidegree dimensions $\dim \B^n_h(A)$ to a one-variable generating
function tracking the Hilbert function of the Verdier--Koszul dual
algebra $A^!$.
For each standard family, the table records the closed-form generating
function, the combinatorial sequence encoded, and the dominant
singularity that controls the exponential growth rate.
Convention: $P_A(x) = \sum_{n \geq 1} \dim(A^!)_n \cdot x^n$.
These functions do not describe $Z^{\mathrm{der}}_{\mathrm{ch}}(A)$.

\emph{Direct expansions.}
Heisenberg: the first generator gives degree~$1$, while the remaining
degree-$n$ terms for $n \geq 2$ are shifted by $p(n-2)$; hence
$P_{\cH}(x) = x + x^2 + x^3 + 2x^4 + 3x^5 + 5x^6 + \cdots$.
Free fermion: $P_{\cF}(x) = x + x^2 + 2x^3 + 3x^4 + 5x^5 + \cdots$
(shifted by one degree).
Virasoro: $M(2) - M(1) = 1$, $M(3) - M(2) = 2$, $M(4) - M(3) = 5$,
$M(5) - M(4) = 12$, $M(6) - M(5) = 30$ (see
Proposition~\ref{prop:virasoro-koszul-acyclic}).
Betagamma: direct square-root expansion gives
$\sqrt{(1+x)/(1-3x)} = 1 + 2x + 4x^2 + 10x^3 + 26x^4 + 80x^5 + \cdots$
(squaring gives $(1 + 2x + 4x^2 + 10x^3 + 26x^4)^2 =
1 + 4x + 12x^2 + 36x^3 + 108x^4 + \cdots$, matching the expansion
of $(1+x)/(1-3x)$ through order $x^4$).
Affine $\widehat{\mathfrak{sl}}_{2,k}$: the Koszul-dual Hilbert
function gives $\dim(A^!)_2 = 5$, \emph{not} $6$: the chiral bar complex
departs from the Chevalley-Eilenberg count on the negative
subalgebra, because the chiral OPE introduces an Orlik-Solomon form
factor absent from the Lie-algebraic resolution, and the standard
Riordan recursion fails.
This is the \emph{modified Riordan} generating function.
Affine $\widehat{\mathfrak{sl}}_{3,k}$: the conditional dimensions
$[1, 8, 36, 204, 1352, 9892, \ldots]$ for $A^!$ in
Section~\ref{sec:sl3-bar-table} are encoded by
Conjecture~\ref{conj:sl3-bar-gf}:
\[
P_{\widehat{\mathfrak{sl}}_3}(x)
=
\frac{4x(2-13x-2x^2)}
{(1-8x)(1-3x-x^2)}.
\]
The dominant pole is $x=1/8$, so the growth is $8^n$. The factor
$1-3x-x^2$ is the DS-invariant spectral subfactor shared with the
$\mathcal{W}_3$ row; it is not the dominant growth pole.
\end{computation}

\begin{table}[h]
\centering
\renewcommand{\arraystretch}{1.3}
\begin{tabular}{@{}lllll@{}}
\toprule
Family & Generating function & OEIS / closed form & Growth rate & Discriminant \\
\midrule
Heisenberg $\cH$ &
$\displaystyle x + \sum_{n \geq 2} p(n{-}2)\, x^n$ &
A000041, shifted after degree $1$ &
$e^{\pi \sqrt{2n/3}}$ (sub-exp.) &
$\prod_{m \geq 1}(1 - x^m)$ \\
Free fermion $\cF$ &
$\displaystyle \sum_{n \geq 1} p(n{-}1)\, x^n$ &
A000041 shifted &
sub-exponential &
$\prod_{m \geq 1}(1 - x^m)$ \\
Betagamma $\beta\gamma$ &
$\displaystyle \sqrt{\frac{1+x}{1-3x}}$ &
algebraic, degree $2$ &
$3^n$ &
$(1 - 3x)(1 + x)$ \\
Affine $\widehat{\mathfrak{sl}}_{2,k}$ &
modified Riordan &
$H^2 = 5$ (not $6$) &
$3^n$ &
$(1 - 3x)(1 + x)$ \\
Virasoro $\mathrm{Vir}_c$ &
$\displaystyle \sum_{n \geq 1}\bigl(M(n{+}1) - M(n)\bigr)\, x^n$ &
A001006 (Motzkin) diff. &
$3^n$ &
$(1 - 3x)(1 + x)$ \\
Affine $\widehat{\mathfrak{sl}}_{3,k}$ &
rational, conditional &
$[1, 8, 36, 204, 1352, \ldots]$ &
$8^n$ &
$(1 - 8x)(1 - 3x - x^2)$ \\
\bottomrule
\end{tabular}
\caption[Koszul-dual Hilbert generating functions across standard families]{Generating functions for Koszul-dual Hilbert functions across standard families. The interacting families $\beta\gamma$, $\widehat{\mathfrak{sl}}_{2,k}$, and $\mathrm{Vir}_c$ share the universal discriminant $(1 - 3x)(1 + x)$ and growth rate $3^n$. The $\widehat{\mathfrak{sl}}_{3,k}$ row has dominant pole $x=1/8$, hence growth $8^n$, and a secondary DS-invariant factor $1 - 3x - x^2$ shared with $\mathcal{W}_3$. The free families $\cH, \cF$ have sub-exponential growth set by the Hardy-Ramanujan asymptotics of the partition function. The $\widehat{\mathfrak{sl}}_{2,k}$ entry is labelled \emph{modified} Riordan because the Koszul-dual Hilbert function gives $\dim(A^!)_2 = 5$ whereas the standard Riordan recursion (or the Chevalley-Eilenberg count on $\mathfrak{g}_-$) gives $6$; the departure is the Orlik-Solomon form factor of the chiral OPE, detailed in Section~\ref{sec:sl3-bar-table}.}
\label{tab:bar-cohomology-gfs}
\end{table}

\begin{remark}[The $3^n$ Koszul-dual growth question]
\label{rem:three-to-the-n-conjecture}
Three of the four interacting families in
Table~\ref{tab:bar-cohomology-gfs} ($\beta\gamma$,
$\widehat{\mathfrak{sl}}_{2,k}$, and $\mathrm{Vir}_c$) share both
the discriminant $(1 - 3x)(1 + x)$ and the exponential growth rate
$3^n$. This is the empirical content of the \emph{$3^n$
Koszul-dual Hilbert growth} question: for interacting chiral algebras in
the single-generator or rank-$1$ regime, the chiral bar complex
realizes its generic growth along the Motzkin path. The
$\widehat{\mathfrak{sl}}_{3,k}$ row exhibits the first departure:
rank-$2$ mixing introduces the dominant pole $x=1/8$ together with
the secondary factor $1-3x-x^2$. The latter is the DS-invariant
spectral factor; the former records the full
$\dim\mathfrak{sl}_3=8$ chain-growth mode. Whether a universal geometric mechanism on
$\overline{M}_{0,n+2}$ (e.g., a Motzkin-path envelope controlled by
the single-generator sector) governs the rank-$1$ case is an open
question. The generating functions here record full
graded dimensions of $A^!$, not Euler characteristics and not derived
centres; individual bar cohomology dimensions depend on the full OPE bracket and are not
reconstructible from $\dim \mathfrak{g}$ alone.
\end{remark}


\section{Finite horizontal Bar--BGG comparison for
\texorpdfstring{$\mathfrak{sl}_2$}{sl-2}}
\label{sec:bgg-sl2-computation}
\index{BGG resolution!from bar complex!sl2@$\mathfrak{sl}_2$|textbf}
\index{bar complex!BGG resolution!sl2@$\mathfrak{sl}_2$}

The finite horizontal matrix-level comparison behind
Theorem~\ref{thm:bgg-from-bar} for
$\widehat{\mathfrak{sl}}_{2,k}$, complementing the
module-bar comparison of
Computation~\ref{comp:bgg-sl2-computation}, uses the bar
differential matrix of Computation~\ref{comp:sl2-bar-matrix}
to isolate the horizontal classical BGG spine, with explicit singular-vector
conditions extracted from bar matrix entries and the
transformation of the resolution under module Koszul duality.

\subsection{Weyl group data and the classical resolution}

\begin{computation}[BGG data for \texorpdfstring{$\mathfrak{sl}_2$}{sl_2}, \texorpdfstring{$\lambda = 0$}{lambda = 0}]
\label{comp:bgg-sl2-data}
\index{BGG resolution!sl2@$\mathfrak{sl}_2$!Weyl orbit}

The Weyl group $W = \{1, s\}$ of $\mathfrak{sl}_2$ acts on
the weight lattice by the dot action
$w \cdot \lambda = w(\lambda + \rho) - \rho$,
with $\rho = 1$ (half the sum of positive roots).
For the trivial highest weight $\lambda = 0$:
\begin{align*}
1 \cdot 0 &= 0, &
s \cdot 0 &= s(0 + 1) - 1 = -1 - 1 = -2.
\end{align*}
The classical BGG resolution
(\cite{BGG76}; \cite{Kac}, \S10.7) is:
\begin{equation}\label{eq:bgg-sl2-classical}
0 \longrightarrow M(-2) \xrightarrow{\;\;d_1\;\;}
M(0) \xrightarrow{\;\;\varepsilon\;\;}
L(0) \longrightarrow 0
\end{equation}
where $M(\mu) = U(\mathfrak{sl}_2) \otimes_{U(\mathfrak{b})}
\mathbb{C}_\mu$ is the Verma module of highest weight~$\mu$,
$\varepsilon$ is the augmentation, and
$d_1(v_{-2}) = f \cdot v_0$ is the embedding via the singular
vector in $M(0)$ at weight~$-2$.
\end{computation}


\subsection{Bar spectral sequence contains the horizontal BGG spine}

\begin{proposition}[Horizontal Bar--BGG spine for
\texorpdfstring{$\widehat{\mathfrak{sl}}_{2,k}$}{sl-hat_2,k}; \ClaimStatusConditional]
\label{prop:bar-bgg-sl2}
\index{bar spectral sequence!BGG!sl2@$\mathfrak{sl}_2$}

For $\widehat{\mathfrak{sl}}_{2,k}$, the horizontal finite
$\mathfrak{sl}_2$ summand of the module bar spectral sequence of
Theorem~\textup{\ref{thm:bgg-from-bar}}, applied to the vacuum module
$L(0)$ and projected to the finite Weyl orbit of the zero weight, has
$E_1$ page
\begin{equation}\label{eq:bar-bgg-sl2-E1}
E_1^{n,0} = \bigoplus_{\ell(w) = n} \mathcal{M}(w \cdot 0),
\qquad E_1^{n,q} = 0 \;\;\text{for } q \neq 0,
\end{equation}
where $w$ ranges over the finite Weyl group
$W = \{1, s\}$. Explicitly:
$E_1^{0,0} = \mathcal{M}(0)$,
$E_1^{1,0} = \mathcal{M}(-2)$, and
$E_1^{n,0} = 0$ for $n \geq 2$.
The differential
$d_1\colon \mathcal{M}(-2) \to \mathcal{M}(0)$
is the screening operator integral
\begin{equation}\label{eq:bar-bgg-sl2-screening}
d_1(v) = \oint_{\gamma} S_\alpha(z)\, v\, \omega_{\log}
\end{equation}
where $S_\alpha$ is the screening current for the positive
root~$\alpha$ \textup{(}Theorem~\textup{\ref{thm:screening-bar})}
and $\gamma$ is the integration cycle in
$\overline{C}_2(X)$. The resulting resolution is the
finite horizontal BGG complex~\eqref{eq:bgg-sl2-classical}. It is
not the full affine BGG resolution, whose indexing set is the affine
Weyl group.
\end{proposition}

\begin{proof}
Apply Theorem~\ref{thm:bgg-from-bar} with
$\mathfrak{g} = \mathfrak{sl}_2$ and $\lambda = 0$.

\emph{Horizontal $E_1$ page.}
The PBW filtration on the chiral envelope
$\widehat{\mathfrak{sl}}_{2,k}$ induces a filtration on
the module bar complex
$\bar{B}_\bullet^{\mathrm{ch}}(\widehat{\mathfrak{sl}}_{2,k},
L(0))$. On the horizontal finite $\mathfrak{sl}_2$ summand, the
module-bar comparison gives
$E_1^{n,*} = \bigoplus_{\ell(w) = n} \mathcal{M}(w \cdot 0)$.
Since $|W| = 2$ with $\ell(1) = 0$ and $\ell(s) = 1$,
the nonzero terms are $E_1^{0,0} = \mathcal{M}(0)$ and
$E_1^{1,0} = \mathcal{M}(-2)$. The displayed page is the
projected horizontal subquotient; it is not a statement about all
affine-Weyl or all conformal-weight summands.

\emph{Differential.}
On this horizontal summand, the
differential $d_1\colon E_1^{1,0} \to E_1^{0,0}$
is the screening operator integral~\eqref{eq:bar-bgg-sl2-screening}.
For $\mathfrak{sl}_2$ with a single positive root~$\alpha$,
the universal property of Verma modules forces the unique
(up to scalar) module map
$\mathcal{M}(-2) \to \mathcal{M}(0)$ to be:
\[
d_1(v_{-2}) = f \cdot v_0
\]
(the singular vector of weight $-\alpha = -2$
in~$\mathcal{M}(0)$).

\emph{Exactness.}
The composition $d_1^2 = 0$ holds trivially
(there is no $E_1^{2,0}$ term). At the level of
configuration space integrals, this trivial vanishing
is the degenerate case of the Arnold relation
(Theorem~\ref{thm:arnold-relations}): with a single
positive root, no relation among screening operators
is needed. Exactness is the classical finite
Bernstein--Gel'fand--Gel'fand resolution
\cite{BGG76}; \cite[\S10.7]{Kac}: the image of
$M(-2)\to M(0)$ is the submodule generated by the singular vector
$f v_0$, and the cokernel is $L(0)$.
\end{proof}


\subsection{Connection to the bar differential matrix}

\begin{computation}[BGG data encoded in bar differential]
\label{comp:bgg-from-bar-matrix}
\index{bar complex!BGG!structure constants}

The $3 \times 9$ bar differential matrix
of~\eqref{eq:sl2-bar-matrix}
(Computation~\ref{comp:sl2-bar-matrix}) encodes the
BGG data as follows.

\emph{Structure constants and singular vectors.}
Each nonzero entry
$d([J^a \mid J^b]) = {f^{ab}}_c\, J^c$ is a structure
constant of $\mathfrak{sl}_2$. These same structure
constants build the singular vector
$f \cdot v_0 \in \mathcal{M}(0)_{-2}$ that generates
the BGG embedding
$\mathcal{M}(-2) \hookrightarrow \mathcal{M}(0)$.
Concretely, the singular-vector conditions
\begin{align}
e \cdot (f \cdot v_0)
 &= [e, f] \cdot v_0 = h \cdot v_0 = 0,
 \label{eq:sv-check-e} \\
h \cdot (f \cdot v_0)
 &= [h, f] \cdot v_0 + f \cdot (h \cdot v_0)
 = -2f \cdot v_0
 \label{eq:sv-check-h}
\end{align}
use the entries $d([e \mid f]) = h$ and
$d([h \mid f]) = -2f$ of~\eqref{eq:sl2-bar-matrix}.
The singular vector is annihilated by $e$ since
$d([e \mid f]) = h$ and $h \cdot v_0 = 0$
(the vacuum has weight~$0$); it has eigenvalue $-2$
under~$h$ since $d([h \mid f]) = -2f$ gives the
coefficient.

\emph{Curvature and genus-$1$ correction.}\quad
The curvature terms $k\kappa^{ab} \cdot \mathbf{1}$
in the bar differential
(Computation~\ref{comp:sl2-bar-matrix}) encode the
genus-$1$ obstruction: at genus~$0$ the bar complex
is strict ($d^2 = 0$); at genus~$1$,
$(d^{(1)})^2 = \kappa \cdot \omega_1$ where
$\kappa = 3(k{+}2)/4$
(Computation~\ref{comp:sl2-genus-table}).
On the finite horizontal complex, the level~$k$ enters through the
Shapovalov/screening normalization
(Proposition~\ref{prop:shapovalov-koszul}), not through the
two-term Weyl combinatorics. In the full affine complex the
dot-action weights also depend on~$k$; the bar chain groups
themselves remain level-independent
(Lemma~\ref{lem:bar-dims-level-independent}).

\emph{Rank and the Koszul dual.}
The rank-$3$ image of $d\colon \bar{B}^2 \to \bar{B}^1$
means every generator of
$\bar{B}^1 = \mathfrak{sl}_2$ is hit.
The $6$-dimensional kernel
(Computation~\ref{comp:sl2-bar-matrix}) encodes the
degree-$2$ relations in the dual coalgebra
$A^i=H^\star B(A)$. After graded linear/Verdier duality, these become
the quadratic relations in $A^!=(A^i)^\vee$
(Theorem~\ref{thm:sl2-koszul-dual}); they are not derived-centre
classes. The kernel identities are mediated by the Arnold relations
(Theorem~\ref{thm:arnold-relations}).
The shifted affine level $-k-4$ is the Feigin--Frenkel parameter of
the same scalar dual lane; it is not the raw bar cohomology object.
\end{computation}


\subsection{Module Koszul duality on the BGG resolution}

\begin{corollary}[BGG involution under Koszul duality;
\ClaimStatusConditional]
\label{cor:bgg-koszul-involution}
\index{BGG resolution!Koszul involution}
\index{module Koszul duality!BGG resolution}

On the finite-type graded-dual lane of
Theorem~\textup{\ref{thm:e1-module-koszul-duality}}, the module
bar functor~$\Phi$ sends the finite horizontal BGG complex at
level~$k$ to the corresponding bar-comodule complex; after the
graded-dual rewrite of Proposition~\ref{prop:vacuum-verma-koszul},
this is the finite horizontal BGG complex at the
Feigin--Frenkel dual level $k' = -k - 2h^\vee = -k - 4$:
\begin{equation}\label{eq:bgg-koszul-involution}
\Phi\bigl(0 \to \mathcal{M}(-2)_k \to
\mathcal{M}(0)_k \to L(0)_k \to 0\bigr)
\;\simeq\;
0 \to \mathcal{M}(-2)_{k'} \to
\mathcal{M}(0)_{k'} \to L(0)_{k'} \to 0.
\end{equation}
At the critical level $k = -h^\vee = -2$, the dual level is again
$-2$ and the scalar curvature
$\kappa(\widehat{\mathfrak{sl}}_{2,-2})=3(-2+2)/4$ vanishes. The
trace-form double pole in the current OPE remains; what vanishes is
the shifted Feigin--Frenkel curvature class.
\end{corollary}

\begin{proof}
By Proposition~\ref{prop:vacuum-verma-koszul},
$\Phi(\mathcal{M}(0)_k) \simeq \mathcal{M}(0)_{-k-4}$.
The same argument applied to
$\mathcal{M}(s \cdot 0)_k = \mathcal{M}(-2)_k$ gives
$\Phi(\mathcal{M}(-2)_k) \simeq \mathcal{M}(-2)_{-k-4}$:
the Feigin--Frenkel involution
acts on the level, not on the finite weight.
 The functoriality of~$\Phi$
 \textup{(}Definition~\textup{\ref{def:module-koszul-functor}}\textup{)}
 preserves the BGG differential:
$\Phi(d_1) = d_1'$, where $d_1'$ is the embedding at
level~$-k-4$.

At $k = -2$: the dual level is $-(-2) - 4 = -2$
(self-dual), and $\kappa(\widehat{\mathfrak{sl}}_{2,-2}) = 3(-2+2)/4 = 0$,
which is the critical-level universality statement
(Corollary~\ref{cor:critical-level-universality}).
The shifted bar curvature is zero at $k=-2$; the finite horizontal
constant-current summand is governed by the Chevalley--Eilenberg
differential $d_{\mathrm{CE}}$ of $\mathfrak{sl}_2$, with
$H^*(\mathfrak{sl}_2,\mathbb{C})=\Lambda(x_3)$
($x_3$ in degree~$3$).
\end{proof}


\begin{remark}[Finite Bar--BGG spine and higher rank]
\label{rem:bgg-computation}
\index{BGG resolution!finite Bar--BGG spine}

The finite horizontal comparison has the form:
\[
\underbrace{\text{OPE data}}_{\S\ref{chap:kac-moody}}
\;\xrightarrow{\;\text{residues}\;}
\underbrace{\text{bar differential}}_{\text{Comp.~\ref{comp:sl2-bar-matrix}}}
\;\xrightarrow{\;\text{screening}\;}
\underbrace{\text{BGG resolution}}_{\text{Prop.~\ref{prop:bar-bgg-sl2}}}
\;\xrightarrow{\;\Phi\;}
\underbrace{\text{dual BGG}}_{\text{Cor.~\ref{cor:bgg-koszul-involution}}}
\]
The OPE structure constants determine the residue differential on the
finite-current summand, while the classical BGG theorem supplies the
representation-theoretic exactness.

For higher-rank algebras, the same finite horizontal construction has
the following shape.
For $\mathfrak{g}$ with Weyl group~$W$, the $E_1$ page
of the bar spectral sequence has
$|W|$~Verma modules
$\{\mathcal{M}(w \cdot \lambda)\}_{w \in W}$
distributed by Weyl group length~$\ell(w)$.
Each simple root~$\alpha_i$ contributes a screening
operator~$S_{\alpha_i}$ that gives the BGG differentials on the
finite horizontal summand. The Arnold relations
(Theorem~\ref{thm:arnold-relations}) ensure
$d_1^2 = 0$: each relation
$\eta_{ij} \wedge \eta_{jk} + \eta_{jk} \wedge \eta_{ki}
+ \eta_{ki} \wedge \eta_{ij} = 0$
in $H^2(\overline{C}_3(X), \Omega^2_{\log})$ yields a
commutation relation among screening operators,
corresponding to the braid relation that
governs composition of BGG maps.

For $\mathfrak{sl}_3$ ($|W| = 6$, two positive roots):
the elements of $W = S_3$ distribute by length as
$(1, 2, 2, 1)$, giving a BGG resolution of length~$3$
with two screening operators $S_{\alpha_1}$,
$S_{\alpha_2}$. The Arnold relation for three
configuration points on $\overline{C}_3(\mathbb{P}^1)$
ensures the compatibility of the two screening charges.
The explicit $\mathfrak{sl}_3$ bar differential
(Computation~\ref{comp:sl3-bar}) supplies the finite-current input.
The full affine lift requires the same screening-normalization and
convergence obligations named in Theorem~\ref{thm:bgg-sl2-bar-explicit}.
\end{remark}

\begin{example}[\texorpdfstring{$\mathfrak{sl}_3$}{sl_3} finite BGG spine; \ClaimStatusProvedElsewhere]
\label{ex:sl3-bgg-computation}
\index{BGG resolution!sl3@$\mathfrak{sl}_3$}
\index{sl3@$\mathfrak{sl}_3$!finite BGG spine}
The finite horizontal spine for $\widehat{\mathfrak{sl}}_3$ uses the
$8$~generators
$\{H_1, H_2, E_1, E_2, E_3, F_1, F_2, F_3\}$
(Computation~\ref{comp:sl3-bar}).

\emph{Weyl group.}
$W = S_3$, with length distribution
$|\{w : \ell(w) = j\}| = 1, 2, 2, 1$ for $j = 0, 1, 2, 3$.
The finite BGG resolution of $L(0)$ has shape:
\begin{equation}\label{eq:sl3-bgg-shape}
0 \to \mathcal{M}(w_0 \cdot 0)
\xrightarrow{d_3}
\mathcal{M}(s_1 s_2 \cdot 0) \oplus \mathcal{M}(s_2 s_1 \cdot 0)
\xrightarrow{d_2}
\mathcal{M}(s_1 \cdot 0) \oplus \mathcal{M}(s_2 \cdot 0)
\xrightarrow{d_1}
\mathcal{M}(0) \to L(0) \to 0
\end{equation}
where $s_1, s_2$ are simple reflections and $w_0 = s_1 s_2 s_1$
is the longest element.

\emph{Screening operators.}
The two screening operators
$S_{\alpha_1}\colon \mathcal{M}(s_1 \cdot 0) \to \mathcal{M}(0)$
and
$S_{\alpha_2}\colon \mathcal{M}(s_2 \cdot 0) \to \mathcal{M}(0)$
are normalized by the classical BGG maps. The bar differential
supplies the current OPE coefficients:
$d[H_1|E_1]=2E_1$ and
$d[E_1|F_1]_{\bar B^1}=H_1$, with the vacuum term
$k|0\rangle$ tracked in $\bar B^0$
(Computation~\ref{comp:sl3-bar}). Proving that the corresponding
configuration-space integrals give the normalized affine screening
maps is the lift obligation.

\emph{Arnold compatibility.}
The Arnold relation for $\overline{C}_3(\mathbb{P}^1)$
ensures $d_1 \circ d_2 = 0$: the two compositions
$S_{\alpha_1} \circ S_{\alpha_2 \alpha_1}$ and
$S_{\alpha_2} \circ S_{\alpha_1 \alpha_2}$
agree up to sign, by the braid relation
$s_1 s_2 s_1 = s_2 s_1 s_2$ and the corresponding
Arnold identity on $3$-point logarithmic forms.

\emph{Koszul duality.}
Under $\Phi$ at level $k' = -k - 6$ (the Feigin--Frenkel dual
for $\mathfrak{sl}_3$, with $h^\vee = 3$), the finite-type
graded-dual lane sends the horizontal BGG complex
\eqref{eq:sl3-bgg-shape} to the corresponding dual-level
horizontal complex with Verma modules $\mathcal{M}(w \cdot 0)_{k'}$.
The Sugawara central charge transforms as
$c = 8k/(k+3) \mapsto c' = 8k'/(k'+3) = 16 - c$,
giving $c + c' = 16 = 2 \dim \mathfrak{sl}_3$.
\end{example}

\medskip

\noindent
These computations feed the Master Table of Computed Invariants
(Table~\ref{tab:master-invariants}).

\section{Synthesis}
\label{sec:bct-supplement}

\begin{remark}[Bar-complex table for $\mathbf H_{\Delta_5}$]
\label{rem:bct-1}
The bar-complex tables here extend to the K3 bialgebra
$\mathbf H_{\Delta_5}$ via the generating function
\[
\chi_{B_{\mathrm{ch}}(\mathbf H_{\Delta_5})}(q,\zeta,p) \;=\; (\Phi_{10}^{\mathrm{un}}(Z))^{-1} \cdot \prod_{(n,l,m)>0}(1 - q^n\zeta^l p^m)^{c(nm,l)},
\]
which (by the Gritsenko--Nikulin 1997--98 denominator identity) reduces
to $(\Phi_{10}^{\mathrm{un}}(Z))^{-1}\cdot\Phi_{10}^{\mathrm{un}}(Z) = 1$, i.e.\ the bar complex is
acyclic at the vacuum sector. The non-trivial entries of the bar-complex
table come from tensoring with each irreducible $\mathbf H_{\Delta_5}$-module:
for the character labelled by a Mukai vector $v\in\Lambda_{\mathrm{Muk}}$
of discriminant $\mathrm{disc}(v) = -2n$, the bar-complex multiplicity
at bidegree $(k,l)$ is the Fourier coefficient of
$(\Phi_{10}^{\mathrm{un}})^{-1}\cdot\theta_v$
where $\theta_v$ is the theta series attached to $v$. Cross-ref Vol.~III,
Chapter~\ref{V3-ch:k3-chiral-bialgebra-platonic},
\S\ref{sec:bar-complex-multiplicities}.
\end{remark}
